% !TeX program = pdflatex
% papers/cristal_xadrez.tex — GERADO por tools/cristal_tex.py; NÃO editar à mão.
% Fonte: cristal/cristal.jsonl (broca-so, última versão por conceito).
% Cada secção carrega o registo original em comentário %CRISTAL — a volta é
% tests/cristal_volta.js (resíduo 0 byte a byte contra a fonte).
\documentclass[11pt,a4paper]{article}
\usepackage[utf8]{inputenc}\usepackage[T1]{fontenc}\usepackage[brazil]{babel}
\usepackage{amsmath,amssymb}\usepackage{textcomp}
\usepackage[margin=2.4cm]{geometry}\usepackage{xcolor}
\definecolor{cinza}{gray}{0.35}
\newcommand{\prov}[1]{\par\smallskip\noindent{\small\color{cinza}#1}\par}
\setcounter{secnumdepth}{0}
% ─────────────────────────────────────────────────────────────────────────────
%  papers/gkcapa.tex — a capa do Reino, mínima, para os papers em `article`.
%
%  O estilo.tex completo é do `report` (títulos de capítulo, skak, …). Aqui fica
%  só o que a capa precisa, para o paper continuar a compilar sozinho com
%  pdflatex. O tradutor próprio (tests/tex) carrega o estilo.tex da raiz / o
%  symlink papers/estilo.tex e avalia o mesmo `\gkcapa`.
% ─────────────────────────────────────────────────────────────────────────────
\definecolor{ouro}{HTML}{B8912F}
\definecolor{ouroclaro}{HTML}{F3E9CE}
\definecolor{tinta}{HTML}{1E1B16}
\definecolor{regua}{HTML}{6E6858}
\providecommand{\gknota}{\fontsize{7.62}{11.05}\selectfont}
\providecommand{\gktexto}{\fontsize{10.50}{15.22}\selectfont}
\providecommand{\gksub}{\fontsize{12.33}{17.87}\selectfont}
\providecommand{\gksec}{\fontsize{14.47}{20.98}\selectfont}
\providecommand{\gkcap}{\fontsize{16.99}{24.63}\selectfont}
\providecommand{\gktit}{\fontsize{23.42}{33.95}\selectfont}
\providecommand{\Prj}{\mathbb{P}}
\providecommand{\Trf}{\mathcal{T}}
\providecommand{\gkcapa}[3]{%
\title{\vspace{-0.6cm}
{\color{ouro}\rule{\textwidth}{1.4pt}}\\[6mm]
\textbf{\gktit\color{tinta} Reino Dourado}\\[3mm]{\gksec\itshape\color{regua} Golden Kingdom}\\[8mm]
{\gkcap\scshape\color{ouro} #1}\\[5mm]
{\gksub\color{regua} #2}\\[2mm]
{\gktexto\color{regua}\itshape #3}\\[7mm]
{\color{ouro}\rule{\textwidth}{0.6pt}}\\[9mm]{\gknota\color{regua}\begin{minipage}{\textwidth}\setlength{\parindent}{0pt}%
\textbf{Aviso de Isenção Algorítmica:} O universo, as leis e as entidades contidas nestas páginas ---
incluindo felinos matriciais, aves de rapina algébricas e amuletos de controle de fluxo --- são
construções de pura ficção formal. Esta obra não descreve a realidade física, não serve como
engenharia reversa do cosmos e não constitui doutrina religiosa. Qualquer semelhança com bugs reais do seu
sistema operacional é mera coincidência (ou falta de reza). Prossiga por sua própria conta e
risco.\end{minipage}}}
\author{}\date{}}

\begin{document}
\gkcapa{O Cristal --- o xadrez do cristal}
       {partidas, mestres, estratégia}
       {corpus recuperado do broca-so; 150 conceitos; projeção verificável da fonte}
\maketitle

%CRISTAL {"arestas":[{"alvo":"xadrez","confianca":1.0,"epistemico":"fato","origem":"wiki","rel":"parte_de"}],"confianca":1.0,"contraexemplos":[],"descricao":"A primeira fase: disputar o CENTRO, desenvolver as peças menores, pôr o rei a salvo (roque). Teoria decorada por milhares de lances — o terreno mais estudado do jogo.","epistemico":"fato","exemplos":[],"id":"abertura_xadrez","memoria":"semantica","meta":{"dominio":"xadrez","fonte":"xadrez"},"origem":"wiki","palavras_chave":[],"sinonimos":[],"tipo":"conceito","titulo":"Abertura"}
\section{Abertura}
A primeira fase: disputar o CENTRO, desenvolver as peças menores, pôr o rei a salvo (roque). Teoria decorada por milhares de lances — o terreno mais estudado do jogo.
\prov{relações: parte\_de xadrez}
\prov{proveniência: wiki \textperiodcentered{} xadrez \textperiodcentered{} fato \textperiodcentered{} confiança 1.0 \textperiodcentered{} domínio xadrez \textperiodcentered{} história 2 versões no jornal}

%CRISTAL {"arestas":[{"alvo":"wu_wei","confianca":1.0,"epistemico":"fato","origem":"wiki","rel":"usa"},{"alvo":"taoismo","confianca":1.0,"epistemico":"fato","origem":"wiki","rel":"parte_de"}],"confianca":1.0,"contraexemplos":[],"descricao":"Lao Tzu: o modelo do sábio — o mais suave vence o mais duro, beneficia tudo sem disputar e ocupa os lugares desprezados. 'Nada é mais suave que a água, no entanto para atacar o rígido nada a supera'. Vencer sem combater, situar-se embaixo.","epistemico":"inferencia","exemplos":[],"id":"agua_taoista","memoria":"semantica","meta":{"autor":"Lao Tzu","dominio":"mestres","fonte":""},"origem":"wiki","palavras_chave":[],"sinonimos":[],"tipo":"conceito","titulo":"A Água como Sabedoria"}
\section{A Água como Sabedoria}
Lao Tzu: o modelo do sábio — o mais suave vence o mais duro, beneficia tudo sem disputar e ocupa os lugares desprezados. 'Nada é mais suave que a água, no entanto para atacar o rígido nada a supera'. Vencer sem combater, situar-se embaixo.
\prov{relações: usa wu\_wei; parte\_de taoismo}
\prov{proveniência: wiki \textperiodcentered{} inferencia \textperiodcentered{} confiança 1.0 \textperiodcentered{} domínio mestres \textperiodcentered{} história 3 versões no jornal}

%CRISTAL {"fusao":[{"arestas":[{"alvo":"sacrificio_xadrez","confianca":1.0,"epistemico":"fato","origem":"wiki","rel":"relaciona"}],"confianca":1.0,"contraexemplos":[],"descricao":"Campeão de ataque combinativo profundíssimo, uniu tática romântica e rigor posicional. Único a morrer com o título; suas combinações são estudo obrigatório.","epistemico":"fato","exemplos":[],"id":"alexander_alekhine","memoria":"semantica","meta":{"autor":"Alekhine","dominio":"xadrez","fonte":"história do xadrez"},"origem":"wiki","palavras_chave":[],"sinonimos":[],"tipo":"conceito","titulo":"Alexander Alekhine"},{"arestas":[{"alvo":"xadrez","confianca":1.0,"epistemico":"fato","origem":"wiki","rel":"relaciona"}],"confianca":1.0,"contraexemplos":[],"descricao":"Campeão mundial de 1927 a 1946, o maior atacante combinativo de sua era — ataques de profundidade e miniaturas fulminantes.","epistemico":"fato","exemplos":[],"id":"alekhine","memoria":"semantica","meta":{"dominio":"xadrez","fonte":"jogadores"},"origem":"wiki","palavras_chave":[],"sinonimos":[],"tipo":"conceito","titulo":"Alexander Alekhine"}],"id":"alexander_alekhine","tipo":"conceito"}
\section{Alexander Alekhine}
Campeão de ataque combinativo profundíssimo, uniu tática romântica e rigor posicional. Único a morrer com o título; suas combinações são estudo obrigatório.
\prov{relações: relaciona sacrificio\_xadrez}
\prov{proveniência: wiki \textperiodcentered{} história do xadrez \textperiodcentered{} fato \textperiodcentered{} confiança 1.0 \textperiodcentered{} domínio xadrez \textperiodcentered{} história 2 versões no jornal}
\prov{a segunda face --- alekhine:}
Campeão mundial de 1927 a 1946, o maior atacante combinativo de sua era — ataques de profundidade e miniaturas fulminantes.
\prov{fusão de curadoria (julgada): absorve alekhine --- as duas faces acima, intactas no registo}

%CRISTAL {"arestas":[{"alvo":"motor_de_xadrez","confianca":1.0,"epistemico":"fato","origem":"wiki","rel":"eh_um"},{"alvo":"aprendizado_por_reforco","confianca":1.0,"epistemico":"fato","origem":"wiki","rel":"usa"},{"alvo":"deep_blue_1997","confianca":1.0,"epistemico":"fato","origem":"wiki","rel":"contradiz"}],"confianca":1.0,"contraexemplos":[],"descricao":"A IA da DeepMind aprende xadrez SÓ jogando contra si mesma (aprendizado por reforço + busca Monte Carlo), sem teoria humana — em horas superou os melhores motores, com um estilo 'intuitivo' e sacrifícios posicionais. O auto-jogo reinventou o jogo.","epistemico":"fato","exemplos":[],"id":"alphazero_2017","memoria":"semantica","meta":{"autor":"DeepMind","dominio":"xadrez","fonte":"história do xadrez"},"origem":"wiki","palavras_chave":[],"sinonimos":[],"tipo":"conceito","titulo":"AlphaZero (2017)"}
\section{AlphaZero (2017)}
A IA da DeepMind aprende xadrez SÓ jogando contra si mesma (aprendizado por reforço + busca Monte Carlo), sem teoria humana — em horas superou os melhores motores, com um estilo 'intuitivo' e sacrifícios posicionais. O auto-jogo reinventou o jogo.
\prov{relações: eh\_um motor\_de\_xadrez; usa aprendizado\_por\_reforco; contradiz deep\_blue\_1997}
\prov{proveniência: wiki \textperiodcentered{} história do xadrez \textperiodcentered{} fato \textperiodcentered{} confiança 1.0 \textperiodcentered{} domínio xadrez \textperiodcentered{} história 4 versões no jornal}

%CRISTAL {"arestas":[{"alvo":"budismo","confianca":1.0,"epistemico":"fato","origem":"wiki","rel":"parte_de"},{"alvo":"meditacao_cura_da_mente_gentil","confianca":1.0,"epistemico":"fato","origem":"wiki","rel":"relaciona"},{"alvo":"aparelho_psiquico","confianca":1.0,"epistemico":"fato","origem":"wiki","rel":"relaciona"},{"alvo":"consciencia_software_universo_gentil","confianca":1.0,"epistemico":"fato","origem":"wiki","rel":"relaciona"}],"confianca":1.0,"contraexemplos":[],"descricao":"Buda: não existe um eu/self substancial, permanente e unitário; o que chamamos 'eu' é um feixe de processos impermanentes (agregados) sem núcleo essencial. Dissolver a ilusão do eu é o cerne do caminho.","epistemico":"inferencia","exemplos":[],"id":"anatta","memoria":"semantica","meta":{"autor":"Buda","dominio":"mestres","fonte":""},"origem":"wiki","palavras_chave":[],"sinonimos":[],"tipo":"conceito","titulo":"Anatta — o Não-Eu"}
\section{Anatta — o Não-Eu}
Buda: não existe um eu/self substancial, permanente e unitário; o que chamamos 'eu' é um feixe de processos impermanentes (agregados) sem núcleo essencial. Dissolver a ilusão do eu é o cerne do caminho.
\prov{relações: parte\_de budismo; relaciona meditacao\_cura\_da\_mente\_gentil; relaciona aparelho\_psiquico; relaciona consciencia\_software\_universo\_gentil}
\prov{proveniência: wiki \textperiodcentered{} inferencia \textperiodcentered{} confiança 1.0 \textperiodcentered{} domínio mestres \textperiodcentered{} história 5 versões no jornal}

%CRISTAL {"arestas":[{"alvo":"xadrez","confianca":1.0,"epistemico":"fato","origem":"wiki","rel":"relaciona"}],"confianca":1.0,"contraexemplos":[],"descricao":"O maior mestre do xadrez romântico (1818–1879), autor da Imortal e da Sempreviva. O ataque e o sacrifício acima do material — a estética antes da técnica posicional.","epistemico":"fato","exemplos":[],"id":"anderssen","memoria":"semantica","meta":{"dominio":"xadrez","fonte":"jogadores"},"origem":"wiki","palavras_chave":[],"sinonimos":[],"tipo":"conceito","titulo":"Adolf Anderssen"}
\section{Adolf Anderssen}
O maior mestre do xadrez romântico (1818–1879), autor da Imortal e da Sempreviva. O ataque e o sacrifício acima do material — a estética antes da técnica posicional.
\prov{relações: relaciona xadrez}
\prov{proveniência: wiki \textperiodcentered{} jogadores \textperiodcentered{} fato \textperiodcentered{} confiança 1.0 \textperiodcentered{} domínio xadrez \textperiodcentered{} história 18 versões no jornal}

%CRISTAL {"arestas":[{"alvo":"budismo","confianca":1.0,"epistemico":"fato","origem":"wiki","rel":"parte_de"},{"alvo":"face_bela_morte","confianca":1.0,"epistemico":"fato","origem":"wiki","rel":"relaciona"}],"confianca":1.0,"contraexemplos":[],"descricao":"Buda: tudo que é condicionado está em fluxo constante; nada é fixo — radicalizada como momentaneidade (tudo surge e cessa a cada instante). Aceitar a impermanência é libertar-se do apego.","epistemico":"inferencia","exemplos":[],"id":"anicca","memoria":"semantica","meta":{"autor":"Buda","dominio":"mestres","fonte":""},"origem":"wiki","palavras_chave":[],"sinonimos":[],"tipo":"conceito","titulo":"Anicca — a Impermanência"}
\section{Anicca — a Impermanência}
Buda: tudo que é condicionado está em fluxo constante; nada é fixo — radicalizada como momentaneidade (tudo surge e cessa a cada instante). Aceitar a impermanência é libertar-se do apego.
\prov{relações: parte\_de budismo; relaciona face\_bela\_morte}
\prov{proveniência: wiki \textperiodcentered{} inferencia \textperiodcentered{} confiança 1.0 \textperiodcentered{} domínio mestres \textperiodcentered{} história 3 versões no jornal}

%CRISTAL {"arestas":[{"alvo":"vedanta","confianca":1.0,"epistemico":"fato","origem":"wiki","rel":"parte_de"},{"alvo":"nao_dualidade","confianca":1.0,"epistemico":"fato","origem":"wiki","rel":"usa"},{"alvo":"consciencia_pura","confianca":1.0,"epistemico":"fato","origem":"wiki","rel":"usa"},{"alvo":"unidade_taoista","confianca":1.0,"epistemico":"fato","origem":"wiki","rel":"relaciona"}],"confianca":1.0,"contraexemplos":[],"descricao":"Shankara (Advaita Vedanta): o núcleo essencial do self individual (atman), entendido como consciência pura, é numericamente idêntico ao Absoluto (brahman) — 'um só, sem um segundo'. A não-dualidade radical.","epistemico":"inferencia","exemplos":[],"id":"atman_brahman","memoria":"semantica","meta":{"autor":"Shankara","dominio":"mestres","fonte":""},"origem":"wiki","palavras_chave":[],"sinonimos":[],"tipo":"conceito","titulo":"Atman = Brahman (Advaita)"}
\section{Atman = Brahman (Advaita)}
Shankara (Advaita Vedanta): o núcleo essencial do self individual (atman), entendido como consciência pura, é numericamente idêntico ao Absoluto (brahman) — 'um só, sem um segundo'. A não-dualidade radical.
\prov{relações: parte\_de vedanta; usa nao\_dualidade; usa consciencia\_pura; relaciona unidade\_taoista}
\prov{proveniência: wiki \textperiodcentered{} inferencia \textperiodcentered{} confiança 1.0 \textperiodcentered{} domínio mestres \textperiodcentered{} história 5 versões no jornal}

%CRISTAL {"arestas":[{"alvo":"informacao_assimetrica","confianca":1.0,"epistemico":"fato","origem":"wiki","rel":"usa"}],"confianca":1.0,"contraexemplos":[],"descricao":"Agir para MODELAR a crença do oponente: o blefe (parecer forte estando fraco), a finta, o sinal custoso (o gasto que só o forte pode bancar, logo crível). A guerra pela percepção.","epistemico":"inferencia","exemplos":[],"id":"blefe_e_sinal","memoria":"semantica","meta":{"dominio":"estrategia","fonte":"estratégia"},"origem":"wiki","palavras_chave":[],"sinonimos":[],"tipo":"conceito","titulo":"Blefe e Sinalização Estratégica"}
\section{Blefe e Sinalização Estratégica}
Agir para MODELAR a crença do oponente: o blefe (parecer forte estando fraco), a finta, o sinal custoso (o gasto que só o forte pode bancar, logo crível). A guerra pela percepção.
\prov{relações: usa informacao\_assimetrica}
\prov{proveniência: wiki \textperiodcentered{} estratégia \textperiodcentered{} inferencia \textperiodcentered{} confiança 1.0 \textperiodcentered{} domínio estrategia \textperiodcentered{} história 2 versões no jornal}

%CRISTAL {"fusao":[{"arestas":[{"alvo":"ruy_lopez","confianca":1.0,"epistemico":"fato","origem":"wiki","rel":"usa"}],"confianca":1.0,"contraexemplos":[],"descricao":"O americano que quebrou a hegemonia soviética (Campeão 1972, vs Spassky, a 'partida da Guerra Fria'). Precisão universal e vontade feroz; elevou o rigor e a profissionalização do jogo.","epistemico":"fato","exemplos":[],"id":"bobby_fischer","memoria":"semantica","meta":{"autor":"Fischer","dominio":"xadrez","fonte":"história do xadrez"},"origem":"wiki","palavras_chave":[],"sinonimos":[],"tipo":"conceito","titulo":"Bobby Fischer"},{"arestas":[{"alvo":"xadrez","confianca":1.0,"epistemico":"fato","origem":"wiki","rel":"relaciona"}],"confianca":1.0,"contraexemplos":[],"descricao":"Campeão mundial de 1972, um dos maiores gênios do jogo — precisão cristalina e força de cálculo. Aos 13 anos jogou a 'Partida do Século'.","epistemico":"fato","exemplos":[],"id":"fischer","memoria":"semantica","meta":{"dominio":"xadrez","fonte":"jogadores"},"origem":"wiki","palavras_chave":[],"sinonimos":[],"tipo":"conceito","titulo":"Bobby Fischer"}],"id":"bobby_fischer","tipo":"conceito"}
\section{Bobby Fischer}
O americano que quebrou a hegemonia soviética (Campeão 1972, vs Spassky, a 'partida da Guerra Fria'). Precisão universal e vontade feroz; elevou o rigor e a profissionalização do jogo.
\prov{relações: usa ruy\_lopez}
\prov{proveniência: wiki \textperiodcentered{} história do xadrez \textperiodcentered{} fato \textperiodcentered{} confiança 1.0 \textperiodcentered{} domínio xadrez \textperiodcentered{} história 2 versões no jornal}
\prov{a segunda face --- fischer:}
Campeão mundial de 1972, um dos maiores gênios do jogo — precisão cristalina e força de cálculo. Aos 13 anos jogou a 'Partida do Século'.
\prov{fusão de curadoria (julgada): absorve fischer --- as duas faces acima, intactas no registo}

%CRISTAL {"arestas":[{"alvo":"jogo_do_covarde","confianca":1.0,"epistemico":"fato","origem":"wiki","rel":"usa"},{"alvo":"dissuasao","confianca":1.0,"epistemico":"fato","origem":"wiki","rel":"especializa"}],"confianca":1.0,"contraexemplos":[],"descricao":"Elevar deliberadamente o risco de catástrofe compartilhada para forçar o outro a ceder primeiro — a estratégia do jogo do covarde levada ao limite. Crível justamente porque perigosa.","epistemico":"inferencia","exemplos":[],"id":"brinkmanship","memoria":"semantica","meta":{"autor":"Schelling","dominio":"estrategia","fonte":"estratégia"},"origem":"wiki","palavras_chave":[],"sinonimos":[],"tipo":"conceito","titulo":"Brinkmanship (a beira do abismo)"}
\section{Brinkmanship (a beira do abismo)}
Elevar deliberadamente o risco de catástrofe compartilhada para forçar o outro a ceder primeiro — a estratégia do jogo do covarde levada ao limite. Crível justamente porque perigosa.
\prov{relações: usa jogo\_do\_covarde; especializa dissuasao}
\prov{proveniência: wiki \textperiodcentered{} estratégia \textperiodcentered{} inferencia \textperiodcentered{} confiança 1.0 \textperiodcentered{} domínio estrategia \textperiodcentered{} história 3 versões no jornal}

%CRISTAL {"arestas":[{"alvo":"sabedoria_oriental","confianca":1.0,"epistemico":"fato","origem":"wiki","rel":"parte_de"},{"alvo":"word","confianca":0.5,"epistemico":"fato","origem":"wiki","rel":"relaciona"}],"confianca":1.0,"contraexemplos":[],"descricao":"A tradição do Buda — o diagnóstico do sofrimento e o caminho de sua cessação, pela compreensão do não-eu, da impermanência e da vacuidade; inclui o Zen.","epistemico":"fato","exemplos":[],"id":"budismo","memoria":"semantica","meta":{"autor":"Buda","dominio":"mestres","fonte":""},"origem":"wiki","palavras_chave":[],"sinonimos":[],"tipo":"conceito","titulo":"Budismo"}
\section{Budismo}
A tradição do Buda — o diagnóstico do sofrimento e o caminho de sua cessação, pela compreensão do não-eu, da impermanência e da vacuidade; inclui o Zen.
\prov{relações: parte\_de sabedoria\_oriental; relaciona word}
\prov{proveniência: wiki \textperiodcentered{} fato \textperiodcentered{} confiança 1.0 \textperiodcentered{} domínio mestres \textperiodcentered{} história 3 versões no jornal}

%CRISTAL {"arestas":[{"alvo":"xadrez","confianca":1.0,"epistemico":"fato","origem":"wiki","rel":"relaciona"}],"confianca":1.0,"contraexemplos":[],"descricao":"Mestre americano (1930–1976), do lado que perdeu a 'Partida do Século' para o jovem Fischer — imortalizado por isso.","epistemico":"fato","exemplos":[],"id":"byrne","memoria":"semantica","meta":{"dominio":"xadrez","fonte":"jogadores"},"origem":"wiki","palavras_chave":[],"sinonimos":[],"tipo":"conceito","titulo":"Donald Byrne"}
\section{Donald Byrne}
Mestre americano (1930–1976), do lado que perdeu a 'Partida do Século' para o jovem Fischer — imortalizado por isso.
\prov{relações: relaciona xadrez}
\prov{proveniência: wiki \textperiodcentered{} jogadores \textperiodcentered{} fato \textperiodcentered{} confiança 1.0 \textperiodcentered{} domínio xadrez \textperiodcentered{} história 6 versões no jornal}

%CRISTAL {"arestas":[{"alvo":"xadrez","confianca":1.0,"epistemico":"fato","origem":"wiki","rel":"relaciona"}],"confianca":1.0,"contraexemplos":[],"descricao":"Grande mestre peruano (1896–1981), autor da 'Imortal Peruana', miniatura de sacrifícios que termina em mate de bispo.","epistemico":"fato","exemplos":[],"id":"canal","memoria":"semantica","meta":{"dominio":"xadrez","fonte":"jogadores"},"origem":"wiki","palavras_chave":[],"sinonimos":[],"tipo":"conceito","titulo":"Esteban Canal"}
\section{Esteban Canal}
Grande mestre peruano (1896–1981), autor da 'Imortal Peruana', miniatura de sacrifícios que termina em mate de bispo.
\prov{relações: relaciona xadrez}
\prov{proveniência: wiki \textperiodcentered{} jogadores \textperiodcentered{} fato \textperiodcentered{} confiança 1.0 \textperiodcentered{} domínio xadrez \textperiodcentered{} história 6 versões no jornal}

%CRISTAL {"arestas":[{"alvo":"sabedoria_oriental","confianca":1.0,"epistemico":"fato","origem":"wiki","rel":"relaciona"},{"alvo":"colapso_observador","confianca":1.0,"epistemico":"fato","origem":"wiki","rel":"usa"},{"alvo":"tao","confianca":1.0,"epistemico":"fato","origem":"wiki","rel":"relaciona"}],"confianca":1.0,"contraexemplos":[],"descricao":"Capra (1975): a tese de que a física moderna (quântica e de partículas) exibe paralelos com o misticismo oriental — o papel do observador, a interconexão universal e o 'vazio' dinâmico ecoando sunyata/Tao. Analogia CONTESTADA — não é física estabelecida.","epistemico":"hipotese","exemplos":[],"id":"capra_tao_fisica","memoria":"semantica","meta":{"autor":"Capra","dominio":"mestres","fonte":""},"origem":"wiki","palavras_chave":[],"sinonimos":[],"tipo":"conceito","titulo":"O Tao da Física (Capra)"}
\section{O Tao da Física (Capra)}
Capra (1975): a tese de que a física moderna (quântica e de partículas) exibe paralelos com o misticismo oriental — o papel do observador, a interconexão universal e o 'vazio' dinâmico ecoando sunyata/Tao. Analogia CONTESTADA — não é física estabelecida.
\prov{relações: relaciona sabedoria\_oriental; usa colapso\_observador; relaciona tao}
\prov{proveniência: wiki \textperiodcentered{} hipotese \textperiodcentered{} confiança 1.0 \textperiodcentered{} domínio mestres \textperiodcentered{} história 4 versões no jornal}

%CRISTAL {"arestas":[{"alvo":"estrategia_posicional","confianca":1.0,"epistemico":"fato","origem":"wiki","rel":"usa"}],"confianca":1.0,"contraexemplos":[],"descricao":"Uma casa que nenhum peão pode mais defender — posto de honra para um cavalo inimigo cravado ali. A estratégia posicional caça e ocupa esses buracos.","epistemico":"inferencia","exemplos":[],"id":"casa_fraca","memoria":"semantica","meta":{"dominio":"xadrez","fonte":"xadrez"},"origem":"wiki","palavras_chave":[],"sinonimos":[],"tipo":"conceito","titulo":"Casa Fraca (buraco)"}
\section{Casa Fraca (buraco)}
Uma casa que nenhum peão pode mais defender — posto de honra para um cavalo inimigo cravado ali. A estratégia posicional caça e ocupa esses buracos.
\prov{relações: usa estrategia\_posicional}
\prov{proveniência: wiki \textperiodcentered{} xadrez \textperiodcentered{} inferencia \textperiodcentered{} confiança 1.0 \textperiodcentered{} domínio xadrez \textperiodcentered{} história 2 versões no jornal}

%CRISTAL {"arestas":[{"alvo":"clausewitz_da_guerra","confianca":1.0,"epistemico":"fato","origem":"wiki","rel":"parte_de"}],"confianca":1.0,"contraexemplos":[],"descricao":"O ponto de força do qual tudo depende — o exército, a aliança, a vontade popular. Concentrar contra ele decide a campanha; dispersar esforço é derrota. O 'ponto único de falha' da estratégia.","epistemico":"inferencia","exemplos":[],"id":"centro_de_gravidade","memoria":"semantica","meta":{"dominio":"estrategia","fonte":"estratégia"},"origem":"wiki","palavras_chave":[],"sinonimos":[],"tipo":"conceito","titulo":"Centro de Gravidade"}
\section{Centro de Gravidade}
O ponto de força do qual tudo depende — o exército, a aliança, a vontade popular. Concentrar contra ele decide a campanha; dispersar esforço é derrota. O 'ponto único de falha' da estratégia.
\prov{relações: parte\_de clausewitz\_da\_guerra}
\prov{proveniência: wiki \textperiodcentered{} estratégia \textperiodcentered{} inferencia \textperiodcentered{} confiança 1.0 \textperiodcentered{} domínio estrategia \textperiodcentered{} história 2 versões no jornal}

%CRISTAL {"arestas":[{"alvo":"estrategia_empresarial","confianca":1.0,"epistemico":"fato","origem":"wiki","rel":"usa"}],"confianca":1.0,"contraexemplos":[],"descricao":"A atratividade de um setor vem de cinco pressões: rivalidade, entrantes, substitutos, poder de fornecedores e de compradores. O mapa para achar onde há lucro a defender.","epistemico":"inferencia","exemplos":[],"id":"cinco_forcas_porter","memoria":"semantica","meta":{"autor":"Michael Porter","dominio":"estrategia","fonte":"estratégia"},"origem":"wiki","palavras_chave":[],"sinonimos":[],"tipo":"conceito","titulo":"Cinco Forças (Porter)"}
\section{Cinco Forças (Porter)}
A atratividade de um setor vem de cinco pressões: rivalidade, entrantes, substitutos, poder de fornecedores e de compradores. O mapa para achar onde há lucro a defender.
\prov{relações: usa estrategia\_empresarial}
\prov{proveniência: wiki \textperiodcentered{} estratégia \textperiodcentered{} inferencia \textperiodcentered{} confiança 1.0 \textperiodcentered{} domínio estrategia \textperiodcentered{} história 2 versões no jornal}

%CRISTAL {"arestas":[{"alvo":"estrategia","confianca":1.0,"epistemico":"fato","origem":"wiki","rel":"explica"}],"confianca":1.0,"contraexemplos":[],"descricao":"A guerra é a continuação da política por outros meios. Introduz a FRICÇÃO (o atrito do real), a NÉVOA da guerra (incerteza) e o centro de gravidade (o ponto cuja queda derruba tudo). O clássico ocidental.","epistemico":"fato","exemplos":[],"id":"clausewitz_da_guerra","memoria":"semantica","meta":{"autor":"Clausewitz","dominio":"estrategia","fonte":"estratégia"},"origem":"wiki","palavras_chave":[],"sinonimos":[],"tipo":"conceito","titulo":"Da Guerra (Clausewitz)"}
\section{Da Guerra (Clausewitz)}
A guerra é a continuação da política por outros meios. Introduz a FRICÇÃO (o atrito do real), a NÉVOA da guerra (incerteza) e o centro de gravidade (o ponto cuja queda derruba tudo). O clássico ocidental.
\prov{relações: explica estrategia}
\prov{proveniência: wiki \textperiodcentered{} estratégia \textperiodcentered{} fato \textperiodcentered{} confiança 1.0 \textperiodcentered{} domínio estrategia \textperiodcentered{} história 2 versões no jornal}

%CRISTAL {"arestas":[{"alvo":"ponto_focal_schelling","confianca":1.0,"epistemico":"fato","origem":"wiki","rel":"relaciona"},{"alvo":"teoria_dos_jogos","confianca":1.0,"epistemico":"fato","origem":"wiki","rel":"usa"},{"alvo":"face_turno_fail_closed","confianca":1.0,"epistemico":"fato","origem":"wiki","rel":"semelhante"}],"confianca":1.0,"contraexemplos":[],"descricao":"Amarrar as próprias mãos VISIVELMENTE para mudar o jogo do outro — queimar a ponte atrás de si torna a retirada impossível e a ameaça real. O poder de se comprometer é o poder de vencer. Ecoa o fail-closed.","epistemico":"inferencia","exemplos":[],"id":"compromisso_credivel","memoria":"semantica","meta":{"autor":"Schelling","dominio":"estrategia","fonte":"estratégia"},"origem":"wiki","palavras_chave":[],"sinonimos":[],"tipo":"conceito","titulo":"Compromisso Crível (Schelling)"}
\section{Compromisso Crível (Schelling)}
Amarrar as próprias mãos VISIVELMENTE para mudar o jogo do outro — queimar a ponte atrás de si torna a retirada impossível e a ameaça real. O poder de se comprometer é o poder de vencer. Ecoa o fail-closed.
\prov{relações: relaciona ponto\_focal\_schelling; usa teoria\_dos\_jogos; semelhante face\_turno\_fail\_closed}
\prov{proveniência: wiki \textperiodcentered{} estratégia \textperiodcentered{} inferencia \textperiodcentered{} confiança 1.0 \textperiodcentered{} domínio estrategia \textperiodcentered{} história 4 versões no jornal}

%CRISTAL {"arestas":[{"alvo":"sabedoria_oriental","confianca":1.0,"epistemico":"fato","origem":"wiki","rel":"parte_de"},{"alvo":"word","confianca":0.5,"epistemico":"fato","origem":"wiki","rel":"relaciona"}],"confianca":1.0,"contraexemplos":[],"descricao":"A tradição dos Analectos — o cultivo moral de si (ren, li, yi) como base da ordem social e política; a via humana da virtude e do aprendizado.","epistemico":"fato","exemplos":[],"id":"confucionismo","memoria":"semantica","meta":{"autor":"Confúcio","dominio":"mestres","fonte":""},"origem":"wiki","palavras_chave":[],"sinonimos":[],"tipo":"conceito","titulo":"Confucionismo (Confúcio)"}
\section{Confucionismo (Confúcio)}
A tradição dos Analectos — o cultivo moral de si (ren, li, yi) como base da ordem social e política; a via humana da virtude e do aprendizado.
\prov{relações: parte\_de sabedoria\_oriental; relaciona word}
\prov{proveniência: wiki \textperiodcentered{} fato \textperiodcentered{} confiança 1.0 \textperiodcentered{} domínio mestres \textperiodcentered{} história 3 versões no jornal}

%CRISTAL {"arestas":[{"alvo":"inteligencia_consciencia_relacional","confianca":1.0,"epistemico":"fato","origem":"wiki","rel":"usa"},{"alvo":"meditacao_cura_da_mente_gentil","confianca":1.0,"epistemico":"fato","origem":"wiki","rel":"relaciona"},{"alvo":"atman_brahman","confianca":1.0,"epistemico":"fato","origem":"wiki","rel":"relaciona"},{"alvo":"tao_da_matematica","confianca":1.0,"epistemico":"fato","origem":"wiki","rel":"parte_de"}],"confianca":1.0,"contraexemplos":[],"descricao":"Lopes: 'algo que observa é algo consciente' — existe uma única Consciência (Deus) que, ao incidir no cérebro (prisma/hardware), se decompõe em pensamentos; a Consciência é 'o Software do Universo', programando até as partículas subatômicas. Meditar é eliminar pensamentos e retornar à Consciência original.","epistemico":"inferencia","exemplos":[],"id":"consciencia_software_universo_gentil","memoria":"semantica","meta":{"autor":"Gentil Lopes","dominio":"mestres","fonte":""},"origem":"wiki","palavras_chave":[],"sinonimos":[],"tipo":"conceito","titulo":"A Consciência como Software do Universo (Lopes)"}
\section{A Consciência como Software do Universo (Lopes)}
Lopes: 'algo que observa é algo consciente' — existe uma única Consciência (Deus) que, ao incidir no cérebro (prisma/hardware), se decompõe em pensamentos; a Consciência é 'o Software do Universo', programando até as partículas subatômicas. Meditar é eliminar pensamentos e retornar à Consciência original.
\prov{relações: usa inteligencia\_consciencia\_relacional; relaciona meditacao\_cura\_da\_mente\_gentil; relaciona atman\_brahman; parte\_de tao\_da\_matematica}
\prov{proveniência: wiki \textperiodcentered{} inferencia \textperiodcentered{} confiança 1.0 \textperiodcentered{} domínio mestres \textperiodcentered{} história 6 versões no jornal}

%CRISTAL {"arestas":[{"alvo":"tatica_xadrez","confianca":1.0,"epistemico":"fato","origem":"wiki","rel":"especializa"}],"confianca":1.0,"contraexemplos":[],"descricao":"Uma peça não pode (ou não deve) sair porque expõe algo mais valioso atrás — absoluta (contra o rei, ilegal mover) ou relativa. Imobiliza e prepara o ataque à peça cravada.","epistemico":"fato","exemplos":[],"id":"cravada","memoria":"semantica","meta":{"dominio":"xadrez","fonte":"xadrez"},"origem":"wiki","palavras_chave":[],"sinonimos":[],"tipo":"conceito","titulo":"Cravada (Pin)"}
\section{Cravada (Pin)}
Uma peça não pode (ou não deve) sair porque expõe algo mais valioso atrás — absoluta (contra o rei, ilegal mover) ou relativa. Imobiliza e prepara o ataque à peça cravada.
\prov{relações: especializa tatica\_xadrez}
\prov{proveniência: wiki \textperiodcentered{} xadrez \textperiodcentered{} fato \textperiodcentered{} confiança 1.0 \textperiodcentered{} domínio xadrez \textperiodcentered{} história 2 versões no jornal}

%CRISTAL {"arestas":[{"alvo":"tianming","confianca":1.0,"epistemico":"fato","origem":"wiki","rel":"usa"},{"alvo":"contrato_social","confianca":1.0,"epistemico":"fato","origem":"wiki","rel":"relaciona"}],"confianca":1.0,"contraexemplos":[],"descricao":"Confúcio: governar (zheng) é retificar (zheng) — o governante virtuoso (de) atrai o povo sem coerção: 'governar pela virtude é como a Estrela Polar, que permanece em seu lugar enquanto as demais lhe prestam homenagem'. A ordem vem do exemplo, não da punição.","epistemico":"inferencia","exemplos":[],"id":"de_zheng","memoria":"semantica","meta":{"autor":"Confúcio","dominio":"mestres","fonte":""},"origem":"wiki","palavras_chave":[],"sinonimos":[],"tipo":"conceito","titulo":"Governo pela Virtude (de/zheng)"}
\section{Governo pela Virtude (de/zheng)}
Confúcio: governar (zheng) é retificar (zheng) — o governante virtuoso (de) atrai o povo sem coerção: 'governar pela virtude é como a Estrela Polar, que permanece em seu lugar enquanto as demais lhe prestam homenagem'. A ordem vem do exemplo, não da punição.
\prov{relações: usa tianming; relaciona contrato\_social}
\prov{proveniência: wiki \textperiodcentered{} inferencia \textperiodcentered{} confiança 1.0 \textperiodcentered{} domínio mestres \textperiodcentered{} história 3 versões no jornal}

%CRISTAL {"arestas":[{"alvo":"xadrez","confianca":1.0,"epistemico":"fato","origem":"wiki","rel":"relaciona"},{"alvo":"motor_de_xadrez","confianca":1.0,"epistemico":"fato","origem":"wiki","rel":"eh_um"},{"alvo":"inteligencia_artificial","confianca":1.0,"epistemico":"fato","origem":"wiki","rel":"parte_de"}],"confianca":1.0,"contraexemplos":[],"descricao":"O supercomputador da IBM que em 1997 venceu Garry Kasparov num match — a primeira vez que uma máquina bateu o campeão mundial. Busca em árvore massiva (minimax + poda) em hardware dedicado.","epistemico":"fato","exemplos":[],"id":"deep_blue","memoria":"semantica","meta":{"dominio":"xadrez","fonte":"jogadores"},"origem":"wiki","palavras_chave":[],"sinonimos":[],"tipo":"conceito","titulo":"Deep Blue (IBM)"}
\section{Deep Blue (IBM)}
O supercomputador da IBM que em 1997 venceu Garry Kasparov num match — a primeira vez que uma máquina bateu o campeão mundial. Busca em árvore massiva (minimax + poda) em hardware dedicado.
\prov{relações: relaciona xadrez; eh\_um motor\_de\_xadrez; parte\_de inteligencia\_artificial}
\prov{proveniência: wiki \textperiodcentered{} jogadores \textperiodcentered{} fato \textperiodcentered{} confiança 1.0 \textperiodcentered{} domínio xadrez \textperiodcentered{} história 24 versões no jornal}

%CRISTAL {"arestas":[{"alvo":"garry_kasparov","confianca":1.0,"epistemico":"fato","origem":"wiki","rel":"relaciona"},{"alvo":"motor_de_xadrez","confianca":1.0,"epistemico":"fato","origem":"wiki","rel":"eh_um"},{"alvo":"minimax_alfabeta","confianca":1.0,"epistemico":"fato","origem":"wiki","rel":"usa"}],"confianca":1.0,"contraexemplos":[],"descricao":"A máquina da IBM derrota o campeão mundial num match — a primeira vez que um computador vence o melhor humano. Força bruta: ~200 milhões de posições/s por minimax alfa-beta e hardware dedicado.","epistemico":"fato","exemplos":[],"id":"deep_blue_1997","memoria":"semantica","meta":{"autor":"IBM","dominio":"xadrez","fonte":"história do xadrez"},"origem":"wiki","palavras_chave":[],"sinonimos":[],"tipo":"conceito","titulo":"Deep Blue vence Kasparov (1997)"}
\section{Deep Blue vence Kasparov (1997)}
A máquina da IBM derrota o campeão mundial num match — a primeira vez que um computador vence o melhor humano. Força bruta: \~{}200 milhões de posições/s por minimax alfa-beta e hardware dedicado.
\prov{relações: relaciona garry\_kasparov; eh\_um motor\_de\_xadrez; usa minimax\_alfabeta}
\prov{proveniência: wiki \textperiodcentered{} história do xadrez \textperiodcentered{} fato \textperiodcentered{} confiança 1.0 \textperiodcentered{} domínio xadrez \textperiodcentered{} história 4 versões no jornal}

%CRISTAL {"arestas":[{"alvo":"abertura_xadrez","confianca":1.0,"epistemico":"fato","origem":"wiki","rel":"especializa"}],"confianca":1.0,"contraexemplos":[],"descricao":"A resposta mais combativa a 1.e4: as pretas desequilibram desde o início, jogando por contra-ataque assimétrico. A abertura mais jogada e analisada do xadrez de elite.","epistemico":"fato","exemplos":[],"id":"defesa_siciliana","memoria":"semantica","meta":{"dominio":"xadrez","fonte":"xadrez"},"origem":"wiki","palavras_chave":[],"sinonimos":[],"tipo":"conceito","titulo":"Defesa Siciliana (1.e4 c5)"}
\section{Defesa Siciliana (1.e4 c5)}
A resposta mais combativa a 1.e4: as pretas desequilibram desde o início, jogando por contra-ataque assimétrico. A abertura mais jogada e analisada do xadrez de elite.
\prov{relações: especializa abertura\_xadrez}
\prov{proveniência: wiki \textperiodcentered{} xadrez \textperiodcentered{} fato \textperiodcentered{} confiança 1.0 \textperiodcentered{} domínio xadrez \textperiodcentered{} história 2 versões no jornal}

%CRISTAL {"arestas":[{"alvo":"borda_critica_vazio","confianca":1.0,"epistemico":"fato","origem":"wiki","rel":"usa"},{"alvo":"sunyata","confianca":1.0,"epistemico":"fato","origem":"wiki","rel":"explica"},{"alvo":"vazio_util","confianca":1.0,"epistemico":"fato","origem":"wiki","rel":"relaciona"},{"alvo":"criacao_livre_fertilidade","confianca":1.0,"epistemico":"fato","origem":"wiki","rel":"relaciona"},{"alvo":"tao_da_matematica","confianca":1.0,"epistemico":"fato","origem":"wiki","rel":"parte_de"}],"confianca":1.0,"contraexemplos":[],"descricao":"Lopes: como o ponto é indimensional, intervalos e cubos 'não ocupam espaço, são Vazios'; logo o universo 0,1^∞, identificado a Deus, é Vazio — 'o próprio Deus é Vazio'. Alinha isto ao suniata/zero de Buda e ao 'Vácuo, o Nada, berço de todos os Possíveis' de Lao Tsé. A realidade coincide com o Vazio (a 'tela em branco'), e tudo é pintura da mente.","epistemico":"inferencia","exemplos":[],"id":"deus_como_vazio_gentil","memoria":"semantica","meta":{"autor":"Gentil Lopes","dominio":"mestres","fonte":""},"origem":"wiki","palavras_chave":[],"sinonimos":[],"tipo":"conceito","titulo":"Deus como o Grande Vazio (Lopes)"}
\section{Deus como o Grande Vazio (Lopes)}
Lopes: como o ponto é indimensional, intervalos e cubos 'não ocupam espaço, são Vazios'; logo o universo 0,1\^{}\ensuremath{\infty}, identificado a Deus, é Vazio — 'o próprio Deus é Vazio'. Alinha isto ao suniata/zero de Buda e ao 'Vácuo, o Nada, berço de todos os Possíveis' de Lao Tsé. A realidade coincide com o Vazio (a 'tela em branco'), e tudo é pintura da mente.
\prov{relações: usa borda\_critica\_vazio; explica sunyata; relaciona vazio\_util; relaciona criacao\_livre\_fertilidade; parte\_de tao\_da\_matematica}
\prov{proveniência: wiki \textperiodcentered{} inferencia \textperiodcentered{} confiança 1.0 \textperiodcentered{} domínio mestres \textperiodcentered{} história 8 versões no jornal}

%CRISTAL {"arestas":[{"alvo":"estrategia","confianca":1.0,"epistemico":"fato","origem":"wiki","rel":"especializa"},{"alvo":"compromisso_credivel","confianca":1.0,"epistemico":"fato","origem":"wiki","rel":"usa"}],"confianca":1.0,"contraexemplos":[],"descricao":"Evitar o ataque tornando a resposta insuportável e CRÍVEL — 'não me ataque porque a retaliação te custa mais'. A lógica da guerra fria e do equilíbrio do terror (MAD).","epistemico":"inferencia","exemplos":[],"id":"dissuasao","memoria":"semantica","meta":{"dominio":"estrategia","fonte":"estratégia"},"origem":"wiki","palavras_chave":[],"sinonimos":[],"tipo":"conceito","titulo":"Dissuasão (Deterrence)"}
\section{Dissuasão (Deterrence)}
Evitar o ataque tornando a resposta insuportável e CRÍVEL — 'não me ataque porque a retaliação te custa mais'. A lógica da guerra fria e do equilíbrio do terror (MAD).
\prov{relações: especializa estrategia; usa compromisso\_credivel}
\prov{proveniência: wiki \textperiodcentered{} estratégia \textperiodcentered{} inferencia \textperiodcentered{} confiança 1.0 \textperiodcentered{} domínio estrategia \textperiodcentered{} história 3 versões no jornal}

%CRISTAL {"arestas":[{"alvo":"xadrez","confianca":1.0,"epistemico":"fato","origem":"wiki","rel":"relaciona"}],"confianca":1.0,"contraexemplos":[],"descricao":"Mestre alemão-americano (1885–1981), autor da célebre 'caminhada do rei' de 1912 — não confundir com o campeão Emanuel Lasker.","epistemico":"fato","exemplos":[],"id":"edward_lasker","memoria":"semantica","meta":{"dominio":"xadrez","fonte":"jogadores"},"origem":"wiki","palavras_chave":[],"sinonimos":[],"tipo":"conceito","titulo":"Edward Lasker"}
\section{Edward Lasker}
Mestre alemão-americano (1885–1981), autor da célebre 'caminhada do rei' de 1912 — não confundir com o campeão Emanuel Lasker.
\prov{relações: relaciona xadrez}
\prov{proveniência: wiki \textperiodcentered{} jogadores \textperiodcentered{} fato \textperiodcentered{} confiança 1.0 \textperiodcentered{} domínio xadrez \textperiodcentered{} história 4 versões no jornal}

%CRISTAL {"arestas":[{"alvo":"xadrez","confianca":1.0,"epistemico":"fato","origem":"wiki","rel":"relaciona"}],"confianca":1.0,"contraexemplos":[],"descricao":"A escala que mede a força relativa pela probabilidade de vitória — diferença de 400 pontos ≈ 10:1. Criada por Arpad Elo para o xadrez, hoje usada de jogos a ranking de modelos.","epistemico":"fato","exemplos":[],"id":"elo_rating","memoria":"semantica","meta":{"autor":"Arpad Elo","dominio":"xadrez","fonte":"xadrez"},"origem":"wiki","palavras_chave":[],"sinonimos":[],"tipo":"conceito","titulo":"Rating Elo"}
\section{Rating Elo}
A escala que mede a força relativa pela probabilidade de vitória — diferença de 400 pontos \ensuremath{\approx} 10:1. Criada por Arpad Elo para o xadrez, hoje usada de jogos a ranking de modelos.
\prov{relações: relaciona xadrez}
\prov{proveniência: wiki \textperiodcentered{} xadrez \textperiodcentered{} fato \textperiodcentered{} confiança 1.0 \textperiodcentered{} domínio xadrez \textperiodcentered{} história 2 versões no jornal}

%CRISTAL {"arestas":[{"alvo":"xadrez","confianca":1.0,"epistemico":"fato","origem":"wiki","rel":"parte_de"},{"alvo":"estrategia_posicional","confianca":1.0,"epistemico":"fato","origem":"wiki","rel":"relaciona"}],"confianca":1.0,"contraexemplos":[],"descricao":"As correntes que o jogo atravessou: a ROMÂNTICA (ataque e sacrifício acima de tudo), a POSICIONAL (Steinitz/Capablanca), a HIPERMODERNA (controlar o centro à distância), a SOVIÉTICA (ciência e preparação) e a era dos MOTORES (verdade calculada). A história como diálogo de idéias.","epistemico":"inferencia","exemplos":[],"id":"escola_de_xadrez_estilos","memoria":"semantica","meta":{"dominio":"xadrez","fonte":"história do xadrez"},"origem":"wiki","palavras_chave":[],"sinonimos":[],"tipo":"conceito","titulo":"Estilos e Escolas de Xadrez"}
\section{Estilos e Escolas de Xadrez}
As correntes que o jogo atravessou: a ROMÂNTICA (ataque e sacrifício acima de tudo), a POSICIONAL (Steinitz/Capablanca), a HIPERMODERNA (controlar o centro à distância), a SOVIÉTICA (ciência e preparação) e a era dos MOTORES (verdade calculada). A história como diálogo de idéias.
\prov{relações: parte\_de xadrez; relaciona estrategia\_posicional}
\prov{proveniência: wiki \textperiodcentered{} história do xadrez \textperiodcentered{} inferencia \textperiodcentered{} confiança 1.0 \textperiodcentered{} domínio xadrez \textperiodcentered{} história 3 versões no jornal}

%CRISTAL {"arestas":[{"alvo":"teoria_dos_jogos","confianca":1.0,"epistemico":"fato","origem":"wiki","rel":"usa"}],"confianca":1.0,"contraexemplos":[],"descricao":"A arte de coordenar meios para fins sob oposição inteligente — decidir ONDE e QUANDO agir para vencer o conflito como um todo, não a batalha isolada. Distingue-se da tática (o como local) e da logística (o sustento).","epistemico":"inferencia","exemplos":[],"id":"estrategia","memoria":"semantica","meta":{"dominio":"estrategia","fonte":"estratégia"},"origem":"wiki","palavras_chave":[],"sinonimos":[],"tipo":"conceito","titulo":"Estratégia"}
\section{Estratégia}
A arte de coordenar meios para fins sob oposição inteligente — decidir ONDE e QUANDO agir para vencer o conflito como um todo, não a batalha isolada. Distingue-se da tática (o como local) e da logística (o sustento).
\prov{relações: usa teoria\_dos\_jogos}
\prov{proveniência: wiki \textperiodcentered{} estratégia \textperiodcentered{} inferencia \textperiodcentered{} confiança 1.0 \textperiodcentered{} domínio estrategia \textperiodcentered{} história 2 versões no jornal}

%CRISTAL {"arestas":[{"alvo":"estrategia","confianca":1.0,"epistemico":"fato","origem":"wiki","rel":"especializa"}],"confianca":1.0,"contraexemplos":[],"descricao":"Como uma organização cria e sustenta vantagem num mercado competitivo — escolher um posicionamento único e coerente, não 'ser melhor em tudo'. A guerra transposta para o mercado.","epistemico":"inferencia","exemplos":[],"id":"estrategia_empresarial","memoria":"semantica","meta":{"autor":"Porter","dominio":"estrategia","fonte":"estratégia"},"origem":"wiki","palavras_chave":[],"sinonimos":[],"tipo":"conceito","titulo":"Estratégia Empresarial"}
\section{Estratégia Empresarial}
Como uma organização cria e sustenta vantagem num mercado competitivo — escolher um posicionamento único e coerente, não 'ser melhor em tudo'. A guerra transposta para o mercado.
\prov{relações: especializa estrategia}
\prov{proveniência: wiki \textperiodcentered{} estratégia \textperiodcentered{} inferencia \textperiodcentered{} confiança 1.0 \textperiodcentered{} domínio estrategia \textperiodcentered{} história 2 versões no jornal}

%CRISTAL {"arestas":[{"alvo":"estrategia","confianca":1.0,"epistemico":"fato","origem":"wiki","rel":"especializa"}],"confianca":1.0,"contraexemplos":[],"descricao":"Vencer desequilibrando o inimigo antes do choque — atacar onde ele não espera, pela linha de menor resistência e menor expectativa. O ataque frontal é o último recurso, não o primeiro.","epistemico":"inferencia","exemplos":[],"id":"estrategia_indireta","memoria":"semantica","meta":{"autor":"Liddell Hart","dominio":"estrategia","fonte":"estratégia"},"origem":"wiki","palavras_chave":[],"sinonimos":[],"tipo":"conceito","titulo":"Aproximação Indireta (Liddell Hart)"}
\section{Aproximação Indireta (Liddell Hart)}
Vencer desequilibrando o inimigo antes do choque — atacar onde ele não espera, pela linha de menor resistência e menor expectativa. O ataque frontal é o último recurso, não o primeiro.
\prov{relações: especializa estrategia}
\prov{proveniência: wiki \textperiodcentered{} estratégia \textperiodcentered{} inferencia \textperiodcentered{} confiança 1.0 \textperiodcentered{} domínio estrategia \textperiodcentered{} história 2 versões no jornal}

%CRISTAL {"arestas":[{"alvo":"xadrez","confianca":1.0,"epistemico":"fato","origem":"wiki","rel":"parte_de"},{"alvo":"estrategia","confianca":1.0,"epistemico":"fato","origem":"wiki","rel":"especializa"}],"confianca":1.0,"contraexemplos":[],"descricao":"O jogo de longo prazo: acumular pequenas vantagens duradouras — estrutura, casas, colunas, atividade das peças — até que virem algo concreto. Steinitz a fundou; Capablanca a fez parecer fácil.","epistemico":"inferencia","exemplos":[],"id":"estrategia_posicional","memoria":"semantica","meta":{"dominio":"xadrez","fonte":"xadrez"},"origem":"wiki","palavras_chave":[],"sinonimos":[],"tipo":"conceito","titulo":"Estratégia Posicional"}
\section{Estratégia Posicional}
O jogo de longo prazo: acumular pequenas vantagens duradouras — estrutura, casas, colunas, atividade das peças — até que virem algo concreto. Steinitz a fundou; Capablanca a fez parecer fácil.
\prov{relações: parte\_de xadrez; especializa estrategia}
\prov{proveniência: wiki \textperiodcentered{} xadrez \textperiodcentered{} inferencia \textperiodcentered{} confiança 1.0 \textperiodcentered{} domínio xadrez \textperiodcentered{} história 3 versões no jornal}

%CRISTAL {"arestas":[{"alvo":"estrategia","confianca":1.0,"epistemico":"fato","origem":"wiki","rel":"parte_de"}],"confianca":1.0,"contraexemplos":[],"descricao":"Estratégia é ganhar a guerra; tática é ganhar a batalha. 'Estratégia sem tática é o caminho mais lento à vitória; tática sem estratégia é o ruído antes da derrota' (atribuído a Sun Tzu). Níveis, não sinônimos.","epistemico":"inferencia","exemplos":[],"id":"estrategia_vs_tatica","memoria":"semantica","meta":{"dominio":"estrategia","fonte":"estratégia"},"origem":"wiki","palavras_chave":[],"sinonimos":[],"tipo":"conceito","titulo":"Estratégia vs Tática"}
\section{Estratégia vs Tática}
Estratégia é ganhar a guerra; tática é ganhar a batalha. 'Estratégia sem tática é o caminho mais lento à vitória; tática sem estratégia é o ruído antes da derrota' (atribuído a Sun Tzu). Níveis, não sinônimos.
\prov{relações: parte\_de estrategia}
\prov{proveniência: wiki \textperiodcentered{} estratégia \textperiodcentered{} inferencia \textperiodcentered{} confiança 1.0 \textperiodcentered{} domínio estrategia \textperiodcentered{} história 2 versões no jornal}

%CRISTAL {"arestas":[{"alvo":"estrategia_posicional","confianca":1.0,"epistemico":"fato","origem":"wiki","rel":"usa"}],"confianca":1.0,"contraexemplos":[],"descricao":"O esqueleto da posição: peões dobrados, isolados, passados, correntes. Como os peões (que não recuam) estão dispostos define planos, casas fracas e o caráter da partida. 'A alma do xadrez' (Philidor).","epistemico":"inferencia","exemplos":[],"id":"estrutura_de_peoes","memoria":"semantica","meta":{"dominio":"xadrez","fonte":"xadrez"},"origem":"wiki","palavras_chave":[],"sinonimos":[],"tipo":"conceito","titulo":"Estrutura de Peões"}
\section{Estrutura de Peões}
O esqueleto da posição: peões dobrados, isolados, passados, correntes. Como os peões (que não recuam) estão dispostos define planos, casas fracas e o caráter da partida. 'A alma do xadrez' (Philidor).
\prov{relações: usa estrategia\_posicional}
\prov{proveniência: wiki \textperiodcentered{} xadrez \textperiodcentered{} inferencia \textperiodcentered{} confiança 1.0 \textperiodcentered{} domínio xadrez \textperiodcentered{} história 2 versões no jornal}

%CRISTAL {"arestas":[{"alvo":"criacao_livre_fertilidade","confianca":1.0,"epistemico":"fato","origem":"wiki","rel":"usa"},{"alvo":"estetica_do_fluir_gentil","confianca":1.0,"epistemico":"fato","origem":"wiki","rel":"relaciona"},{"alvo":"borda_critica_vazio","confianca":1.0,"epistemico":"fato","origem":"wiki","rel":"relaciona"}],"confianca":1.0,"contraexemplos":[],"descricao":"Lopes: boas teorias matemáticas 'nascem, crescem, reproduzem-se e não morrem' — são organismos vivos; e ao erguer uma estrutura nova sobre outra (ℝ→ℂ) deve-se sacrificar propriedades caras da anterior, senão a construção afunda. Da renúncia vêm a beleza e a aplicabilidade.","epistemico":"inferencia","exemplos":[],"id":"estruturas_vivas_sacrificio","memoria":"semantica","meta":{"autor":"Gentil Lopes","dominio":"mestres","fonte":""},"origem":"wiki","palavras_chave":[],"sinonimos":[],"tipo":"conceito","titulo":"Estruturas Vivas, Criar é Sacrificar (Lopes)"}
\section{Estruturas Vivas, Criar é Sacrificar (Lopes)}
Lopes: boas teorias matemáticas 'nascem, crescem, reproduzem-se e não morrem' — são organismos vivos; e ao erguer uma estrutura nova sobre outra (\ensuremath{\mathbb{R}}\ensuremath{\to}\ensuremath{\mathbb{C}}) deve-se sacrificar propriedades caras da anterior, senão a construção afunda. Da renúncia vêm a beleza e a aplicabilidade.
\prov{relações: usa criacao\_livre\_fertilidade; relaciona estetica\_do\_fluir\_gentil; relaciona borda\_critica\_vazio}
\prov{proveniência: wiki \textperiodcentered{} inferencia \textperiodcentered{} confiança 1.0 \textperiodcentered{} domínio mestres \textperiodcentered{} história 5 versões no jornal}

%CRISTAL {"arestas":[{"alvo":"xadrez","confianca":1.0,"epistemico":"fato","origem":"wiki","rel":"parte_de"}],"confianca":1.0,"contraexemplos":[],"descricao":"A fase com poucas peças, onde o rei vira atacante e a técnica exata reina — muitos finais estão matematicamente RESOLVIDOS (tablebases). 'Para melhorar, estude os finais primeiro' (Capablanca).","epistemico":"fato","exemplos":[],"id":"final_de_xadrez","memoria":"semantica","meta":{"dominio":"xadrez","fonte":"xadrez"},"origem":"wiki","palavras_chave":[],"sinonimos":[],"tipo":"conceito","titulo":"Finais"}
\section{Finais}
A fase com poucas peças, onde o rei vira atacante e a técnica exata reina — muitos finais estão matematicamente RESOLVIDOS (tablebases). 'Para melhorar, estude os finais primeiro' (Capablanca).
\prov{relações: parte\_de xadrez}
\prov{proveniência: wiki \textperiodcentered{} xadrez \textperiodcentered{} fato \textperiodcentered{} confiança 1.0 \textperiodcentered{} domínio xadrez \textperiodcentered{} história 2 versões no jornal}

%CRISTAL {"arestas":[{"alvo":"abertura_xadrez","confianca":1.0,"epistemico":"fato","origem":"wiki","rel":"especializa"}],"confianca":1.0,"contraexemplos":[],"descricao":"Oferecer um peão do flanco para dominar o centro — quase sempre recusado. A abertura fechada por excelência, campo de batalha dos Campeonatos Mundiais.","epistemico":"fato","exemplos":[],"id":"gambito_da_dama","memoria":"semantica","meta":{"dominio":"xadrez","fonte":"xadrez"},"origem":"wiki","palavras_chave":[],"sinonimos":[],"tipo":"conceito","titulo":"Gambito da Dama (1.d4 d5 2.c4)"}
\section{Gambito da Dama (1.d4 d5 2.c4)}
Oferecer um peão do flanco para dominar o centro — quase sempre recusado. A abertura fechada por excelência, campo de batalha dos Campeonatos Mundiais.
\prov{relações: especializa abertura\_xadrez}
\prov{proveniência: wiki \textperiodcentered{} xadrez \textperiodcentered{} fato \textperiodcentered{} confiança 1.0 \textperiodcentered{} domínio xadrez \textperiodcentered{} história 2 versões no jornal}

%CRISTAL {"arestas":[{"alvo":"tatica_xadrez","confianca":1.0,"epistemico":"fato","origem":"wiki","rel":"especializa"}],"confianca":1.0,"contraexemplos":[],"descricao":"Uma peça ataca duas ou mais ao mesmo tempo — o cavalo é o mestre disso. Ganha material porque o oponente só salva uma. A tática mais elementar e mais comum.","epistemico":"fato","exemplos":[],"id":"garfo_xadrez","memoria":"semantica","meta":{"dominio":"xadrez","fonte":"xadrez"},"origem":"wiki","palavras_chave":[],"sinonimos":[],"tipo":"conceito","titulo":"Garfo (Fork)"}
\section{Garfo (Fork)}
Uma peça ataca duas ou mais ao mesmo tempo — o cavalo é o mestre disso. Ganha material porque o oponente só salva uma. A tática mais elementar e mais comum.
\prov{relações: especializa tatica\_xadrez}
\prov{proveniência: wiki \textperiodcentered{} xadrez \textperiodcentered{} fato \textperiodcentered{} confiança 1.0 \textperiodcentered{} domínio xadrez \textperiodcentered{} história 2 versões no jornal}

%CRISTAL {"fusao":[{"arestas":[{"alvo":"xadrez","confianca":1.0,"epistemico":"fato","origem":"wiki","rel":"parte_de"}],"confianca":1.0,"contraexemplos":[],"descricao":"O mais dominante da história (Campeão 1985–2000), rating recorde por décadas: preparação de abertura devastadora e agressão dinâmica. Perdeu para o Deep Blue em 1997 — o marco homem×máquina.","epistemico":"fato","exemplos":[],"id":"garry_kasparov","memoria":"semantica","meta":{"autor":"Kasparov","dominio":"xadrez","fonte":"história do xadrez"},"origem":"wiki","palavras_chave":[],"sinonimos":[],"tipo":"conceito","titulo":"Garry Kasparov"},{"arestas":[{"alvo":"xadrez","confianca":1.0,"epistemico":"fato","origem":"wiki","rel":"relaciona"}],"confianca":1.0,"contraexemplos":[],"descricao":"Campeão mundial de 1985 a 2000, considerado um dos maiores de todos os tempos — e o humano que perdeu o match histórico para o Deep Blue em 1997.","epistemico":"fato","exemplos":[],"id":"kasparov","memoria":"semantica","meta":{"dominio":"xadrez","fonte":"jogadores"},"origem":"wiki","palavras_chave":[],"sinonimos":[],"tipo":"conceito","titulo":"Garry Kasparov"}],"id":"garry_kasparov","tipo":"conceito"}
\section{Garry Kasparov}
O mais dominante da história (Campeão 1985–2000), rating recorde por décadas: preparação de abertura devastadora e agressão dinâmica. Perdeu para o Deep Blue em 1997 — o marco homem$\times$máquina.
\prov{relações: parte\_de xadrez}
\prov{proveniência: wiki \textperiodcentered{} história do xadrez \textperiodcentered{} fato \textperiodcentered{} confiança 1.0 \textperiodcentered{} domínio xadrez \textperiodcentered{} história 2 versões no jornal}
\prov{a segunda face --- kasparov:}
Campeão mundial de 1985 a 2000, considerado um dos maiores de todos os tempos — e o humano que perdeu o match histórico para o Deep Blue em 1997.
\prov{fusão de curadoria (julgada): absorve kasparov --- as duas faces acima, intactas no registo}

%CRISTAL {"arestas":[{"alvo":"estrategia","confianca":1.0,"epistemico":"fato","origem":"wiki","rel":"generaliza"}],"confianca":1.0,"contraexemplos":[],"descricao":"O nível acima do militar: alinhar TODOS os recursos de uma nação (econômicos, diplomáticos, morais) a um fim de longo prazo — inclusive a paz que se quer no fim. A estratégia dos fins últimos.","epistemico":"inferencia","exemplos":[],"id":"grande_estrategia","memoria":"semantica","meta":{"dominio":"estrategia","fonte":"estratégia"},"origem":"wiki","palavras_chave":[],"sinonimos":[],"tipo":"conceito","titulo":"Grande Estratégia"}
\section{Grande Estratégia}
O nível acima do militar: alinhar TODOS os recursos de uma nação (econômicos, diplomáticos, morais) a um fim de longo prazo — inclusive a paz que se quer no fim. A estratégia dos fins últimos.
\prov{relações: generaliza estrategia}
\prov{proveniência: wiki \textperiodcentered{} estratégia \textperiodcentered{} inferencia \textperiodcentered{} confiança 1.0 \textperiodcentered{} domínio estrategia \textperiodcentered{} história 2 versões no jornal}

%CRISTAL {"arestas":[{"alvo":"li_ritual","confianca":1.0,"epistemico":"fato","origem":"wiki","rel":"usa"},{"alvo":"confucionismo","confianca":1.0,"epistemico":"fato","origem":"wiki","rel":"parte_de"}],"confianca":1.0,"contraexemplos":[],"descricao":"Confúcio: o ideal social, mas 'harmonia sem uniformidade' — 'o cavalheiro harmoniza-se sem ser um eco'. É o fruto mais valioso dos ritos, porém não pode ser buscada sozinha, sem ser regulada pelo li. Consenso, não conformismo.","epistemico":"inferencia","exemplos":[],"id":"he_harmonia","memoria":"semantica","meta":{"autor":"Confúcio","dominio":"mestres","fonte":""},"origem":"wiki","palavras_chave":[],"sinonimos":[],"tipo":"conceito","titulo":"He — Harmonia (sem uniformidade)"}
\section{He — Harmonia (sem uniformidade)}
Confúcio: o ideal social, mas 'harmonia sem uniformidade' — 'o cavalheiro harmoniza-se sem ser um eco'. É o fruto mais valioso dos ritos, porém não pode ser buscada sozinha, sem ser regulada pelo li. Consenso, não conformismo.
\prov{relações: usa li\_ritual; parte\_de confucionismo}
\prov{proveniência: wiki \textperiodcentered{} inferencia \textperiodcentered{} confiança 1.0 \textperiodcentered{} domínio mestres \textperiodcentered{} história 3 versões no jornal}

%CRISTAL {"arestas":[{"alvo":"estrategia_posicional","confianca":1.0,"epistemico":"fato","origem":"wiki","rel":"usa"}],"confianca":1.0,"contraexemplos":[],"descricao":"Quem dita o jogo, força respostas e ataca — o tempo é vantagem. Vale material: muitos gambitos trocam um peão pela iniciativa. Perder a iniciativa é passar a só reagir.","epistemico":"inferencia","exemplos":[],"id":"iniciativa_xadrez","memoria":"semantica","meta":{"dominio":"xadrez","fonte":"xadrez"},"origem":"wiki","palavras_chave":[],"sinonimos":[],"tipo":"conceito","titulo":"Iniciativa"}
\section{Iniciativa}
Quem dita o jogo, força respostas e ataca — o tempo é vantagem. Vale material: muitos gambitos trocam um peão pela iniciativa. Perder a iniciativa é passar a só reagir.
\prov{relações: usa estrategia\_posicional}
\prov{proveniência: wiki \textperiodcentered{} xadrez \textperiodcentered{} inferencia \textperiodcentered{} confiança 1.0 \textperiodcentered{} domínio xadrez \textperiodcentered{} história 2 versões no jornal}

%CRISTAL {"arestas":[{"alvo":"yin_yang","confianca":1.0,"epistemico":"fato","origem":"wiki","rel":"explica"},{"alvo":"matematica_arte_engenharia","confianca":1.0,"epistemico":"fato","origem":"wiki","rel":"usa"},{"alvo":"nao_dualidade_binaria_gentil","confianca":1.0,"epistemico":"fato","origem":"wiki","rel":"usa"},{"alvo":"tao_da_matematica","confianca":1.0,"epistemico":"fato","origem":"wiki","rel":"parte_de"}],"confianca":1.0,"contraexemplos":[],"descricao":"Lopes detecta um isomorfismo entre os trigramas/hexagramas do I Ching e as sequências binárias de sua 'construção de Deus' (yin=0, yang=1): os 8 trigramas são os vértices de um cubo (ℤ³), os 64 hexagramas os de um hipercubo 6-D (2⁶). A lei de mutação do Tao (obscuro↔luminoso) = o intercâmbio yin/yang. Ele matematiza o Oriente.","epistemico":"inferencia","exemplos":[],"id":"isomorfismo_iching_binario","memoria":"semantica","meta":{"autor":"Gentil Lopes","dominio":"mestres","fonte":""},"origem":"wiki","palavras_chave":[],"sinonimos":[],"tipo":"conceito","titulo":"I Ching ↔ Hipercubo Binário (Lopes)"}
\section{I Ching \ensuremath{\leftrightarrow} Hipercubo Binário (Lopes)}
Lopes detecta um isomorfismo entre os trigramas/hexagramas do I Ching e as sequências binárias de sua 'construção de Deus' (yin=0, yang=1): os 8 trigramas são os vértices de um cubo (\ensuremath{\mathbb{Z}}\ensuremath{^{3}}), os 64 hexagramas os de um hipercubo 6-D (2\ensuremath{^{6}}). A lei de mutação do Tao (obscuro\ensuremath{\leftrightarrow}luminoso) = o intercâmbio yin/yang. Ele matematiza o Oriente.
\prov{relações: explica yin\_yang; usa matematica\_arte\_engenharia; usa nao\_dualidade\_binaria\_gentil; parte\_de tao\_da\_matematica}
\prov{proveniência: wiki \textperiodcentered{} inferencia \textperiodcentered{} confiança 1.0 \textperiodcentered{} domínio mestres \textperiodcentered{} história 6 versões no jornal}

%CRISTAL {"arestas":[{"alvo":"final_de_xadrez","confianca":1.0,"epistemico":"fato","origem":"wiki","rel":"relaciona"}],"confianca":1.0,"contraexemplos":[],"descricao":"O cubano (Campeão 1921–27) de clareza sobrenatural nos finais e simplicidade — parecia nunca errar. A 'máquina de xadrez' humana; técnica pura e economia de esforço.","epistemico":"fato","exemplos":[],"id":"jose_capablanca","memoria":"semantica","meta":{"autor":"Capablanca","dominio":"xadrez","fonte":"história do xadrez"},"origem":"wiki","palavras_chave":[],"sinonimos":[],"tipo":"conceito","titulo":"José Raúl Capablanca"}
\section{José Raúl Capablanca}
O cubano (Campeão 1921–27) de clareza sobrenatural nos finais e simplicidade — parecia nunca errar. A 'máquina de xadrez' humana; técnica pura e economia de esforço.
\prov{relações: relaciona final\_de\_xadrez}
\prov{proveniência: wiki \textperiodcentered{} história do xadrez \textperiodcentered{} fato \textperiodcentered{} confiança 1.0 \textperiodcentered{} domínio xadrez \textperiodcentered{} história 2 versões no jornal}

%CRISTAL {"arestas":[{"alvo":"confucionismo","confianca":1.0,"epistemico":"fato","origem":"wiki","rel":"parte_de"},{"alvo":"eudaimonia","confianca":1.0,"epistemico":"fato","origem":"wiki","rel":"relaciona"},{"alvo":"ren","confianca":1.0,"epistemico":"fato","origem":"wiki","rel":"usa"}],"confianca":1.0,"contraexemplos":[],"descricao":"Confúcio: o 'cavalheiro' definido por caráter e cultivo moral, não por nascimento — nunca abandona a benevolência 'nem no espaço de uma refeição'. Contrasta com o homem pequeno (xiaoren): 'o cavalheiro concorda sem ser um eco; o vulgar ecoa sem concordar'.","epistemico":"inferencia","exemplos":[],"id":"junzi","memoria":"semantica","meta":{"autor":"Confúcio","dominio":"mestres","fonte":""},"origem":"wiki","palavras_chave":[],"sinonimos":[],"tipo":"conceito","titulo":"Junzi — o Homem Nobre / Exemplar"}
\section{Junzi — o Homem Nobre / Exemplar}
Confúcio: o 'cavalheiro' definido por caráter e cultivo moral, não por nascimento — nunca abandona a benevolência 'nem no espaço de uma refeição'. Contrasta com o homem pequeno (xiaoren): 'o cavalheiro concorda sem ser um eco; o vulgar ecoa sem concordar'.
\prov{relações: parte\_de confucionismo; relaciona eudaimonia; usa ren}
\prov{proveniência: wiki \textperiodcentered{} inferencia \textperiodcentered{} confiança 1.0 \textperiodcentered{} domínio mestres \textperiodcentered{} história 4 versões no jornal}

%CRISTAL {"arestas":[{"alvo":"vedanta","confianca":1.0,"epistemico":"fato","origem":"wiki","rel":"parte_de"}],"confianca":1.0,"contraexemplos":[],"descricao":"Tradição indiana: karma é a lei de causa-efeito moral pela qual as ações condicionam experiências futuras e alimentam o samsara (ciclo de renascimentos); dharma é o dever/ordem cósmica que a pessoa cumpre.","epistemico":"inferencia","exemplos":[],"id":"karma_dharma","memoria":"semantica","meta":{"autor":"—","dominio":"mestres","fonte":""},"origem":"wiki","palavras_chave":[],"sinonimos":[],"tipo":"conceito","titulo":"Karma e Dharma"}
\section{Karma e Dharma}
Tradição indiana: karma é a lei de causa-efeito moral pela qual as ações condicionam experiências futuras e alimentam o samsara (ciclo de renascimentos); dharma é o dever/ordem cósmica que a pessoa cumpre.
\prov{relações: parte\_de vedanta}
\prov{proveniência: wiki \textperiodcentered{} inferencia \textperiodcentered{} confiança 1.0 \textperiodcentered{} domínio mestres \textperiodcentered{} história 2 versões no jornal}

%CRISTAL {"arestas":[{"alvo":"garry_kasparov","confianca":1.0,"epistemico":"fato","origem":"wiki","rel":"parte_de"},{"alvo":"tatica_xadrez","confianca":1.0,"epistemico":"fato","origem":"wiki","rel":"usa"}],"confianca":1.0,"contraexemplos":[],"descricao":"Talvez a melhor partida de torneio já jogada: um sacrifício de torre e uma caçada ao rei através de todo o tabuleiro, calculada à perfeição. O ápice do cálculo humano.","epistemico":"fato","exemplos":[],"id":"kasparov_topalov_1999","memoria":"semantica","meta":{"autor":"Kasparov","dominio":"xadrez","fonte":"história do xadrez"},"origem":"wiki","palavras_chave":[],"sinonimos":[],"tipo":"conceito","titulo":"Kasparov–Topalov (Wijk aan Zee, 1999)"}
\section{Kasparov–Topalov (Wijk aan Zee, 1999)}
Talvez a melhor partida de torneio já jogada: um sacrifício de torre e uma caçada ao rei através de todo o tabuleiro, calculada à perfeição. O ápice do cálculo humano.
\prov{relações: parte\_de garry\_kasparov; usa tatica\_xadrez}
\prov{proveniência: wiki \textperiodcentered{} história do xadrez \textperiodcentered{} fato \textperiodcentered{} confiança 1.0 \textperiodcentered{} domínio xadrez \textperiodcentered{} história 3 versões no jornal}

%CRISTAL {"arestas":[{"alvo":"budismo","confianca":1.0,"epistemico":"fato","origem":"wiki","rel":"parte_de"},{"alvo":"borda_critica_vazio","confianca":1.0,"epistemico":"fato","origem":"wiki","rel":"relaciona"}],"confianca":1.0,"contraexemplos":[],"descricao":"Zen (Rinzai): o koan é um enigma paradoxal ('o som de uma só mão') que curto-circuita o pensamento racional-conceitual para provocar o satori/kensho — o despertar súbito, o ver diretamente a própria natureza, ruptura da percepção dualista.","epistemico":"inferencia","exemplos":[],"id":"koan_satori","memoria":"semantica","meta":{"autor":"Zen Rinzai","dominio":"mestres","fonte":""},"origem":"wiki","palavras_chave":[],"sinonimos":[],"tipo":"conceito","titulo":"Koan e Satori (Zen)"}
\section{Koan e Satori (Zen)}
Zen (Rinzai): o koan é um enigma paradoxal ('o som de uma só mão') que curto-circuita o pensamento racional-conceitual para provocar o satori/kensho — o despertar súbito, o ver diretamente a própria natureza, ruptura da percepção dualista.
\prov{relações: parte\_de budismo; relaciona borda\_critica\_vazio}
\prov{proveniência: wiki \textperiodcentered{} inferencia \textperiodcentered{} confiança 1.0 \textperiodcentered{} domínio mestres \textperiodcentered{} história 3 versões no jornal}

%CRISTAL {"arestas":[{"alvo":"confucionismo","confianca":1.0,"epistemico":"fato","origem":"wiki","rel":"parte_de"},{"alvo":"contrato_social","confianca":1.0,"epistemico":"fato","origem":"wiki","rel":"relaciona"}],"confianca":1.0,"contraexemplos":[],"descricao":"Confúcio: o corpo de ritos e formas corretas de conduta que ordenam a sociedade e as relações — mas só vale com participação sincera; a mera forma vazia 'reflete um mal espiritual'. A benevolência (ren) realiza-se através do li.","epistemico":"inferencia","exemplos":[],"id":"li_ritual","memoria":"semantica","meta":{"autor":"Confúcio","dominio":"mestres","fonte":""},"origem":"wiki","palavras_chave":[],"sinonimos":[],"tipo":"conceito","titulo":"Li — Ritual / Propriedade"}
\section{Li — Ritual / Propriedade}
Confúcio: o corpo de ritos e formas corretas de conduta que ordenam a sociedade e as relações — mas só vale com participação sincera; a mera forma vazia 'reflete um mal espiritual'. A benevolência (ren) realiza-se através do li.
\prov{relações: parte\_de confucionismo; relaciona contrato\_social}
\prov{proveniência: wiki \textperiodcentered{} inferencia \textperiodcentered{} confiança 1.0 \textperiodcentered{} domínio mestres \textperiodcentered{} história 3 versões no jornal}

%CRISTAL {"arestas":[{"alvo":"xadrez","confianca":1.0,"epistemico":"fato","origem":"wiki","rel":"parte_de"}],"confianca":1.0,"contraexemplos":[],"descricao":"O norueguês (Campeão 2013–2023), maior rating de todos os tempos: universalidade e técnica de espremer posições 'iguais' até a vitória. A era da preparação com motores.","epistemico":"fato","exemplos":[],"id":"magnus_carlsen","memoria":"semantica","meta":{"autor":"Carlsen","dominio":"xadrez","fonte":"história do xadrez"},"origem":"wiki","palavras_chave":[],"sinonimos":[],"tipo":"conceito","titulo":"Magnus Carlsen"}
\section{Magnus Carlsen}
O norueguês (Campeão 2013–2023), maior rating de todos os tempos: universalidade e técnica de espremer posições 'iguais' até a vitória. A era da preparação com motores.
\prov{relações: parte\_de xadrez}
\prov{proveniência: wiki \textperiodcentered{} história do xadrez \textperiodcentered{} fato \textperiodcentered{} confiança 1.0 \textperiodcentered{} domínio xadrez \textperiodcentered{} história 2 versões no jornal}

%CRISTAL {"arestas":[{"alvo":"vedanta","confianca":1.0,"epistemico":"fato","origem":"wiki","rel":"parte_de"},{"alvo":"nao_dualidade","confianca":1.0,"epistemico":"fato","origem":"wiki","rel":"relaciona"},{"alvo":"colapso_observador","confianca":1.0,"epistemico":"fato","origem":"wiki","rel":"relaciona"}],"confianca":1.0,"contraexemplos":[],"descricao":"Advaita Vedanta: o poder que faz a realidade una (brahman) aparecer como o mundo múltiplo e as distinções sujeito-objeto — a superimposição nascida da ignorância (avidya) que vela a não-dualidade.","epistemico":"inferencia","exemplos":[],"id":"maya","memoria":"semantica","meta":{"autor":"Shankara","dominio":"mestres","fonte":""},"origem":"wiki","palavras_chave":[],"sinonimos":[],"tipo":"conceito","titulo":"Maya — a Ilusão"}
\section{Maya — a Ilusão}
Advaita Vedanta: o poder que faz a realidade una (brahman) aparecer como o mundo múltiplo e as distinções sujeito-objeto — a superimposição nascida da ignorância (avidya) que vela a não-dualidade.
\prov{relações: parte\_de vedanta; relaciona nao\_dualidade; relaciona colapso\_observador}
\prov{proveniência: wiki \textperiodcentered{} inferencia \textperiodcentered{} confiança 1.0 \textperiodcentered{} domínio mestres \textperiodcentered{} história 4 versões no jornal}

%CRISTAL {"arestas":[{"alvo":"xadrez","confianca":1.0,"epistemico":"fato","origem":"wiki","rel":"parte_de"}],"confianca":1.0,"contraexemplos":[],"descricao":"O 'patriarca' e engenheiro do xadrez científico: preparação metódica, análise profunda de aberturas e adiamento. Fundou a escola soviética que dominou o século e formou Karpov e Kasparov.","epistemico":"fato","exemplos":[],"id":"mikhail_botvinnik","memoria":"semantica","meta":{"autor":"Botvinnik","dominio":"xadrez","fonte":"história do xadrez"},"origem":"wiki","palavras_chave":[],"sinonimos":[],"tipo":"conceito","titulo":"Mikhail Botvinnik e a Escola Soviética"}
\section{Mikhail Botvinnik e a Escola Soviética}
O 'patriarca' e engenheiro do xadrez científico: preparação metódica, análise profunda de aberturas e adiamento. Fundou a escola soviética que dominou o século e formou Karpov e Kasparov.
\prov{relações: parte\_de xadrez}
\prov{proveniência: wiki \textperiodcentered{} história do xadrez \textperiodcentered{} fato \textperiodcentered{} confiança 1.0 \textperiodcentered{} domínio xadrez \textperiodcentered{} história 2 versões no jornal}

%CRISTAL {"arestas":[{"alvo":"sacrificio_xadrez","confianca":1.0,"epistemico":"fato","origem":"wiki","rel":"relaciona"}],"confianca":1.0,"contraexemplos":[],"descricao":"O 'Mágico de Riga': sacrifícios intuitivos e caos calculado que ninguém conseguia refutar no tabuleiro. O ataque como arte psicológica — arrastava o oponente à floresta onde 2+2=5.","epistemico":"fato","exemplos":[],"id":"mikhail_tal","memoria":"semantica","meta":{"autor":"Tal","dominio":"xadrez","fonte":"história do xadrez"},"origem":"wiki","palavras_chave":[],"sinonimos":[],"tipo":"conceito","titulo":"Mikhail Tal"}
\section{Mikhail Tal}
O 'Mágico de Riga': sacrifícios intuitivos e caos calculado que ninguém conseguia refutar no tabuleiro. O ataque como arte psicológica — arrastava o oponente à floresta onde 2+2=5.
\prov{relações: relaciona sacrificio\_xadrez}
\prov{proveniência: wiki \textperiodcentered{} história do xadrez \textperiodcentered{} fato \textperiodcentered{} confiança 1.0 \textperiodcentered{} domínio xadrez \textperiodcentered{} história 2 versões no jornal}

%CRISTAL {"arestas":[{"alvo":"minimax","confianca":1.0,"epistemico":"fato","origem":"wiki","rel":"especializa"},{"alvo":"forma_extensiva","confianca":1.0,"epistemico":"fato","origem":"wiki","rel":"usa"}],"confianca":1.0,"contraexemplos":[],"descricao":"O minimax da teoria dos jogos, viabilizado: a poda alfa-beta descarta ramos que não podem mudar a decisão, cortando a árvore sem perder a jogada ótima. O coração dos motores clássicos.","epistemico":"fato","exemplos":[],"id":"minimax_alfabeta","memoria":"semantica","meta":{"dominio":"xadrez","fonte":"xadrez"},"origem":"wiki","palavras_chave":[],"sinonimos":[],"tipo":"conceito","titulo":"Minimax com Poda Alfa-Beta"}
\section{Minimax com Poda Alfa-Beta}
O minimax da teoria dos jogos, viabilizado: a poda alfa-beta descarta ramos que não podem mudar a decisão, cortando a árvore sem perder a jogada ótima. O coração dos motores clássicos.
\prov{relações: especializa minimax; usa forma\_extensiva}
\prov{proveniência: wiki \textperiodcentered{} xadrez \textperiodcentered{} fato \textperiodcentered{} confiança 1.0 \textperiodcentered{} domínio xadrez \textperiodcentered{} história 3 versões no jornal}

%CRISTAL {"arestas":[{"alvo":"xadrez","confianca":1.0,"epistemico":"fato","origem":"wiki","rel":"usa"},{"alvo":"minimax_alfabeta","confianca":1.0,"epistemico":"fato","origem":"wiki","rel":"usa"},{"alvo":"inteligencia_artificial","confianca":1.0,"epistemico":"fato","origem":"wiki","rel":"parte_de"}],"confianca":1.0,"contraexemplos":[],"descricao":"O programa que joga: busca na árvore (minimax + poda alfa-beta) com função de avaliação, ou, no AlphaZero, rede neural + busca Monte Carlo aprendida por auto-jogo. Hoje muito acima do melhor humano.","epistemico":"fato","exemplos":[],"id":"motor_de_xadrez","memoria":"semantica","meta":{"dominio":"xadrez","fonte":"xadrez"},"origem":"wiki","palavras_chave":[],"sinonimos":[],"tipo":"conceito","titulo":"Motor de Xadrez"}
\section{Motor de Xadrez}
O programa que joga: busca na árvore (minimax + poda alfa-beta) com função de avaliação, ou, no AlphaZero, rede neural + busca Monte Carlo aprendida por auto-jogo. Hoje muito acima do melhor humano.
\prov{relações: usa xadrez; usa minimax\_alfabeta; parte\_de inteligencia\_artificial}
\prov{proveniência: wiki \textperiodcentered{} xadrez \textperiodcentered{} fato \textperiodcentered{} confiança 1.0 \textperiodcentered{} domínio xadrez \textperiodcentered{} história 4 versões no jornal}

%CRISTAL {"arestas":[{"alvo":"budismo","confianca":1.0,"epistemico":"fato","origem":"wiki","rel":"parte_de"},{"alvo":"contemplacao","confianca":1.0,"epistemico":"fato","origem":"wiki","rel":"usa"},{"alvo":"meditacao_cura_da_mente_gentil","confianca":1.0,"epistemico":"fato","origem":"wiki","rel":"relaciona"}],"confianca":1.0,"contraexemplos":[],"descricao":"Zen (Soto): mushin é a mente sem apego/interferência egoica; para Dogen, a prática é 'só sentar' (shikantaza) como prática-realização — o zazen não é meio para a iluminação, mas a própria iluminação atualizada, dissolvendo a dualidade meio/fim.","epistemico":"inferencia","exemplos":[],"id":"mushin_zazen","memoria":"semantica","meta":{"autor":"Dogen","dominio":"mestres","fonte":""},"origem":"wiki","palavras_chave":[],"sinonimos":[],"tipo":"conceito","titulo":"Mushin e Zazen (Dogen)"}
\section{Mushin e Zazen (Dogen)}
Zen (Soto): mushin é a mente sem apego/interferência egoica; para Dogen, a prática é 'só sentar' (shikantaza) como prática-realização — o zazen não é meio para a iluminação, mas a própria iluminação atualizada, dissolvendo a dualidade meio/fim.
\prov{relações: parte\_de budismo; usa contemplacao; relaciona meditacao\_cura\_da\_mente\_gentil}
\prov{proveniência: wiki \textperiodcentered{} inferencia \textperiodcentered{} confiança 1.0 \textperiodcentered{} domínio mestres \textperiodcentered{} história 4 versões no jornal}

%CRISTAL {"arestas":[{"alvo":"nao_dualidade","confianca":1.0,"epistemico":"fato","origem":"wiki","rel":"usa"},{"alvo":"yin_yang","confianca":1.0,"epistemico":"fato","origem":"wiki","rel":"explica"},{"alvo":"deus_como_vazio_gentil","confianca":1.0,"epistemico":"fato","origem":"wiki","rel":"usa"},{"alvo":"simbolo_vs_significado","confianca":1.0,"epistemico":"fato","origem":"wiki","rel":"relaciona"}],"confianca":1.0,"contraexemplos":[],"descricao":"Lopes: Deus absoluto = Vazio =; Deus manifesto = dualidade = 0,1^∞ − 0. Toda sequência tem sua dual (bem/mal), que só existem 'da perspectiva humana, não da Divina' — o mesmo par yin/yang, cujos opostos complementares escondem uma única energia (o Tao).","epistemico":"inferencia","exemplos":[],"id":"nao_dualidade_binaria_gentil","memoria":"semantica","meta":{"autor":"Gentil Lopes","dominio":"mestres","fonte":""},"origem":"wiki","palavras_chave":[],"sinonimos":[],"tipo":"conceito","titulo":"A Não-Dualidade Binária (Lopes)"}
\section{A Não-Dualidade Binária (Lopes)}
Lopes: Deus absoluto = Vazio =; Deus manifesto = dualidade = 0,1\^{}\ensuremath{\infty} \ensuremath{-} 0. Toda sequência tem sua dual (bem/mal), que só existem 'da perspectiva humana, não da Divina' — o mesmo par yin/yang, cujos opostos complementares escondem uma única energia (o Tao).
\prov{relações: usa nao\_dualidade; explica yin\_yang; usa deus\_como\_vazio\_gentil; relaciona simbolo\_vs\_significado}
\prov{proveniência: wiki \textperiodcentered{} inferencia \textperiodcentered{} confiança 1.0 \textperiodcentered{} domínio mestres \textperiodcentered{} história 7 versões no jornal}

%CRISTAL {"arestas":[{"alvo":"vedanta","confianca":1.0,"epistemico":"fato","origem":"wiki","rel":"parte_de"},{"alvo":"contemplacao","confianca":1.0,"epistemico":"fato","origem":"wiki","rel":"usa"},{"alvo":"atman_brahman","confianca":1.0,"epistemico":"fato","origem":"wiki","rel":"usa"}],"confianca":1.0,"contraexemplos":[],"descricao":"Upanishads / Shankara: o método apofático que conhece o Absoluto removendo toda identificação com corpo, mente e mundo; a negação não é fim em si — aponta para a identidade atman-brahman. A teologia negativa do Oriente.","epistemico":"inferencia","exemplos":[],"id":"neti_neti","memoria":"semantica","meta":{"autor":"Shankara","dominio":"mestres","fonte":""},"origem":"wiki","palavras_chave":[],"sinonimos":[],"tipo":"conceito","titulo":"Neti Neti — 'nem isto, nem aquilo'"}
\section{Neti Neti — 'nem isto, nem aquilo'}
Upanishads / Shankara: o método apofático que conhece o Absoluto removendo toda identificação com corpo, mente e mundo; a negação não é fim em si — aponta para a identidade atman-brahman. A teologia negativa do Oriente.
\prov{relações: parte\_de vedanta; usa contemplacao; usa atman\_brahman}
\prov{proveniência: wiki \textperiodcentered{} inferencia \textperiodcentered{} confiança 1.0 \textperiodcentered{} domínio mestres \textperiodcentered{} história 4 versões no jornal}

%CRISTAL {"arestas":[{"alvo":"clausewitz_da_guerra","confianca":1.0,"epistemico":"fato","origem":"wiki","rel":"parte_de"}],"confianca":1.0,"contraexemplos":[],"descricao":"Névoa: agir com informação incompleta e falsa. Fricção: mil pequenos atritos que fazem o simples ser difícil. Por que planos raramente sobrevivem ao contato — a estratégia real convive com o desconhecido.","epistemico":"inferencia","exemplos":[],"id":"nevoa_e_friccao","memoria":"semantica","meta":{"dominio":"estrategia","fonte":"estratégia"},"origem":"wiki","palavras_chave":[],"sinonimos":[],"tipo":"conceito","titulo":"Névoa e Fricção da Guerra"}
\section{Névoa e Fricção da Guerra}
Névoa: agir com informação incompleta e falsa. Fricção: mil pequenos atritos que fazem o simples ser difícil. Por que planos raramente sobrevivem ao contato — a estratégia real convive com o desconhecido.
\prov{relações: parte\_de clausewitz\_da\_guerra}
\prov{proveniência: wiki \textperiodcentered{} estratégia \textperiodcentered{} inferencia \textperiodcentered{} confiança 1.0 \textperiodcentered{} domínio estrategia \textperiodcentered{} história 2 versões no jornal}

%CRISTAL {"arestas":[{"alvo":"budismo","confianca":1.0,"epistemico":"fato","origem":"wiki","rel":"parte_de"},{"alvo":"meditacao_cura_da_mente_gentil","confianca":1.0,"epistemico":"fato","origem":"wiki","rel":"relaciona"},{"alvo":"quatro_nobres_verdades","confianca":1.0,"epistemico":"fato","origem":"wiki","rel":"usa"},{"alvo":"karma_dharma","confianca":1.0,"epistemico":"fato","origem":"wiki","rel":"relaciona"}],"confianca":1.0,"contraexemplos":[],"descricao":"Buda: a 'extinção' dos fogos do desejo, aversão e ilusão — cessação do apego e, com ele, do sofrimento (dukkha). No Mahayana, não uma fuga do mundo, mas um modo iluminado de estar-no-mundo aqui e agora. Inclui moksha (a libertação do samsara).","epistemico":"inferencia","exemplos":[],"id":"nirvana","memoria":"semantica","meta":{"autor":"Buda","dominio":"mestres","fonte":""},"origem":"wiki","palavras_chave":[],"sinonimos":[],"tipo":"conceito","titulo":"Nirvana — a Extinção do Apego"}
\section{Nirvana — a Extinção do Apego}
Buda: a 'extinção' dos fogos do desejo, aversão e ilusão — cessação do apego e, com ele, do sofrimento (dukkha). No Mahayana, não uma fuga do mundo, mas um modo iluminado de estar-no-mundo aqui e agora. Inclui moksha (a libertação do samsara).
\prov{relações: parte\_de budismo; relaciona meditacao\_cura\_da\_mente\_gentil; usa quatro\_nobres\_verdades; relaciona karma\_dharma}
\prov{proveniência: wiki \textperiodcentered{} inferencia \textperiodcentered{} confiança 1.0 \textperiodcentered{} domínio mestres \textperiodcentered{} história 5 versões no jornal}

%CRISTAL {"arestas":[{"alvo":"xadrez","confianca":1.0,"epistemico":"fato","origem":"wiki","rel":"parte_de"}],"confianca":1.0,"contraexemplos":[],"descricao":"O sistema para escrever lances: colunas a–h, linhas 1–8, letra da peça (C=cavalo, B=bispo…). 'Cf3', 'e4', 'O-O' (roque). A linguagem que torna toda partida reproduzível e estudável.","epistemico":"fato","exemplos":[],"id":"notacao_algebrica","memoria":"semantica","meta":{"dominio":"xadrez","fonte":"xadrez"},"origem":"wiki","palavras_chave":[],"sinonimos":[],"tipo":"conceito","titulo":"Notação Algébrica"}
\section{Notação Algébrica}
O sistema para escrever lances: colunas a–h, linhas 1–8, letra da peça (C=cavalo, B=bispo…). 'Cf3', 'e4', 'O-O' (roque). A linguagem que torna toda partida reproduzível e estudável.
\prov{relações: parte\_de xadrez}
\prov{proveniência: wiki \textperiodcentered{} xadrez \textperiodcentered{} fato \textperiodcentered{} confiança 1.0 \textperiodcentered{} domínio xadrez \textperiodcentered{} história 2 versões no jornal}

%CRISTAL {"arestas":[{"alvo":"xadrez","confianca":1.0,"epistemico":"fato","origem":"wiki","rel":"explica"}],"confianca":1.0,"contraexemplos":[],"descricao":"A estimativa da complexidade do xadrez: ~10¹²⁰ partidas possíveis — mais que átomos no universo observável. Por que a força bruta pura é impossível e a poda/heurística é obrigatória.","epistemico":"fato","exemplos":[],"id":"numero_de_shannon","memoria":"semantica","meta":{"autor":"Shannon","dominio":"xadrez","fonte":"xadrez"},"origem":"wiki","palavras_chave":[],"sinonimos":[],"tipo":"conceito","titulo":"Número de Shannon"}
\section{Número de Shannon}
A estimativa da complexidade do xadrez: \~{}10\ensuremath{^{120}} partidas possíveis — mais que átomos no universo observável. Por que a força bruta pura é impossível e a poda/heurística é obrigatória.
\prov{relações: explica xadrez}
\prov{proveniência: wiki \textperiodcentered{} xadrez \textperiodcentered{} fato \textperiodcentered{} confiança 1.0 \textperiodcentered{} domínio xadrez \textperiodcentered{} história 2 versões no jornal}

%CRISTAL {"arestas":[{"alvo":"estrategia","confianca":1.0,"epistemico":"fato","origem":"wiki","rel":"usa"}],"confianca":1.0,"contraexemplos":[],"descricao":"Observar–Orientar–Decidir–Agir, mais rápido que o inimigo: quem fecha o ciclo primeiro entra DENTRO do ciclo do outro e o desorienta. A estratégia como velocidade de adaptação.","epistemico":"inferencia","exemplos":[],"id":"ooda_loop","memoria":"semantica","meta":{"autor":"John Boyd","dominio":"estrategia","fonte":"estratégia"},"origem":"wiki","palavras_chave":[],"sinonimos":[],"tipo":"conceito","titulo":"Ciclo OODA (Boyd)"}
\section{Ciclo OODA (Boyd)}
Observar–Orientar–Decidir–Agir, mais rápido que o inimigo: quem fecha o ciclo primeiro entra DENTRO do ciclo do outro e o desorienta. A estratégia como velocidade de adaptação.
\prov{relações: usa estrategia}
\prov{proveniência: wiki \textperiodcentered{} estratégia \textperiodcentered{} inferencia \textperiodcentered{} confiança 1.0 \textperiodcentered{} domínio estrategia \textperiodcentered{} história 2 versões no jornal}

%CRISTAL {"arestas":[{"alvo":"final_de_xadrez","confianca":1.0,"epistemico":"fato","origem":"wiki","rel":"usa"}],"confianca":1.0,"contraexemplos":[],"descricao":"Reis frente a frente com uma casa entre eles: quem NÃO tem a vez de jogar tem a oposição e controla as casas-chave. A geometria decisiva dos finais de rei e peão.","epistemico":"fato","exemplos":[],"id":"oposicao_final","memoria":"semantica","meta":{"dominio":"xadrez","fonte":"xadrez"},"origem":"wiki","palavras_chave":[],"sinonimos":[],"tipo":"conceito","titulo":"Oposição (finais de rei e peão)"}
\section{Oposição (finais de rei e peão)}
Reis frente a frente com uma casa entre eles: quem NÃO tem a vez de jogar tem a oposição e controla as casas-chave. A geometria decisiva dos finais de rei e peão.
\prov{relações: usa final\_de\_xadrez}
\prov{proveniência: wiki \textperiodcentered{} xadrez \textperiodcentered{} fato \textperiodcentered{} confiança 1.0 \textperiodcentered{} domínio xadrez \textperiodcentered{} história 2 versões no jornal}

%CRISTAL {"arestas":[{"alvo":"orbita_de_partida","confianca":1.0,"epistemico":"fato","origem":"wiki","rel":"eh_um"},{"alvo":"abertura_xadrez","confianca":1.0,"epistemico":"fato","origem":"wiki","rel":"usa"},{"alvo":"xeque_mate","confianca":1.0,"epistemico":"fato","origem":"wiki","rel":"relaciona"},{"alvo":"sacrificio_xadrez","confianca":1.0,"epistemico":"fato","origem":"wiki","rel":"usa"},{"alvo":"alekhine","confianca":1.0,"epistemico":"fato","origem":"wiki","rel":"relaciona"},{"alvo":"lyapunov_da_orbita","confianca":1.0,"epistemico":"fato","origem":"wiki","rel":"usa"}],"confianca":1.0,"contraexemplos":[],"descricao":"Sacrifício de dama (Dxe6+!!) que rasga o roque e o bispo dá mate em Bg6# — uma miniatura de 11 lances do campeão. Jogadores: Alexander Alekhine × Vasic (1931). A órbita percorre 22 posições e termina em xeque-mate após 21 lances.","epistemico":"fato","exemplos":[],"id":"partida_alekhine_vasic","memoria":"semantica","meta":{"abertura":"Defesa Francesa","ano":1931,"classe_orbita":"borda","dominio":"xadrez","fen_final":"r1bqk2r/p1pn2p1/1p2pnBp/8/3P4/B1P5/P1P2PPP/R3K1NR b - - 0 1","fonte":"partidas reais","jogadores":"Alexander Alekhine × Vasic","lances_san":["e4","e6","d4","d5","Nc3","Bb4","Bd3","Bxc3+","bxc3","h6","Ba3","Nd7","Qe2","dxe4","Bxe4","Ngf6","Bd3","b6","Qxe6+","fxe6","Bg6#"],"lyapunov":0.4758,"mineral_H":0.2878,"mineral_classe":"borda","mineral_theta_g":85.7,"n_lances":21,"orbita_fen":["rnbqkbnr/pppppppp/8/8/8/8/PPPPPPPP/RNBQKBNR w - - 0 1","rnbqkbnr/pppppppp/8/8/4P3/8/PPPP1PPP/RNBQKBNR b - e3 0 1","rnbqkbnr/pppp1ppp/4p3/8/4P3/8/PPPP1PPP/RNBQKBNR w - - 0 1","rnbqkbnr/pppp1ppp/4p3/8/3PP3/8/PPP2PPP/RNBQKBNR b - d3 0 1","rnbqkbnr/ppp2ppp/4p3/3p4/3PP3/8/PPP2PPP/RNBQKBNR w - d6 0 1","rnbqkbnr/ppp2ppp/4p3/3p4/3PP3/2N5/PPP2PPP/R1BQKBNR b - - 0 1","rnbqk1nr/ppp2ppp/4p3/3p4/1b1PP3/2N5/PPP2PPP/R1BQKBNR w - - 0 1","rnbqk1nr/ppp2ppp/4p3/3p4/1b1PP3/2NB4/PPP2PPP/R1BQK1NR b - - 0 1","rnbqk1nr/ppp2ppp/4p3/3p4/3PP3/2bB4/PPP2PPP/R1BQK1NR w - - 0 1","rnbqk1nr/ppp2ppp/4p3/3p4/3PP3/2PB4/P1P2PPP/R1BQK1NR b - - 0 1","rnbqk1nr/ppp2pp1/4p2p/3p4/3PP3/2PB4/P1P2PPP/R1BQK1NR w - - 0 1","rnbqk1nr/ppp2pp1/4p2p/3p4/3PP3/B1PB4/P1P2PPP/R2QK1NR b - - 0 1","r1bqk1nr/pppn1pp1/4p2p/3p4/3PP3/B1PB4/P1P2PPP/R2QK1NR w - - 0 1","r1bqk1nr/pppn1pp1/4p2p/3p4/3PP3/B1PB4/P1P1QPPP/R3K1NR b - - 0 1","r1bqk1nr/pppn1pp1/4p2p/8/3Pp3/B1PB4/P1P1QPPP/R3K1NR w - - 0 1","r1bqk1nr/pppn1pp1/4p2p/8/3PB3/B1P5/P1P1QPPP/R3K1NR b - - 0 1","r1bqk2r/pppn1pp1/4pn1p/8/3PB3/B1P5/P1P1QPPP/R3K1NR w - - 0 1","r1bqk2r/pppn1pp1/4pn1p/8/3P4/B1PB4/P1P1QPPP/R3K1NR b - - 0 1","r1bqk2r/p1pn1pp1/1p2pn1p/8/3P4/B1PB4/P1P1QPPP/R3K1NR w - - 0 1","r1bqk2r/p1pn1pp1/1p2Qn1p/8/3P4/B1PB4/P1P2PPP/R3K1NR b - - 0 1","r1bqk2r/p1pn2p1/1p2pn1p/8/3P4/B1PB4/P1P2PPP/R3K1NR w - - 0 1","r1bqk2r/p1pn2p1/1p2pnBp/8/3P4/B1P5/P1P2PPP/R3K1NR b - - 0 1"],"resultado":"xeque-mate","tabuleiro_final":"8 r . b q k . . r\n7 p . p n . . p .\n6 . p . . p n B p\n5 . . . . . . . .\n4 . . . P . . . .\n3 B . P . . . . .\n2 P . P . . P P P\n1 R . . . K . N R\n  a b c d e f g h"},"origem":"wiki","palavras_chave":[],"sinonimos":[],"tipo":"conceito","titulo":"Alekhine × Vasic (1931)"}
\section{Alekhine $\times$ Vasic (1931)}
Sacrifício de dama (Dxe6+!!) que rasga o roque e o bispo dá mate em Bg6\# — uma miniatura de 11 lances do campeão. Jogadores: Alexander Alekhine $\times$ Vasic (1931). A órbita percorre 22 posições e termina em xeque-mate após 21 lances.
\prov{relações: eh\_um orbita\_de\_partida; usa abertura\_xadrez; relaciona xeque\_mate; usa sacrificio\_xadrez; relaciona alekhine; usa lyapunov\_da\_orbita}
\prov{proveniência: wiki \textperiodcentered{} partidas reais \textperiodcentered{} fato \textperiodcentered{} confiança 1.0 \textperiodcentered{} domínio xadrez \textperiodcentered{} história 14 versões no jornal}

%CRISTAL {"arestas":[{"alvo":"orbita_de_partida","confianca":1.0,"epistemico":"fato","origem":"wiki","rel":"eh_um"},{"alvo":"abertura_xadrez","confianca":1.0,"epistemico":"fato","origem":"wiki","rel":"usa"},{"alvo":"xeque_mate","confianca":1.0,"epistemico":"fato","origem":"wiki","rel":"relaciona"},{"alvo":"tatica_xadrez","confianca":1.0,"epistemico":"fato","origem":"wiki","rel":"usa"},{"alvo":"lyapunov_da_orbita","confianca":1.0,"epistemico":"fato","origem":"wiki","rel":"usa"}],"confianca":1.0,"contraexemplos":[],"descricao":"O xeque-mate mais rápido possível — dois lances das brancas abrem a diagonal para Dh4#. A prova de que a segurança do rei vem antes de tudo. Jogadores: (padrão didático). A órbita percorre 5 posições e termina em xeque-mate após 4 lances.","epistemico":"fato","exemplos":[],"id":"partida_bobo","memoria":"semantica","meta":{"abertura":"Abertura Bird invertida","ano":null,"classe_orbita":"cristal","dominio":"xadrez","fen_final":"rnb1kbnr/pppp1ppp/8/4p3/6Pq/5P2/PPPPP2P/RNBQKBNR w - - 0 1","fonte":"partidas reais","jogadores":"(padrão didático)","lances_san":["f3","e5","g4","Qh4#"],"lyapunov":0.1535,"n_lances":4,"orbita_fen":["rnbqkbnr/pppppppp/8/8/8/8/PPPPPPPP/RNBQKBNR w - - 0 1","rnbqkbnr/pppppppp/8/8/8/5P2/PPPPP1PP/RNBQKBNR b - - 0 1","rnbqkbnr/pppp1ppp/8/4p3/8/5P2/PPPPP1PP/RNBQKBNR w - e6 0 1","rnbqkbnr/pppp1ppp/8/4p3/6P1/5P2/PPPPP2P/RNBQKBNR b - g3 0 1","rnb1kbnr/pppp1ppp/8/4p3/6Pq/5P2/PPPPP2P/RNBQKBNR w - - 0 1"],"resultado":"xeque-mate","tabuleiro_final":"8 r n b . k b n r\n7 p p p p . p p p\n6 . . . . . . . .\n5 . . . . p . . .\n4 . . . . . . P q\n3 . . . . . P . .\n2 P P P P P . . P\n1 R N B Q K B N R\n  a b c d e f g h"},"origem":"wiki","palavras_chave":[],"sinonimos":[],"tipo":"conceito","titulo":"Mate do Bobo (o mais rápido)"}
\section{Mate do Bobo (o mais rápido)}
O xeque-mate mais rápido possível — dois lances das brancas abrem a diagonal para Dh4\#. A prova de que a segurança do rei vem antes de tudo. Jogadores: (padrão didático). A órbita percorre 5 posições e termina em xeque-mate após 4 lances.
\prov{relações: eh\_um orbita\_de\_partida; usa abertura\_xadrez; relaciona xeque\_mate; usa tatica\_xadrez; usa lyapunov\_da\_orbita}
\prov{proveniência: wiki \textperiodcentered{} partidas reais \textperiodcentered{} fato \textperiodcentered{} confiança 1.0 \textperiodcentered{} domínio xadrez \textperiodcentered{} história 48 versões no jornal}

%CRISTAL {"arestas":[{"alvo":"paul_morphy","confianca":1.0,"epistemico":"fato","origem":"wiki","rel":"parte_de"},{"alvo":"principios_de_abertura","confianca":1.0,"epistemico":"fato","origem":"wiki","rel":"explica"}],"confianca":1.0,"contraexemplos":[],"descricao":"Morphy, num camarote em Paris, dá uma aula de desenvolvimento e iniciativa contra dois nobres, terminando com sacrifício de dama e mate. A partida didática mais famosa da história.","epistemico":"fato","exemplos":[],"id":"partida_da_opera","memoria":"semantica","meta":{"autor":"Paul Morphy","dominio":"xadrez","fonte":"história do xadrez"},"origem":"wiki","palavras_chave":[],"sinonimos":[],"tipo":"conceito","titulo":"A Partida da Ópera (Morphy, 1858)"}
\section{A Partida da Ópera (Morphy, 1858)}
Morphy, num camarote em Paris, dá uma aula de desenvolvimento e iniciativa contra dois nobres, terminando com sacrifício de dama e mate. A partida didática mais famosa da história.
\prov{relações: parte\_de paul\_morphy; explica principios\_de\_abertura}
\prov{proveniência: wiki \textperiodcentered{} história do xadrez \textperiodcentered{} fato \textperiodcentered{} confiança 1.0 \textperiodcentered{} domínio xadrez \textperiodcentered{} história 3 versões no jornal}

%CRISTAL {"arestas":[{"alvo":"orbita_de_partida","confianca":1.0,"epistemico":"fato","origem":"wiki","rel":"eh_um"},{"alvo":"abertura_xadrez","confianca":1.0,"epistemico":"fato","origem":"wiki","rel":"usa"},{"alvo":"xeque_mate","confianca":1.0,"epistemico":"fato","origem":"wiki","rel":"relaciona"},{"alvo":"sacrificio_xadrez","confianca":1.0,"epistemico":"fato","origem":"wiki","rel":"usa"},{"alvo":"deep_blue","confianca":1.0,"epistemico":"fato","origem":"wiki","rel":"relaciona"},{"alvo":"kasparov","confianca":1.0,"epistemico":"fato","origem":"wiki","rel":"relaciona"},{"alvo":"lyapunov_da_orbita","confianca":1.0,"epistemico":"fato","origem":"wiki","rel":"usa"}],"confianca":1.0,"contraexemplos":[],"descricao":"O jogo em que uma MÁQUINA venceu o campeão mundial num match: Deep Blue sacrifica o cavalo (Nxe6!) contra a Caro-Kann e Kasparov abandona no 19º lance. O marco da IA no xadrez (Kasparov abandonou — a órbita para 'em jogo', com vantagem decisiva das brancas). Jogadores: Deep Blue × Garry Kasparov (1997). A órbita percorre 38 posições e termina em em jogo após 37 lances.","epistemico":"fato","exemplos":[],"id":"partida_deepblue","memoria":"semantica","meta":{"abertura":"Defesa Caro-Kann","ano":1997,"classe_orbita":"borda","dominio":"xadrez","fen_final":"r1k4r/p2nb1p1/2b4p/1p1n1p2/2PP4/3Q1NB1/1P3PPP/R5K1 b - c3 0 1","fonte":"partidas reais","jogadores":"Deep Blue × Garry Kasparov","lances_san":["e4","c6","d4","d5","Nc3","dxe4","Nxe4","Nd7","Ng5","Ngf6","Bd3","e6","N1f3","h6","Nxe6","Qe7","O-O","fxe6","Bg6+","Kd8","Bf4","b5","a4","Bb7","Re1","Nd5","Bg3","Kc8","axb5","cxb5","Qd3","Bc6","Bf5","exf5","Rxe7","Bxe7","c4"],"lyapunov":0.3872,"mineral_H":-0.1401,"mineral_classe":"cristal","mineral_theta_g":126.1,"n_lances":37,"orbita_fen":["rnbqkbnr/pppppppp/8/8/8/8/PPPPPPPP/RNBQKBNR w - - 0 1","rnbqkbnr/pppppppp/8/8/4P3/8/PPPP1PPP/RNBQKBNR b - e3 0 1","rnbqkbnr/pp1ppppp/2p5/8/4P3/8/PPPP1PPP/RNBQKBNR w - - 0 1","rnbqkbnr/pp1ppppp/2p5/8/3PP3/8/PPP2PPP/RNBQKBNR b - d3 0 1","rnbqkbnr/pp2pppp/2p5/3p4/3PP3/8/PPP2PPP/RNBQKBNR w - d6 0 1","rnbqkbnr/pp2pppp/2p5/3p4/3PP3/2N5/PPP2PPP/R1BQKBNR b - - 0 1","rnbqkbnr/pp2pppp/2p5/8/3Pp3/2N5/PPP2PPP/R1BQKBNR w - - 0 1","rnbqkbnr/pp2pppp/2p5/8/3PN3/8/PPP2PPP/R1BQKBNR b - - 0 1","r1bqkbnr/pp1npppp/2p5/8/3PN3/8/PPP2PPP/R1BQKBNR w - - 0 1","r1bqkbnr/pp1npppp/2p5/6N1/3P4/8/PPP2PPP/R1BQKBNR b - - 0 1","r1bqkb1r/pp1npppp/2p2n2/6N1/3P4/8/PPP2PPP/R1BQKBNR w - - 0 1","r1bqkb1r/pp1npppp/2p2n2/6N1/3P4/3B4/PPP2PPP/R1BQK1NR b - - 0 1","r1bqkb1r/pp1n1ppp/2p1pn2/6N1/3P4/3B4/PPP2PPP/R1BQK1NR w - - 0 1","r1bqkb1r/pp1n1ppp/2p1pn2/6N1/3P4/3B1N2/PPP2PPP/R1BQK2R b - - 0 1","r1bqkb1r/pp1n1pp1/2p1pn1p/6N1/3P4/3B1N2/PPP2PPP/R1BQK2R w - - 0 1","r1bqkb1r/pp1n1pp1/2p1Nn1p/8/3P4/3B1N2/PPP2PPP/R1BQK2R b - - 0 1","r1b1kb1r/pp1nqpp1/2p1Nn1p/8/3P4/3B1N2/PPP2PPP/R1BQK2R w - - 0 1","r1b1kb1r/pp1nqpp1/2p1Nn1p/8/3P4/3B1N2/PPP2PPP/R1BQ1RK1 b - - 0 1","r1b1kb1r/pp1nq1p1/2p1pn1p/8/3P4/3B1N2/PPP2PPP/R1BQ1RK1 w - - 0 1","r1b1kb1r/pp1nq1p1/2p1pnBp/8/3P4/5N2/PPP2PPP/R1BQ1RK1 b - - 0 1","r1bk1b1r/pp1nq1p1/2p1pnBp/8/3P4/5N2/PPP2PPP/R1BQ1RK1 w - - 0 1","r1bk1b1r/pp1nq1p1/2p1pnBp/8/3P1B2/5N2/PPP2PPP/R2Q1RK1 b - - 0 1","r1bk1b1r/p2nq1p1/2p1pnBp/1p6/3P1B2/5N2/PPP2PPP/R2Q1RK1 w - b6 0 1","r1bk1b1r/p2nq1p1/2p1pnBp/1p6/P2P1B2/5N2/1PP2PPP/R2Q1RK1 b - a3 0 1","r2k1b1r/pb1nq1p1/2p1pnBp/1p6/P2P1B2/5N2/1PP2PPP/R2Q1RK1 w - - 0 1","r2k1b1r/pb1nq1p1/2p1pnBp/1p6/P2P1B2/5N2/1PP2PPP/R2QR1K1 b - - 0 1","r2k1b1r/pb1nq1p1/2p1p1Bp/1p1n4/P2P1B2/5N2/1PP2PPP/R2QR1K1 w - - 0 1","r2k1b1r/pb1nq1p1/2p1p1Bp/1p1n4/P2P4/5NB1/1PP2PPP/R2QR1K1 b - - 0 1","r1k2b1r/pb1nq1p1/2p1p1Bp/1p1n4/P2P4/5NB1/1PP2PPP/R2QR1K1 w - - 0 1","r1k2b1r/pb1nq1p1/2p1p1Bp/1P1n4/3P4/5NB1/1PP2PPP/R2QR1K1 b - - 0 1","r1k2b1r/pb1nq1p1/4p1Bp/1p1n4/3P4/5NB1/1PP2PPP/R2QR1K1 w - - 0 1","r1k2b1r/pb1nq1p1/4p1Bp/1p1n4/3P4/3Q1NB1/1PP2PPP/R3R1K1 b - - 0 1","r1k2b1r/p2nq1p1/2b1p1Bp/1p1n4/3P4/3Q1NB1/1PP2PPP/R3R1K1 w - - 0 1","r1k2b1r/p2nq1p1/2b1p2p/1p1n1B2/3P4/3Q1NB1/1PP2PPP/R3R1K1 b - - 0 1","r1k2b1r/p2nq1p1/2b4p/1p1n1p2/3P4/3Q1NB1/1PP2PPP/R3R1K1 w - - 0 1","r1k2b1r/p2nR1p1/2b4p/1p1n1p2/3P4/3Q1NB1/1PP2PPP/R5K1 b - - 0 1","r1k4r/p2nb1p1/2b4p/1p1n1p2/3P4/3Q1NB1/1PP2PPP/R5K1 w - - 0 1","r1k4r/p2nb1p1/2b4p/1p1n1p2/2PP4/3Q1NB1/1P3PPP/R5K1 b - c3 0 1"],"resultado":"em jogo","tabuleiro_final":"8 r . k . . . . r\n7 p . . n b . p .\n6 . . b . . . . p\n5 . p . n . p . .\n4 . . P P . . . .\n3 . . . Q . N B .\n2 . P . . . P P P\n1 R . . . . . K .\n  a b c d e f g h"},"origem":"wiki","palavras_chave":[],"sinonimos":[],"tipo":"conceito","titulo":"Deep Blue × Kasparov (1997, Jogo 6)"}
\section{Deep Blue $\times$ Kasparov (1997, Jogo 6)}
O jogo em que uma MÁQUINA venceu o campeão mundial num match: Deep Blue sacrifica o cavalo (Nxe6!) contra a Caro-Kann e Kasparov abandona no 19º lance. O marco da IA no xadrez (Kasparov abandonou — a órbita para 'em jogo', com vantagem decisiva das brancas). Jogadores: Deep Blue $\times$ Garry Kasparov (1997). A órbita percorre 38 posições e termina em em jogo após 37 lances.
\prov{relações: eh\_um orbita\_de\_partida; usa abertura\_xadrez; relaciona xeque\_mate; usa sacrificio\_xadrez; relaciona deep\_blue; relaciona kasparov; usa lyapunov\_da\_orbita}
\prov{proveniência: wiki \textperiodcentered{} partidas reais \textperiodcentered{} fato \textperiodcentered{} confiança 1.0 \textperiodcentered{} domínio xadrez \textperiodcentered{} história 47 versões no jornal}

%CRISTAL {"arestas":[{"alvo":"bobby_fischer","confianca":1.0,"epistemico":"fato","origem":"wiki","rel":"parte_de"},{"alvo":"sacrificio_xadrez","confianca":1.0,"epistemico":"fato","origem":"wiki","rel":"usa"}],"confianca":1.0,"contraexemplos":[],"descricao":"Fischer, aos 13 anos, sacrifica a dama e conduz uma combinação imortal com peças menores. O anúncio de um gênio; a demonstração de que material não é tudo.","epistemico":"fato","exemplos":[],"id":"partida_do_seculo","memoria":"semantica","meta":{"autor":"Fischer","dominio":"xadrez","fonte":"história do xadrez"},"origem":"wiki","palavras_chave":[],"sinonimos":[],"tipo":"conceito","titulo":"A Partida do Século (D. Byrne–Fischer, 1956)"}
\section{A Partida do Século (D. Byrne–Fischer, 1956)}
Fischer, aos 13 anos, sacrifica a dama e conduz uma combinação imortal com peças menores. O anúncio de um gênio; a demonstração de que material não é tudo.
\prov{relações: parte\_de bobby\_fischer; usa sacrificio\_xadrez}
\prov{proveniência: wiki \textperiodcentered{} história do xadrez \textperiodcentered{} fato \textperiodcentered{} confiança 1.0 \textperiodcentered{} domínio xadrez \textperiodcentered{} história 3 versões no jornal}

%CRISTAL {"arestas":[{"alvo":"orbita_de_partida","confianca":1.0,"epistemico":"fato","origem":"wiki","rel":"eh_um"},{"alvo":"abertura_xadrez","confianca":1.0,"epistemico":"fato","origem":"wiki","rel":"usa"},{"alvo":"xeque_mate","confianca":1.0,"epistemico":"fato","origem":"wiki","rel":"relaciona"},{"alvo":"sacrificio_xadrez","confianca":1.0,"epistemico":"fato","origem":"wiki","rel":"usa"},{"alvo":"anderssen","confianca":1.0,"epistemico":"fato","origem":"wiki","rel":"relaciona"},{"alvo":"lyapunov_da_orbita","confianca":1.0,"epistemico":"fato","origem":"wiki","rel":"usa"}],"confianca":1.0,"contraexemplos":[],"descricao":"A 'Evergreen' — a irmã da Imortal. Anderssen sacrifica a dama (Dxd7+!!) e conclui com o duplo sacrifício de bispo num mate de peças menores. A joia do romantismo. Jogadores: Adolf Anderssen × Jean Dufresne (1852). A órbita percorre 48 posições e termina em xeque-mate após 47 lances.","epistemico":"fato","exemplos":[],"id":"partida_evergreen","memoria":"semantica","meta":{"abertura":"Gambito Evans","ano":1852,"classe_orbita":"borda","dominio":"xadrez","fen_final":"1r3kr1/pbpBBp1p/1b3P2/8/8/2P2q2/P4PPP/3R2K1 b - - 0 1","fonte":"partidas reais","jogadores":"Adolf Anderssen × Jean Dufresne","lances_san":["e4","e5","Nf3","Nc6","Bc4","Bc5","b4","Bxb4","c3","Ba5","d4","exd4","O-O","d3","Qb3","Qf6","e5","Qg6","Re1","Nge7","Ba3","b5","Qxb5","Rb8","Qa4","Bb6","Nbd2","Bb7","Ne4","Qf5","Bxd3","Qh5","Nf6+","gxf6","exf6","Rg8","Rad1","Qxf3","Rxe7+","Nxe7","Qxd7+","Kxd7","Bf5+","Ke8","Bd7+","Kf8","Bxe7#"],"lyapunov":0.3692,"mineral_H":-0.0328,"mineral_classe":"cristal","mineral_theta_g":180.0,"n_lances":47,"orbita_fen":["rnbqkbnr/pppppppp/8/8/8/8/PPPPPPPP/RNBQKBNR w - - 0 1","rnbqkbnr/pppppppp/8/8/4P3/8/PPPP1PPP/RNBQKBNR b - e3 0 1","rnbqkbnr/pppp1ppp/8/4p3/4P3/8/PPPP1PPP/RNBQKBNR w - e6 0 1","rnbqkbnr/pppp1ppp/8/4p3/4P3/5N2/PPPP1PPP/RNBQKB1R b - - 0 1","r1bqkbnr/pppp1ppp/2n5/4p3/4P3/5N2/PPPP1PPP/RNBQKB1R w - - 0 1","r1bqkbnr/pppp1ppp/2n5/4p3/2B1P3/5N2/PPPP1PPP/RNBQK2R b - - 0 1","r1bqk1nr/pppp1ppp/2n5/2b1p3/2B1P3/5N2/PPPP1PPP/RNBQK2R w - - 0 1","r1bqk1nr/pppp1ppp/2n5/2b1p3/1PB1P3/5N2/P1PP1PPP/RNBQK2R b - b3 0 1","r1bqk1nr/pppp1ppp/2n5/4p3/1bB1P3/5N2/P1PP1PPP/RNBQK2R w - - 0 1","r1bqk1nr/pppp1ppp/2n5/4p3/1bB1P3/2P2N2/P2P1PPP/RNBQK2R b - - 0 1","r1bqk1nr/pppp1ppp/2n5/b3p3/2B1P3/2P2N2/P2P1PPP/RNBQK2R w - - 0 1","r1bqk1nr/pppp1ppp/2n5/b3p3/2BPP3/2P2N2/P4PPP/RNBQK2R b - d3 0 1","r1bqk1nr/pppp1ppp/2n5/b7/2BpP3/2P2N2/P4PPP/RNBQK2R w - - 0 1","r1bqk1nr/pppp1ppp/2n5/b7/2BpP3/2P2N2/P4PPP/RNBQ1RK1 b - - 0 1","r1bqk1nr/pppp1ppp/2n5/b7/2B1P3/2Pp1N2/P4PPP/RNBQ1RK1 w - - 0 1","r1bqk1nr/pppp1ppp/2n5/b7/2B1P3/1QPp1N2/P4PPP/RNB2RK1 b - - 0 1","r1b1k1nr/pppp1ppp/2n2q2/b7/2B1P3/1QPp1N2/P4PPP/RNB2RK1 w - - 0 1","r1b1k1nr/pppp1ppp/2n2q2/b3P3/2B5/1QPp1N2/P4PPP/RNB2RK1 b - - 0 1","r1b1k1nr/pppp1ppp/2n3q1/b3P3/2B5/1QPp1N2/P4PPP/RNB2RK1 w - - 0 1","r1b1k1nr/pppp1ppp/2n3q1/b3P3/2B5/1QPp1N2/P4PPP/RNB1R1K1 b - - 0 1","r1b1k2r/ppppnppp/2n3q1/b3P3/2B5/1QPp1N2/P4PPP/RNB1R1K1 w - - 0 1","r1b1k2r/ppppnppp/2n3q1/b3P3/2B5/BQPp1N2/P4PPP/RN2R1K1 b - - 0 1","r1b1k2r/p1ppnppp/2n3q1/bp2P3/2B5/BQPp1N2/P4PPP/RN2R1K1 w - b6 0 1","r1b1k2r/p1ppnppp/2n3q1/bQ2P3/2B5/B1Pp1N2/P4PPP/RN2R1K1 b - - 0 1","1rb1k2r/p1ppnppp/2n3q1/bQ2P3/2B5/B1Pp1N2/P4PPP/RN2R1K1 w - - 0 1","1rb1k2r/p1ppnppp/2n3q1/b3P3/Q1B5/B1Pp1N2/P4PPP/RN2R1K1 b - - 0 1","1rb1k2r/p1ppnppp/1bn3q1/4P3/Q1B5/B1Pp1N2/P4PPP/RN2R1K1 w - - 0 1","1rb1k2r/p1ppnppp/1bn3q1/4P3/Q1B5/B1Pp1N2/P2N1PPP/R3R1K1 b - - 0 1","1r2k2r/pbppnppp/1bn3q1/4P3/Q1B5/B1Pp1N2/P2N1PPP/R3R1K1 w - - 0 1","1r2k2r/pbppnppp/1bn3q1/4P3/Q1B1N3/B1Pp1N2/P4PPP/R3R1K1 b - - 0 1","1r2k2r/pbppnppp/1bn5/4Pq2/Q1B1N3/B1Pp1N2/P4PPP/R3R1K1 w - - 0 1","1r2k2r/pbppnppp/1bn5/4Pq2/Q3N3/B1PB1N2/P4PPP/R3R1K1 b - - 0 1","1r2k2r/pbppnppp/1bn5/4P2q/Q3N3/B1PB1N2/P4PPP/R3R1K1 w - - 0 1","1r2k2r/pbppnppp/1bn2N2/4P2q/Q7/B1PB1N2/P4PPP/R3R1K1 b - - 0 1","1r2k2r/pbppnp1p/1bn2p2/4P2q/Q7/B1PB1N2/P4PPP/R3R1K1 w - - 0 1","1r2k2r/pbppnp1p/1bn2P2/7q/Q7/B1PB1N2/P4PPP/R3R1K1 b - - 0 1","1r2k1r1/pbppnp1p/1bn2P2/7q/Q7/B1PB1N2/P4PPP/R3R1K1 w - - 0 1","1r2k1r1/pbppnp1p/1bn2P2/7q/Q7/B1PB1N2/P4PPP/3RR1K1 b - - 0 1","1r2k1r1/pbppnp1p/1bn2P2/8/Q7/B1PB1q2/P4PPP/3RR1K1 w - - 0 1","1r2k1r1/pbppRp1p/1bn2P2/8/Q7/B1PB1q2/P4PPP/3R2K1 b - - 0 1","1r2k1r1/pbppnp1p/1b3P2/8/Q7/B1PB1q2/P4PPP/3R2K1 w - - 0 1","1r2k1r1/pbpQnp1p/1b3P2/8/8/B1PB1q2/P4PPP/3R2K1 b - - 0 1","1r4r1/pbpknp1p/1b3P2/8/8/B1PB1q2/P4PPP/3R2K1 w - - 0 1","1r4r1/pbpknp1p/1b3P2/5B2/8/B1P2q2/P4PPP/3R2K1 b - - 0 1","1r2k1r1/pbp1np1p/1b3P2/5B2/8/B1P2q2/P4PPP/3R2K1 w - - 0 1","1r2k1r1/pbpBnp1p/1b3P2/8/8/B1P2q2/P4PPP/3R2K1 b - - 0 1","1r3kr1/pbpBnp1p/1b3P2/8/8/B1P2q2/P4PPP/3R2K1 w - - 0 1","1r3kr1/pbpBBp1p/1b3P2/8/8/2P2q2/P4PPP/3R2K1 b - - 0 1"],"resultado":"xeque-mate","tabuleiro_final":"8 . r . . . k r .\n7 p b p B B p . p\n6 . b . . . P . .\n5 . . . . . . . .\n4 . . . . . . . .\n3 . . P . . q . .\n2 P . . . . P P P\n1 . . . R . . K .\n  a b c d e f g h"},"origem":"wiki","palavras_chave":[],"sinonimos":[],"tipo":"conceito","titulo":"A Sempreviva (Anderssen, 1852)"}
\section{A Sempreviva (Anderssen, 1852)}
A 'Evergreen' — a irmã da Imortal. Anderssen sacrifica a dama (Dxd7+!!) e conclui com o duplo sacrifício de bispo num mate de peças menores. A joia do romantismo. Jogadores: Adolf Anderssen $\times$ Jean Dufresne (1852). A órbita percorre 48 posições e termina em xeque-mate após 47 lances.
\prov{relações: eh\_um orbita\_de\_partida; usa abertura\_xadrez; relaciona xeque\_mate; usa sacrificio\_xadrez; relaciona anderssen; usa lyapunov\_da\_orbita}
\prov{proveniência: wiki \textperiodcentered{} partidas reais \textperiodcentered{} fato \textperiodcentered{} confiança 1.0 \textperiodcentered{} domínio xadrez \textperiodcentered{} história 41 versões no jornal}

%CRISTAL {"arestas":[{"alvo":"orbita_de_partida","confianca":1.0,"epistemico":"fato","origem":"wiki","rel":"eh_um"},{"alvo":"abertura_xadrez","confianca":1.0,"epistemico":"fato","origem":"wiki","rel":"usa"},{"alvo":"xeque_mate","confianca":1.0,"epistemico":"fato","origem":"wiki","rel":"relaciona"},{"alvo":"sacrificio_xadrez","confianca":1.0,"epistemico":"fato","origem":"wiki","rel":"usa"},{"alvo":"byrne","confianca":1.0,"epistemico":"fato","origem":"wiki","rel":"relaciona"},{"alvo":"fischer","confianca":1.0,"epistemico":"fato","origem":"wiki","rel":"relaciona"},{"alvo":"lyapunov_da_orbita","confianca":1.0,"epistemico":"fato","origem":"wiki","rel":"usa"}],"confianca":1.0,"contraexemplos":[],"descricao":"Fischer, aos 13 anos, sacrifica a dama (Be6!!) e conduz um moinho de peças menores até Tc2#. A partida que anunciou o gênio — 'a Partida do Século'. Jogadores: Donald Byrne × Bobby Fischer (1956). A órbita percorre 83 posições e termina em xeque-mate após 82 lances.","epistemico":"fato","exemplos":[],"id":"partida_fischer_byrne","memoria":"semantica","meta":{"abertura":"Defesa Grünfeld","ano":1956,"classe_orbita":"borda","dominio":"xadrez","fen_final":"1Q6/5pk1/2p3p1/1p2N2p/1b5P/1bn5/2r3P1/2K5 w - - 0 1","fonte":"partidas reais","jogadores":"Donald Byrne × Bobby Fischer","lances_san":["Nf3","Nf6","c4","g6","Nc3","Bg7","d4","O-O","Bf4","d5","Qb3","dxc4","Qxc4","c6","e4","Nbd7","Rd1","Nb6","Qc5","Bg4","Bg5","Na4","Qa3","Nxc3","bxc3","Nxe4","Bxe7","Qb6","Bc4","Nxc3","Bc5","Rfe8+","Kf1","Be6","Bxb6","Bxc4+","Kg1","Ne2+","Kf1","Nxd4+","Kg1","Ne2+","Kf1","Nc3+","Kg1","axb6","Qb4","Ra4","Qxb6","Nxd1","h3","Rxa2","Kh2","Nxf2","Re1","Rxe1","Qd8+","Bf8","Nxe1","Bd5","Nf3","Ne4","Qb8","b5","h4","h5","Ne5","Kg7","Kg1","Bc5+","Kf1","Ng3+","Ke1","Bb4+","Kd1","Bb3+","Kc1","Ne2+","Kb1","Nc3+","Kc1","Rc2#"],"lyapunov":0.3011,"mineral_H":-0.0323,"mineral_classe":"cristal","mineral_theta_g":180.0,"n_lances":82,"orbita_fen":["rnbqkbnr/pppppppp/8/8/8/8/PPPPPPPP/RNBQKBNR w - - 0 1","rnbqkbnr/pppppppp/8/8/8/5N2/PPPPPPPP/RNBQKB1R b - - 0 1","rnbqkb1r/pppppppp/5n2/8/8/5N2/PPPPPPPP/RNBQKB1R w - - 0 1","rnbqkb1r/pppppppp/5n2/8/2P5/5N2/PP1PPPPP/RNBQKB1R b - c3 0 1","rnbqkb1r/pppppp1p/5np1/8/2P5/5N2/PP1PPPPP/RNBQKB1R w - - 0 1","rnbqkb1r/pppppp1p/5np1/8/2P5/2N2N2/PP1PPPPP/R1BQKB1R b - - 0 1","rnbqk2r/ppppppbp/5np1/8/2P5/2N2N2/PP1PPPPP/R1BQKB1R w - - 0 1","rnbqk2r/ppppppbp/5np1/8/2PP4/2N2N2/PP2PPPP/R1BQKB1R b - d3 0 1","rnbq1rk1/ppppppbp/5np1/8/2PP4/2N2N2/PP2PPPP/R1BQKB1R w - - 0 1","rnbq1rk1/ppppppbp/5np1/8/2PP1B2/2N2N2/PP2PPPP/R2QKB1R b - - 0 1","rnbq1rk1/ppp1ppbp/5np1/3p4/2PP1B2/2N2N2/PP2PPPP/R2QKB1R w - d6 0 1","rnbq1rk1/ppp1ppbp/5np1/3p4/2PP1B2/1QN2N2/PP2PPPP/R3KB1R b - - 0 1","rnbq1rk1/ppp1ppbp/5np1/8/2pP1B2/1QN2N2/PP2PPPP/R3KB1R w - - 0 1","rnbq1rk1/ppp1ppbp/5np1/8/2QP1B2/2N2N2/PP2PPPP/R3KB1R b - - 0 1","rnbq1rk1/pp2ppbp/2p2np1/8/2QP1B2/2N2N2/PP2PPPP/R3KB1R w - - 0 1","rnbq1rk1/pp2ppbp/2p2np1/8/2QPPB2/2N2N2/PP3PPP/R3KB1R b - e3 0 1","r1bq1rk1/pp1nppbp/2p2np1/8/2QPPB2/2N2N2/PP3PPP/R3KB1R w - - 0 1","r1bq1rk1/pp1nppbp/2p2np1/8/2QPPB2/2N2N2/PP3PPP/3RKB1R b - - 0 1","r1bq1rk1/pp2ppbp/1np2np1/8/2QPPB2/2N2N2/PP3PPP/3RKB1R w - - 0 1","r1bq1rk1/pp2ppbp/1np2np1/2Q5/3PPB2/2N2N2/PP3PPP/3RKB1R b - - 0 1","r2q1rk1/pp2ppbp/1np2np1/2Q5/3PPBb1/2N2N2/PP3PPP/3RKB1R w - - 0 1","r2q1rk1/pp2ppbp/1np2np1/2Q3B1/3PP1b1/2N2N2/PP3PPP/3RKB1R b - - 0 1","r2q1rk1/pp2ppbp/2p2np1/2Q3B1/n2PP1b1/2N2N2/PP3PPP/3RKB1R w - - 0 1","r2q1rk1/pp2ppbp/2p2np1/6B1/n2PP1b1/Q1N2N2/PP3PPP/3RKB1R b - - 0 1","r2q1rk1/pp2ppbp/2p2np1/6B1/3PP1b1/Q1n2N2/PP3PPP/3RKB1R w - - 0 1","r2q1rk1/pp2ppbp/2p2np1/6B1/3PP1b1/Q1P2N2/P4PPP/3RKB1R b - - 0 1","r2q1rk1/pp2ppbp/2p3p1/6B1/3Pn1b1/Q1P2N2/P4PPP/3RKB1R w - - 0 1","r2q1rk1/pp2Bpbp/2p3p1/8/3Pn1b1/Q1P2N2/P4PPP/3RKB1R b - - 0 1","r4rk1/pp2Bpbp/1qp3p1/8/3Pn1b1/Q1P2N2/P4PPP/3RKB1R w - - 0 1","r4rk1/pp2Bpbp/1qp3p1/8/2BPn1b1/Q1P2N2/P4PPP/3RK2R b - - 0 1","r4rk1/pp2Bpbp/1qp3p1/8/2BP2b1/Q1n2N2/P4PPP/3RK2R w - - 0 1","r4rk1/pp3pbp/1qp3p1/2B5/2BP2b1/Q1n2N2/P4PPP/3RK2R b - - 0 1","r3r1k1/pp3pbp/1qp3p1/2B5/2BP2b1/Q1n2N2/P4PPP/3RK2R w - - 0 1","r3r1k1/pp3pbp/1qp3p1/2B5/2BP2b1/Q1n2N2/P4PPP/3R1K1R b - - 0 1","r3r1k1/pp3pbp/1qp1b1p1/2B5/2BP4/Q1n2N2/P4PPP/3R1K1R w - - 0 1","r3r1k1/pp3pbp/1Bp1b1p1/8/2BP4/Q1n2N2/P4PPP/3R1K1R b - - 0 1","r3r1k1/pp3pbp/1Bp3p1/8/2bP4/Q1n2N2/P4PPP/3R1K1R w - - 0 1","r3r1k1/pp3pbp/1Bp3p1/8/2bP4/Q1n2N2/P4PPP/3R2KR b - - 0 1","r3r1k1/pp3pbp/1Bp3p1/8/2bP4/Q4N2/P3nPPP/3R2KR w - - 0 1","r3r1k1/pp3pbp/1Bp3p1/8/2bP4/Q4N2/P3nPPP/3R1K1R b - - 0 1","r3r1k1/pp3pbp/1Bp3p1/8/2bn4/Q4N2/P4PPP/3R1K1R w - - 0 1","r3r1k1/pp3pbp/1Bp3p1/8/2bn4/Q4N2/P4PPP/3R2KR b - - 0 1","r3r1k1/pp3pbp/1Bp3p1/8/2b5/Q4N2/P3nPPP/3R2KR w - - 0 1","r3r1k1/pp3pbp/1Bp3p1/8/2b5/Q4N2/P3nPPP/3R1K1R b - - 0 1","r3r1k1/pp3pbp/1Bp3p1/8/2b5/Q1n2N2/P4PPP/3R1K1R w - - 0 1","r3r1k1/pp3pbp/1Bp3p1/8/2b5/Q1n2N2/P4PPP/3R2KR b - - 0 1","r3r1k1/1p3pbp/1pp3p1/8/2b5/Q1n2N2/P4PPP/3R2KR w - - 0 1","r3r1k1/1p3pbp/1pp3p1/8/1Qb5/2n2N2/P4PPP/3R2KR b - - 0 1","4r1k1/1p3pbp/1pp3p1/8/rQb5/2n2N2/P4PPP/3R2KR w - - 0 1","4r1k1/1p3pbp/1Qp3p1/8/r1b5/2n2N2/P4PPP/3R2KR b - - 0 1","4r1k1/1p3pbp/1Qp3p1/8/r1b5/5N2/P4PPP/3n2KR w - - 0 1","4r1k1/1p3pbp/1Qp3p1/8/r1b5/5N1P/P4PP1/3n2KR b - - 0 1","4r1k1/1p3pbp/1Qp3p1/8/2b5/5N1P/r4PP1/3n2KR w - - 0 1","4r1k1/1p3pbp/1Qp3p1/8/2b5/5N1P/r4PPK/3n3R b - - 0 1","4r1k1/1p3pbp/1Qp3p1/8/2b5/5N1P/r4nPK/7R w - - 0 1","4r1k1/1p3pbp/1Qp3p1/8/2b5/5N1P/r4nPK/4R3 b - - 0 1","6k1/1p3pbp/1Qp3p1/8/2b5/5N1P/r4nPK/4r3 w - - 0 1","3Q2k1/1p3pbp/2p3p1/8/2b5/5N1P/r4nPK/4r3 b - - 0 1","3Q1bk1/1p3p1p/2p3p1/8/2b5/5N1P/r4nPK/4r3 w - - 0 1","3Q1bk1/1p3p1p/2p3p1/8/2b5/7P/r4nPK/4N3 b - - 0 1","3Q1bk1/1p3p1p/2p3p1/3b4/8/7P/r4nPK/4N3 w - - 0 1","3Q1bk1/1p3p1p/2p3p1/3b4/8/5N1P/r4nPK/8 b - - 0 1","3Q1bk1/1p3p1p/2p3p1/3b4/4n3/5N1P/r5PK/8 w - - 0 1","1Q3bk1/1p3p1p/2p3p1/3b4/4n3/5N1P/r5PK/8 b - - 0 1","1Q3bk1/5p1p/2p3p1/1p1b4/4n3/5N1P/r5PK/8 w - b6 0 1","1Q3bk1/5p1p/2p3p1/1p1b4/4n2P/5N2/r5PK/8 b - - 0 1","1Q3bk1/5p2/2p3p1/1p1b3p/4n2P/5N2/r5PK/8 w - h6 0 1","1Q3bk1/5p2/2p3p1/1p1bN2p/4n2P/8/r5PK/8 b - - 0 1","1Q3b2/5pk1/2p3p1/1p1bN2p/4n2P/8/r5PK/8 w - - 0 1","1Q3b2/5pk1/2p3p1/1p1bN2p/4n2P/8/r5P1/6K1 b - - 0 1","1Q6/5pk1/2p3p1/1pbbN2p/4n2P/8/r5P1/6K1 w - - 0 1","1Q6/5pk1/2p3p1/1pbbN2p/4n2P/8/r5P1/5K2 b - - 0 1","1Q6/5pk1/2p3p1/1pbbN2p/7P/6n1/r5P1/5K2 w - - 0 1","1Q6/5pk1/2p3p1/1pbbN2p/7P/6n1/r5P1/4K3 b - - 0 1","1Q6/5pk1/2p3p1/1p1bN2p/1b5P/6n1/r5P1/4K3 w - - 0 1","1Q6/5pk1/2p3p1/1p1bN2p/1b5P/6n1/r5P1/3K4 b - - 0 1","1Q6/5pk1/2p3p1/1p2N2p/1b5P/1b4n1/r5P1/3K4 w - - 0 1","1Q6/5pk1/2p3p1/1p2N2p/1b5P/1b4n1/r5P1/2K5 b - - 0 1","1Q6/5pk1/2p3p1/1p2N2p/1b5P/1b6/r3n1P1/2K5 w - - 0 1","1Q6/5pk1/2p3p1/1p2N2p/1b5P/1b6/r3n1P1/1K6 b - - 0 1","1Q6/5pk1/2p3p1/1p2N2p/1b5P/1bn5/r5P1/1K6 w - - 0 1","1Q6/5pk1/2p3p1/1p2N2p/1b5P/1bn5/r5P1/2K5 b - - 0 1","1Q6/5pk1/2p3p1/1p2N2p/1b5P/1bn5/2r3P1/2K5 w - - 0 1"],"resultado":"xeque-mate","tabuleiro_final":"8 . Q . . . . . .\n7 . . . . . p k .\n6 . . p . . . p .\n5 . p . . N . . p\n4 . b . . . . . P\n3 . b n . . . . .\n2 . . r . . . P .\n1 . . K . . . . .\n  a b c d e f g h"},"origem":"wiki","palavras_chave":[],"sinonimos":[],"tipo":"conceito","titulo":"A Partida do Século (Byrne × Fischer, 1956)"}
\section{A Partida do Século (Byrne $\times$ Fischer, 1956)}
Fischer, aos 13 anos, sacrifica a dama (Be6!!) e conduz um moinho de peças menores até Tc2\#. A partida que anunciou o gênio — 'a Partida do Século'. Jogadores: Donald Byrne $\times$ Bobby Fischer (1956). A órbita percorre 83 posições e termina em xeque-mate após 82 lances.
\prov{relações: eh\_um orbita\_de\_partida; usa abertura\_xadrez; relaciona xeque\_mate; usa sacrificio\_xadrez; relaciona byrne; relaciona fischer; usa lyapunov\_da\_orbita}
\prov{proveniência: wiki \textperiodcentered{} partidas reais \textperiodcentered{} fato \textperiodcentered{} confiança 1.0 \textperiodcentered{} domínio xadrez \textperiodcentered{} história 24 versões no jornal}

%CRISTAL {"arestas":[{"alvo":"xadrez","confianca":1.0,"epistemico":"fato","origem":"wiki","rel":"parte_de"},{"alvo":"orbita_de_partida","confianca":1.0,"epistemico":"fato","origem":"wiki","rel":"eh_um"},{"alvo":"abertura_xadrez","confianca":1.0,"epistemico":"fato","origem":"wiki","rel":"usa"},{"alvo":"xeque_mate","confianca":1.0,"epistemico":"fato","origem":"wiki","rel":"relaciona"},{"alvo":"sacrificio_xadrez","confianca":1.0,"epistemico":"fato","origem":"wiki","rel":"usa"},{"alvo":"anderssen","confianca":1.0,"epistemico":"fato","origem":"wiki","rel":"relaciona"},{"alvo":"lyapunov_da_orbita","confianca":1.0,"epistemico":"fato","origem":"wiki","rel":"usa"}],"confianca":1.0,"contraexemplos":[],"descricao":"A 'Partida Imortal': Anderssen sacrifica bispo, as DUAS torres e a dama para dar mate com as peças menores. O hino romântico ao ataque acima do material. Jogadores: Adolf Anderssen × Lionel Kieseritzky (1851). A órbita percorre 46 posições e termina em xeque-mate após 45 lances.","epistemico":"fato","exemplos":[],"id":"partida_imortal","memoria":"semantica","meta":{"abertura":"Gambito do Rei","ano":1851,"classe_orbita":"borda","dominio":"xadrez","fen_final":"r1bk3r/p2pBpNp/n4n2/1p1NP2P/6P1/3P4/P1P1K3/q5b1 b - - 0 1","fonte":"partidas reais","jogadores":"Adolf Anderssen × Lionel Kieseritzky","lances_san":["e4","e5","f4","exf4","Bc4","Qh4+","Kf1","b5","Bxb5","Nf6","Nf3","Qh6","d3","Nh5","Nh4","Qg5","Nf5","c6","g4","Nf6","Rg1","cxb5","h4","Qg6","h5","Qg5","Qf3","Ng8","Bxf4","Qf6","Nc3","Bc5","Nd5","Qxb2","Bd6","Bxg1","e5","Qxa1+","Ke2","Na6","Nxg7+","Kd8","Qf6+","Nxf6","Be7#"],"lyapunov":0.3435,"mineral_H":0.1195,"mineral_classe":"cristal","mineral_theta_g":180.0,"n_lances":45,"orbita_fen":["rnbqkbnr/pppppppp/8/8/8/8/PPPPPPPP/RNBQKBNR w - - 0 1","rnbqkbnr/pppppppp/8/8/4P3/8/PPPP1PPP/RNBQKBNR b - e3 0 1","rnbqkbnr/pppp1ppp/8/4p3/4P3/8/PPPP1PPP/RNBQKBNR w - e6 0 1","rnbqkbnr/pppp1ppp/8/4p3/4PP2/8/PPPP2PP/RNBQKBNR b - f3 0 1","rnbqkbnr/pppp1ppp/8/8/4Pp2/8/PPPP2PP/RNBQKBNR w - - 0 1","rnbqkbnr/pppp1ppp/8/8/2B1Pp2/8/PPPP2PP/RNBQK1NR b - - 0 1","rnb1kbnr/pppp1ppp/8/8/2B1Pp1q/8/PPPP2PP/RNBQK1NR w - - 0 1","rnb1kbnr/pppp1ppp/8/8/2B1Pp1q/8/PPPP2PP/RNBQ1KNR b - - 0 1","rnb1kbnr/p1pp1ppp/8/1p6/2B1Pp1q/8/PPPP2PP/RNBQ1KNR w - b6 0 1","rnb1kbnr/p1pp1ppp/8/1B6/4Pp1q/8/PPPP2PP/RNBQ1KNR b - - 0 1","rnb1kb1r/p1pp1ppp/5n2/1B6/4Pp1q/8/PPPP2PP/RNBQ1KNR w - - 0 1","rnb1kb1r/p1pp1ppp/5n2/1B6/4Pp1q/5N2/PPPP2PP/RNBQ1K1R b - - 0 1","rnb1kb1r/p1pp1ppp/5n1q/1B6/4Pp2/5N2/PPPP2PP/RNBQ1K1R w - - 0 1","rnb1kb1r/p1pp1ppp/5n1q/1B6/4Pp2/3P1N2/PPP3PP/RNBQ1K1R b - - 0 1","rnb1kb1r/p1pp1ppp/7q/1B5n/4Pp2/3P1N2/PPP3PP/RNBQ1K1R w - - 0 1","rnb1kb1r/p1pp1ppp/7q/1B5n/4Pp1N/3P4/PPP3PP/RNBQ1K1R b - - 0 1","rnb1kb1r/p1pp1ppp/8/1B4qn/4Pp1N/3P4/PPP3PP/RNBQ1K1R w - - 0 1","rnb1kb1r/p1pp1ppp/8/1B3Nqn/4Pp2/3P4/PPP3PP/RNBQ1K1R b - - 0 1","rnb1kb1r/p2p1ppp/2p5/1B3Nqn/4Pp2/3P4/PPP3PP/RNBQ1K1R w - - 0 1","rnb1kb1r/p2p1ppp/2p5/1B3Nqn/4PpP1/3P4/PPP4P/RNBQ1K1R b - g3 0 1","rnb1kb1r/p2p1ppp/2p2n2/1B3Nq1/4PpP1/3P4/PPP4P/RNBQ1K1R w - - 0 1","rnb1kb1r/p2p1ppp/2p2n2/1B3Nq1/4PpP1/3P4/PPP4P/RNBQ1KR1 b - - 0 1","rnb1kb1r/p2p1ppp/5n2/1p3Nq1/4PpP1/3P4/PPP4P/RNBQ1KR1 w - - 0 1","rnb1kb1r/p2p1ppp/5n2/1p3Nq1/4PpPP/3P4/PPP5/RNBQ1KR1 b - h3 0 1","rnb1kb1r/p2p1ppp/5nq1/1p3N2/4PpPP/3P4/PPP5/RNBQ1KR1 w - - 0 1","rnb1kb1r/p2p1ppp/5nq1/1p3N1P/4PpP1/3P4/PPP5/RNBQ1KR1 b - - 0 1","rnb1kb1r/p2p1ppp/5n2/1p3NqP/4PpP1/3P4/PPP5/RNBQ1KR1 w - - 0 1","rnb1kb1r/p2p1ppp/5n2/1p3NqP/4PpP1/3P1Q2/PPP5/RNB2KR1 b - - 0 1","rnb1kbnr/p2p1ppp/8/1p3NqP/4PpP1/3P1Q2/PPP5/RNB2KR1 w - - 0 1","rnb1kbnr/p2p1ppp/8/1p3NqP/4PBP1/3P1Q2/PPP5/RN3KR1 b - - 0 1","rnb1kbnr/p2p1ppp/5q2/1p3N1P/4PBP1/3P1Q2/PPP5/RN3KR1 w - - 0 1","rnb1kbnr/p2p1ppp/5q2/1p3N1P/4PBP1/2NP1Q2/PPP5/R4KR1 b - - 0 1","rnb1k1nr/p2p1ppp/5q2/1pb2N1P/4PBP1/2NP1Q2/PPP5/R4KR1 w - - 0 1","rnb1k1nr/p2p1ppp/5q2/1pbN1N1P/4PBP1/3P1Q2/PPP5/R4KR1 b - - 0 1","rnb1k1nr/p2p1ppp/8/1pbN1N1P/4PBP1/3P1Q2/PqP5/R4KR1 w - - 0 1","rnb1k1nr/p2p1ppp/3B4/1pbN1N1P/4P1P1/3P1Q2/PqP5/R4KR1 b - - 0 1","rnb1k1nr/p2p1ppp/3B4/1p1N1N1P/4P1P1/3P1Q2/PqP5/R4Kb1 w - - 0 1","rnb1k1nr/p2p1ppp/3B4/1p1NPN1P/6P1/3P1Q2/PqP5/R4Kb1 b - - 0 1","rnb1k1nr/p2p1ppp/3B4/1p1NPN1P/6P1/3P1Q2/P1P5/q4Kb1 w - - 0 1","rnb1k1nr/p2p1ppp/3B4/1p1NPN1P/6P1/3P1Q2/P1P1K3/q5b1 b - - 0 1","r1b1k1nr/p2p1ppp/n2B4/1p1NPN1P/6P1/3P1Q2/P1P1K3/q5b1 w - - 0 1","r1b1k1nr/p2p1pNp/n2B4/1p1NP2P/6P1/3P1Q2/P1P1K3/q5b1 b - - 0 1","r1bk2nr/p2p1pNp/n2B4/1p1NP2P/6P1/3P1Q2/P1P1K3/q5b1 w - - 0 1","r1bk2nr/p2p1pNp/n2B1Q2/1p1NP2P/6P1/3P4/P1P1K3/q5b1 b - - 0 1","r1bk3r/p2p1pNp/n2B1n2/1p1NP2P/6P1/3P4/P1P1K3/q5b1 w - - 0 1","r1bk3r/p2pBpNp/n4n2/1p1NP2P/6P1/3P4/P1P1K3/q5b1 b - - 0 1"],"resultado":"xeque-mate","tabuleiro_final":"8 r . b k . . . r\n7 p . . p B p N p\n6 n . . . . n . .\n5 . p . N P . . P\n4 . . . . . . P .\n3 . . . P . . . .\n2 P . P . K . . .\n1 q . . . . . b .\n  a b c d e f g h"},"origem":"wiki","palavras_chave":[],"sinonimos":[],"tipo":"conceito","titulo":"A Imortal (Anderssen, 1851)"}
\section{A Imortal (Anderssen, 1851)}
A 'Partida Imortal': Anderssen sacrifica bispo, as DUAS torres e a dama para dar mate com as peças menores. O hino romântico ao ataque acima do material. Jogadores: Adolf Anderssen $\times$ Lionel Kieseritzky (1851). A órbita percorre 46 posições e termina em xeque-mate após 45 lances.
\prov{relações: parte\_de xadrez; eh\_um orbita\_de\_partida; usa abertura\_xadrez; relaciona xeque\_mate; usa sacrificio\_xadrez; relaciona anderssen; usa lyapunov\_da\_orbita}
\prov{proveniência: wiki \textperiodcentered{} partidas reais \textperiodcentered{} fato \textperiodcentered{} confiança 1.0 \textperiodcentered{} domínio xadrez \textperiodcentered{} história 60 versões no jornal}

%CRISTAL {"arestas":[{"alvo":"orbita_de_partida","confianca":1.0,"epistemico":"fato","origem":"wiki","rel":"eh_um"},{"alvo":"abertura_xadrez","confianca":1.0,"epistemico":"fato","origem":"wiki","rel":"usa"},{"alvo":"xeque_mate","confianca":1.0,"epistemico":"fato","origem":"wiki","rel":"relaciona"},{"alvo":"sacrificio_xadrez","confianca":1.0,"epistemico":"fato","origem":"wiki","rel":"usa"},{"alvo":"edward_lasker","confianca":1.0,"epistemico":"fato","origem":"wiki","rel":"relaciona"},{"alvo":"lyapunov_da_orbita","confianca":1.0,"epistemico":"fato","origem":"wiki","rel":"usa"}],"confianca":1.0,"contraexemplos":[],"descricao":"Sacrifício de dama (Dxh7+!!) e o rei preto é arrastado do h7 até o g1, onde leva mate — uma caminhada forçada de sete casas, Kd2#. O rei-hunt mais famoso. Jogadores: Edward Lasker × George Thomas (1912). A órbita percorre 36 posições e termina em xeque-mate após 35 lances.","epistemico":"fato","exemplos":[],"id":"partida_lasker_thomas","memoria":"semantica","meta":{"abertura":"Holandesa","ano":1912,"classe_orbita":"borda","dominio":"xadrez","fen_final":"rn3r2/pbppq1p1/1p2pN2/8/3P2NP/6P1/PPPKBP1R/R5k1 b - - 0 1","fonte":"partidas reais","jogadores":"Edward Lasker × George Thomas","lances_san":["d4","e6","Nf3","f5","Nc3","Nf6","Bg5","Be7","Bxf6","Bxf6","e4","fxe4","Nxe4","b6","Ne5","O-O","Bd3","Bb7","Qh5","Qe7","Qxh7+","Kxh7","Nxf6+","Kh6","Neg4+","Kg5","h4+","Kf4","g3+","Kf3","Be2+","Kg2","Rh2+","Kg1","Kd2#"],"lyapunov":0.3399,"mineral_H":-0.0358,"mineral_classe":"cristal","mineral_theta_g":180.0,"n_lances":35,"orbita_fen":["rnbqkbnr/pppppppp/8/8/8/8/PPPPPPPP/RNBQKBNR w - - 0 1","rnbqkbnr/pppppppp/8/8/3P4/8/PPP1PPPP/RNBQKBNR b - d3 0 1","rnbqkbnr/pppp1ppp/4p3/8/3P4/8/PPP1PPPP/RNBQKBNR w - - 0 1","rnbqkbnr/pppp1ppp/4p3/8/3P4/5N2/PPP1PPPP/RNBQKB1R b - - 0 1","rnbqkbnr/pppp2pp/4p3/5p2/3P4/5N2/PPP1PPPP/RNBQKB1R w - f6 0 1","rnbqkbnr/pppp2pp/4p3/5p2/3P4/2N2N2/PPP1PPPP/R1BQKB1R b - - 0 1","rnbqkb1r/pppp2pp/4pn2/5p2/3P4/2N2N2/PPP1PPPP/R1BQKB1R w - - 0 1","rnbqkb1r/pppp2pp/4pn2/5pB1/3P4/2N2N2/PPP1PPPP/R2QKB1R b - - 0 1","rnbqk2r/ppppb1pp/4pn2/5pB1/3P4/2N2N2/PPP1PPPP/R2QKB1R w - - 0 1","rnbqk2r/ppppb1pp/4pB2/5p2/3P4/2N2N2/PPP1PPPP/R2QKB1R b - - 0 1","rnbqk2r/pppp2pp/4pb2/5p2/3P4/2N2N2/PPP1PPPP/R2QKB1R w - - 0 1","rnbqk2r/pppp2pp/4pb2/5p2/3PP3/2N2N2/PPP2PPP/R2QKB1R b - e3 0 1","rnbqk2r/pppp2pp/4pb2/8/3Pp3/2N2N2/PPP2PPP/R2QKB1R w - - 0 1","rnbqk2r/pppp2pp/4pb2/8/3PN3/5N2/PPP2PPP/R2QKB1R b - - 0 1","rnbqk2r/p1pp2pp/1p2pb2/8/3PN3/5N2/PPP2PPP/R2QKB1R w - - 0 1","rnbqk2r/p1pp2pp/1p2pb2/4N3/3PN3/8/PPP2PPP/R2QKB1R b - - 0 1","rnbq1rk1/p1pp2pp/1p2pb2/4N3/3PN3/8/PPP2PPP/R2QKB1R w - - 0 1","rnbq1rk1/p1pp2pp/1p2pb2/4N3/3PN3/3B4/PPP2PPP/R2QK2R b - - 0 1","rn1q1rk1/pbpp2pp/1p2pb2/4N3/3PN3/3B4/PPP2PPP/R2QK2R w - - 0 1","rn1q1rk1/pbpp2pp/1p2pb2/4N2Q/3PN3/3B4/PPP2PPP/R3K2R b - - 0 1","rn3rk1/pbppq1pp/1p2pb2/4N2Q/3PN3/3B4/PPP2PPP/R3K2R w - - 0 1","rn3rk1/pbppq1pQ/1p2pb2/4N3/3PN3/3B4/PPP2PPP/R3K2R b - - 0 1","rn3r2/pbppq1pk/1p2pb2/4N3/3PN3/3B4/PPP2PPP/R3K2R w - - 0 1","rn3r2/pbppq1pk/1p2pN2/4N3/3P4/3B4/PPP2PPP/R3K2R b - - 0 1","rn3r2/pbppq1p1/1p2pN1k/4N3/3P4/3B4/PPP2PPP/R3K2R w - - 0 1","rn3r2/pbppq1p1/1p2pN1k/8/3P2N1/3B4/PPP2PPP/R3K2R b - - 0 1","rn3r2/pbppq1p1/1p2pN2/6k1/3P2N1/3B4/PPP2PPP/R3K2R w - - 0 1","rn3r2/pbppq1p1/1p2pN2/6k1/3P2NP/3B4/PPP2PP1/R3K2R b - h3 0 1","rn3r2/pbppq1p1/1p2pN2/8/3P1kNP/3B4/PPP2PP1/R3K2R w - - 0 1","rn3r2/pbppq1p1/1p2pN2/8/3P1kNP/3B2P1/PPP2P2/R3K2R b - - 0 1","rn3r2/pbppq1p1/1p2pN2/8/3P2NP/3B1kP1/PPP2P2/R3K2R w - - 0 1","rn3r2/pbppq1p1/1p2pN2/8/3P2NP/5kP1/PPP1BP2/R3K2R b - - 0 1","rn3r2/pbppq1p1/1p2pN2/8/3P2NP/6P1/PPP1BPk1/R3K2R w - - 0 1","rn3r2/pbppq1p1/1p2pN2/8/3P2NP/6P1/PPP1BPkR/R3K3 b - - 0 1","rn3r2/pbppq1p1/1p2pN2/8/3P2NP/6P1/PPP1BP1R/R3K1k1 w - - 0 1","rn3r2/pbppq1p1/1p2pN2/8/3P2NP/6P1/PPPKBP1R/R5k1 b - - 0 1"],"resultado":"xeque-mate","tabuleiro_final":"8 r n . . . r . .\n7 p b p p q . p .\n6 . p . . p N . .\n5 . . . . . . . .\n4 . . . P . . N P\n3 . . . . . . P .\n2 P P P K B P . R\n1 R . . . . . k .\n  a b c d e f g h"},"origem":"wiki","palavras_chave":[],"sinonimos":[],"tipo":"conceito","titulo":"A Caminhada do Rei (Ed. Lasker × Thomas, 1912)"}
\section{A Caminhada do Rei (Ed. Lasker $\times$ Thomas, 1912)}
Sacrifício de dama (Dxh7+!!) e o rei preto é arrastado do h7 até o g1, onde leva mate — uma caminhada forçada de sete casas, Kd2\#. O rei-hunt mais famoso. Jogadores: Edward Lasker $\times$ George Thomas (1912). A órbita percorre 36 posições e termina em xeque-mate após 35 lances.
\prov{relações: eh\_um orbita\_de\_partida; usa abertura\_xadrez; relaciona xeque\_mate; usa sacrificio\_xadrez; relaciona edward\_lasker; usa lyapunov\_da\_orbita}
\prov{proveniência: wiki \textperiodcentered{} partidas reais \textperiodcentered{} fato \textperiodcentered{} confiança 1.0 \textperiodcentered{} domínio xadrez \textperiodcentered{} história 14 versões no jornal}

%CRISTAL {"arestas":[{"alvo":"orbita_de_partida","confianca":1.0,"epistemico":"fato","origem":"wiki","rel":"eh_um"},{"alvo":"abertura_xadrez","confianca":1.0,"epistemico":"fato","origem":"wiki","rel":"usa"},{"alvo":"xeque_mate","confianca":1.0,"epistemico":"fato","origem":"wiki","rel":"relaciona"},{"alvo":"sacrificio_xadrez","confianca":1.0,"epistemico":"fato","origem":"wiki","rel":"usa"},{"alvo":"lyapunov_da_orbita","confianca":1.0,"epistemico":"fato","origem":"wiki","rel":"usa"}],"confianca":1.0,"contraexemplos":[],"descricao":"O 'pseudo-sacrifício de Légal': ignora-se a dama (Nxe5) porque o mate com as peças menores vem primeiro. O padrão de mate mais antigo com nome próprio. Jogadores: Sire de Légal × Saint Brie (1750). A órbita percorre 14 posições e termina em xeque-mate após 13 lances.","epistemico":"fato","exemplos":[],"id":"partida_legal","memoria":"semantica","meta":{"abertura":"Abertura Italiana","ano":1750,"classe_orbita":"borda","dominio":"xadrez","fen_final":"rn1q1bnr/ppp1kB1p/3p2p1/3NN3/4P3/8/PPPP1PPP/R1BbK2R b - - 0 1","fonte":"partidas reais","jogadores":"Sire de Légal × Saint Brie","lances_san":["e4","e5","Bc4","d6","Nf3","Bg4","Nc3","g6","Nxe5","Bxd1","Bxf7+","Ke7","Nd5#"],"lyapunov":0.4445,"mineral_H":0.4643,"mineral_classe":"borda","mineral_theta_g":48.7,"n_lances":13,"orbita_fen":["rnbqkbnr/pppppppp/8/8/8/8/PPPPPPPP/RNBQKBNR w - - 0 1","rnbqkbnr/pppppppp/8/8/4P3/8/PPPP1PPP/RNBQKBNR b - e3 0 1","rnbqkbnr/pppp1ppp/8/4p3/4P3/8/PPPP1PPP/RNBQKBNR w - e6 0 1","rnbqkbnr/pppp1ppp/8/4p3/2B1P3/8/PPPP1PPP/RNBQK1NR b - - 0 1","rnbqkbnr/ppp2ppp/3p4/4p3/2B1P3/8/PPPP1PPP/RNBQK1NR w - - 0 1","rnbqkbnr/ppp2ppp/3p4/4p3/2B1P3/5N2/PPPP1PPP/RNBQK2R b - - 0 1","rn1qkbnr/ppp2ppp/3p4/4p3/2B1P1b1/5N2/PPPP1PPP/RNBQK2R w - - 0 1","rn1qkbnr/ppp2ppp/3p4/4p3/2B1P1b1/2N2N2/PPPP1PPP/R1BQK2R b - - 0 1","rn1qkbnr/ppp2p1p/3p2p1/4p3/2B1P1b1/2N2N2/PPPP1PPP/R1BQK2R w - - 0 1","rn1qkbnr/ppp2p1p/3p2p1/4N3/2B1P1b1/2N5/PPPP1PPP/R1BQK2R b - - 0 1","rn1qkbnr/ppp2p1p/3p2p1/4N3/2B1P3/2N5/PPPP1PPP/R1BbK2R w - - 0 1","rn1qkbnr/ppp2B1p/3p2p1/4N3/4P3/2N5/PPPP1PPP/R1BbK2R b - - 0 1","rn1q1bnr/ppp1kB1p/3p2p1/4N3/4P3/2N5/PPPP1PPP/R1BbK2R w - - 0 1","rn1q1bnr/ppp1kB1p/3p2p1/3NN3/4P3/8/PPPP1PPP/R1BbK2R b - - 0 1"],"resultado":"xeque-mate","tabuleiro_final":"8 r n . q . b n r\n7 p p p . k B . p\n6 . . . p . . p .\n5 . . . N N . . .\n4 . . . . P . . .\n3 . . . . . . . .\n2 P P P P . P P P\n1 R . B b K . . R\n  a b c d e f g h"},"origem":"wiki","palavras_chave":[],"sinonimos":[],"tipo":"conceito","titulo":"Mate de Légal (1750)"}
\section{Mate de Légal (1750)}
O 'pseudo-sacrifício de Légal': ignora-se a dama (Nxe5) porque o mate com as peças menores vem primeiro. O padrão de mate mais antigo com nome próprio. Jogadores: Sire de Légal $\times$ Saint Brie (1750). A órbita percorre 14 posições e termina em xeque-mate após 13 lances.
\prov{relações: eh\_um orbita\_de\_partida; usa abertura\_xadrez; relaciona xeque\_mate; usa sacrificio\_xadrez; usa lyapunov\_da\_orbita}
\prov{proveniência: wiki \textperiodcentered{} partidas reais \textperiodcentered{} fato \textperiodcentered{} confiança 1.0 \textperiodcentered{} domínio xadrez \textperiodcentered{} história 48 versões no jornal}

%CRISTAL {"arestas":[{"alvo":"orbita_de_partida","confianca":1.0,"epistemico":"fato","origem":"wiki","rel":"eh_um"},{"alvo":"abertura_xadrez","confianca":1.0,"epistemico":"fato","origem":"wiki","rel":"usa"},{"alvo":"xeque_mate","confianca":1.0,"epistemico":"fato","origem":"wiki","rel":"relaciona"},{"alvo":"sacrificio_xadrez","confianca":1.0,"epistemico":"fato","origem":"wiki","rel":"usa"},{"alvo":"cravada","confianca":1.0,"epistemico":"fato","origem":"wiki","rel":"usa"},{"alvo":"morphy","confianca":1.0,"epistemico":"fato","origem":"wiki","rel":"relaciona"},{"alvo":"lyapunov_da_orbita","confianca":1.0,"epistemico":"fato","origem":"wiki","rel":"usa"}],"confianca":1.0,"contraexemplos":[],"descricao":"A partida mais famosa da história, jogada num camarote da ópera em Paris. Morphy desenvolve cada peça com propósito e conclui com o sacrifício da dama (Db8+!!) seguido de Td8#: desenvolvimento e iniciativa levados ao mate. Jogadores: Paul Morphy × Duque de Brunswick e Conde Isouard (1858). A órbita percorre 34 posições e termina em xeque-mate após 33 lances.","epistemico":"fato","exemplos":[],"id":"partida_opera","memoria":"semantica","meta":{"abertura":"Defesa Philidor","ano":1858,"classe_orbita":"caos","dominio":"xadrez","fen_final":"1n1Rkb1r/p4ppp/4q3/4p1B1/4P3/8/PPP2PPP/2K5 b - - 0 1","fonte":"partidas reais","jogadores":"Paul Morphy × Duque de Brunswick e Conde Isouard","lances_san":["e4","e5","Nf3","d6","d4","Bg4","dxe5","Bxf3","Qxf3","dxe5","Bc4","Nf6","Qb3","Qe7","Nc3","c6","Bg5","b5","Nxb5","cxb5","Bxb5+","Nbd7","O-O-O","Rd8","Rxd7","Rxd7","Rd1","Qe6","Bxd7+","Nxd7","Qb8+","Nxb8","Rd8#"],"lyapunov":0.5832,"mineral_H":0.1588,"mineral_classe":"cristal","mineral_theta_g":180.0,"n_lances":33,"orbita_fen":["rnbqkbnr/pppppppp/8/8/8/8/PPPPPPPP/RNBQKBNR w - - 0 1","rnbqkbnr/pppppppp/8/8/4P3/8/PPPP1PPP/RNBQKBNR b - e3 0 1","rnbqkbnr/pppp1ppp/8/4p3/4P3/8/PPPP1PPP/RNBQKBNR w - e6 0 1","rnbqkbnr/pppp1ppp/8/4p3/4P3/5N2/PPPP1PPP/RNBQKB1R b - - 0 1","rnbqkbnr/ppp2ppp/3p4/4p3/4P3/5N2/PPPP1PPP/RNBQKB1R w - - 0 1","rnbqkbnr/ppp2ppp/3p4/4p3/3PP3/5N2/PPP2PPP/RNBQKB1R b - d3 0 1","rn1qkbnr/ppp2ppp/3p4/4p3/3PP1b1/5N2/PPP2PPP/RNBQKB1R w - - 0 1","rn1qkbnr/ppp2ppp/3p4/4P3/4P1b1/5N2/PPP2PPP/RNBQKB1R b - - 0 1","rn1qkbnr/ppp2ppp/3p4/4P3/4P3/5b2/PPP2PPP/RNBQKB1R w - - 0 1","rn1qkbnr/ppp2ppp/3p4/4P3/4P3/5Q2/PPP2PPP/RNB1KB1R b - - 0 1","rn1qkbnr/ppp2ppp/8/4p3/4P3/5Q2/PPP2PPP/RNB1KB1R w - - 0 1","rn1qkbnr/ppp2ppp/8/4p3/2B1P3/5Q2/PPP2PPP/RNB1K2R b - - 0 1","rn1qkb1r/ppp2ppp/5n2/4p3/2B1P3/5Q2/PPP2PPP/RNB1K2R w - - 0 1","rn1qkb1r/ppp2ppp/5n2/4p3/2B1P3/1Q6/PPP2PPP/RNB1K2R b - - 0 1","rn2kb1r/ppp1qppp/5n2/4p3/2B1P3/1Q6/PPP2PPP/RNB1K2R w - - 0 1","rn2kb1r/ppp1qppp/5n2/4p3/2B1P3/1QN5/PPP2PPP/R1B1K2R b - - 0 1","rn2kb1r/pp2qppp/2p2n2/4p3/2B1P3/1QN5/PPP2PPP/R1B1K2R w - - 0 1","rn2kb1r/pp2qppp/2p2n2/4p1B1/2B1P3/1QN5/PPP2PPP/R3K2R b - - 0 1","rn2kb1r/p3qppp/2p2n2/1p2p1B1/2B1P3/1QN5/PPP2PPP/R3K2R w - b6 0 1","rn2kb1r/p3qppp/2p2n2/1N2p1B1/2B1P3/1Q6/PPP2PPP/R3K2R b - - 0 1","rn2kb1r/p3qppp/5n2/1p2p1B1/2B1P3/1Q6/PPP2PPP/R3K2R w - - 0 1","rn2kb1r/p3qppp/5n2/1B2p1B1/4P3/1Q6/PPP2PPP/R3K2R b - - 0 1","r3kb1r/p2nqppp/5n2/1B2p1B1/4P3/1Q6/PPP2PPP/R3K2R w - - 0 1","r3kb1r/p2nqppp/5n2/1B2p1B1/4P3/1Q6/PPP2PPP/2KR3R b - - 0 1","3rkb1r/p2nqppp/5n2/1B2p1B1/4P3/1Q6/PPP2PPP/2KR3R w - - 0 1","3rkb1r/p2Rqppp/5n2/1B2p1B1/4P3/1Q6/PPP2PPP/2K4R b - - 0 1","4kb1r/p2rqppp/5n2/1B2p1B1/4P3/1Q6/PPP2PPP/2K4R w - - 0 1","4kb1r/p2rqppp/5n2/1B2p1B1/4P3/1Q6/PPP2PPP/2KR4 b - - 0 1","4kb1r/p2r1ppp/4qn2/1B2p1B1/4P3/1Q6/PPP2PPP/2KR4 w - - 0 1","4kb1r/p2B1ppp/4qn2/4p1B1/4P3/1Q6/PPP2PPP/2KR4 b - - 0 1","4kb1r/p2n1ppp/4q3/4p1B1/4P3/1Q6/PPP2PPP/2KR4 w - - 0 1","1Q2kb1r/p2n1ppp/4q3/4p1B1/4P3/8/PPP2PPP/2KR4 b - - 0 1","1n2kb1r/p4ppp/4q3/4p1B1/4P3/8/PPP2PPP/2KR4 w - - 0 1","1n1Rkb1r/p4ppp/4q3/4p1B1/4P3/8/PPP2PPP/2K5 b - - 0 1"],"resultado":"xeque-mate","tabuleiro_final":"8 . n . R k b . r\n7 p . . . . p p p\n6 . . . . q . . .\n5 . . . . p . B .\n4 . . . . P . . .\n3 . . . . . . . .\n2 P P P . . P P P\n1 . . K . . . . .\n  a b c d e f g h"},"origem":"wiki","palavras_chave":[],"sinonimos":[],"tipo":"conceito","titulo":"Partida da Ópera (Morphy, 1858)"}
\section{Partida da Ópera (Morphy, 1858)}
A partida mais famosa da história, jogada num camarote da ópera em Paris. Morphy desenvolve cada peça com propósito e conclui com o sacrifício da dama (Db8+!!) seguido de Td8\#: desenvolvimento e iniciativa levados ao mate. Jogadores: Paul Morphy $\times$ Duque de Brunswick e Conde Isouard (1858). A órbita percorre 34 posições e termina em xeque-mate após 33 lances.
\prov{relações: eh\_um orbita\_de\_partida; usa abertura\_xadrez; relaciona xeque\_mate; usa sacrificio\_xadrez; usa cravada; relaciona morphy; usa lyapunov\_da\_orbita}
\prov{proveniência: wiki \textperiodcentered{} partidas reais \textperiodcentered{} fato \textperiodcentered{} confiança 1.0 \textperiodcentered{} domínio xadrez \textperiodcentered{} história 66 versões no jornal}

%CRISTAL {"arestas":[{"alvo":"orbita_de_partida","confianca":1.0,"epistemico":"fato","origem":"wiki","rel":"eh_um"},{"alvo":"abertura_xadrez","confianca":1.0,"epistemico":"fato","origem":"wiki","rel":"usa"},{"alvo":"xeque_mate","confianca":1.0,"epistemico":"fato","origem":"wiki","rel":"relaciona"},{"alvo":"tatica_xadrez","confianca":1.0,"epistemico":"fato","origem":"wiki","rel":"usa"},{"alvo":"lyapunov_da_orbita","confianca":1.0,"epistemico":"fato","origem":"wiki","rel":"usa"}],"confianca":1.0,"contraexemplos":[],"descricao":"O mate em quatro que todo iniciante aprende (e teme): a dama e o bispo miram f7, o ponto mais fraco, e a dama entra em Dxf7#. Jogadores: (padrão didático). A órbita percorre 8 posições e termina em xeque-mate após 7 lances.","epistemico":"fato","exemplos":[],"id":"partida_pastor","memoria":"semantica","meta":{"abertura":"Abertura Italiana","ano":null,"classe_orbita":"borda","dominio":"xadrez","fen_final":"r1bqkb1r/pppp1Qpp/2n2n2/4p3/2B1P3/8/PPPP1PPP/RNB1K1NR b - - 0 1","fonte":"partidas reais","jogadores":"(padrão didático)","lances_san":["e4","e5","Bc4","Nc6","Qh5","Nf6","Qxf7#"],"lyapunov":0.318,"mineral_H":0.8238,"mineral_classe":"caos","mineral_theta_g":34.3,"n_lances":7,"orbita_fen":["rnbqkbnr/pppppppp/8/8/8/8/PPPPPPPP/RNBQKBNR w - - 0 1","rnbqkbnr/pppppppp/8/8/4P3/8/PPPP1PPP/RNBQKBNR b - e3 0 1","rnbqkbnr/pppp1ppp/8/4p3/4P3/8/PPPP1PPP/RNBQKBNR w - e6 0 1","rnbqkbnr/pppp1ppp/8/4p3/2B1P3/8/PPPP1PPP/RNBQK1NR b - - 0 1","r1bqkbnr/pppp1ppp/2n5/4p3/2B1P3/8/PPPP1PPP/RNBQK1NR w - - 0 1","r1bqkbnr/pppp1ppp/2n5/4p2Q/2B1P3/8/PPPP1PPP/RNB1K1NR b - - 0 1","r1bqkb1r/pppp1ppp/2n2n2/4p2Q/2B1P3/8/PPPP1PPP/RNB1K1NR w - - 0 1","r1bqkb1r/pppp1Qpp/2n2n2/4p3/2B1P3/8/PPPP1PPP/RNB1K1NR b - - 0 1"],"resultado":"xeque-mate","tabuleiro_final":"8 r . b q k b . r\n7 p p p p . Q p p\n6 . . n . . n . .\n5 . . . . p . . .\n4 . . B . P . . .\n3 . . . . . . . .\n2 P P P P . P P P\n1 R N B . K . N R\n  a b c d e f g h"},"origem":"wiki","palavras_chave":[],"sinonimos":[],"tipo":"conceito","titulo":"Mate do Pastor"}
\section{Mate do Pastor}
O mate em quatro que todo iniciante aprende (e teme): a dama e o bispo miram f7, o ponto mais fraco, e a dama entra em Dxf7\#. Jogadores: (padrão didático). A órbita percorre 8 posições e termina em xeque-mate após 7 lances.
\prov{relações: eh\_um orbita\_de\_partida; usa abertura\_xadrez; relaciona xeque\_mate; usa tatica\_xadrez; usa lyapunov\_da\_orbita}
\prov{proveniência: wiki \textperiodcentered{} partidas reais \textperiodcentered{} fato \textperiodcentered{} confiança 1.0 \textperiodcentered{} domínio xadrez \textperiodcentered{} história 48 versões no jornal}

%CRISTAL {"arestas":[{"alvo":"orbita_de_partida","confianca":1.0,"epistemico":"fato","origem":"wiki","rel":"eh_um"},{"alvo":"abertura_xadrez","confianca":1.0,"epistemico":"fato","origem":"wiki","rel":"usa"},{"alvo":"xeque_mate","confianca":1.0,"epistemico":"fato","origem":"wiki","rel":"relaciona"},{"alvo":"sacrificio_xadrez","confianca":1.0,"epistemico":"fato","origem":"wiki","rel":"usa"},{"alvo":"canal","confianca":1.0,"epistemico":"fato","origem":"wiki","rel":"relaciona"},{"alvo":"lyapunov_da_orbita","confianca":1.0,"epistemico":"fato","origem":"wiki","rel":"usa"}],"confianca":1.0,"contraexemplos":[],"descricao":"Canal sacrifica torre, dama e a outra torre numa simultânea e conclui com Ba6# — o bispo dá mate sozinho no rei preto engaiolado. A 'Imortal Peruana'. Jogadores: Esteban Canal × N.N. (1934). A órbita percorre 28 posições e termina em xeque-mate após 27 lances.","epistemico":"fato","exemplos":[],"id":"partida_peruana","memoria":"semantica","meta":{"abertura":"Defesa Escandinava","ano":1934,"classe_orbita":"caos","dominio":"xadrez","fen_final":"2kr2nr/p2n1ppp/B1p1p3/8/1P1P1B2/2N4P/1PPK1PP1/7q b - - 0 1","fonte":"partidas reais","jogadores":"Esteban Canal × N.N.","lances_san":["e4","d5","exd5","Qxd5","Nc3","Qa5","d4","c6","Nf3","Bg4","Bf4","e6","h3","Bxf3","Qxf3","Bb4","Be2","Nd7","a3","O-O-O","axb4","Qxa1+","Kd2","Qxh1","Qxc6+","bxc6","Ba6#"],"lyapunov":0.5574,"mineral_H":0.1976,"mineral_classe":"cristal","mineral_theta_g":180.0,"n_lances":27,"orbita_fen":["rnbqkbnr/pppppppp/8/8/8/8/PPPPPPPP/RNBQKBNR w - - 0 1","rnbqkbnr/pppppppp/8/8/4P3/8/PPPP1PPP/RNBQKBNR b - e3 0 1","rnbqkbnr/ppp1pppp/8/3p4/4P3/8/PPPP1PPP/RNBQKBNR w - d6 0 1","rnbqkbnr/ppp1pppp/8/3P4/8/8/PPPP1PPP/RNBQKBNR b - - 0 1","rnb1kbnr/ppp1pppp/8/3q4/8/8/PPPP1PPP/RNBQKBNR w - - 0 1","rnb1kbnr/ppp1pppp/8/3q4/8/2N5/PPPP1PPP/R1BQKBNR b - - 0 1","rnb1kbnr/ppp1pppp/8/q7/8/2N5/PPPP1PPP/R1BQKBNR w - - 0 1","rnb1kbnr/ppp1pppp/8/q7/3P4/2N5/PPP2PPP/R1BQKBNR b - d3 0 1","rnb1kbnr/pp2pppp/2p5/q7/3P4/2N5/PPP2PPP/R1BQKBNR w - - 0 1","rnb1kbnr/pp2pppp/2p5/q7/3P4/2N2N2/PPP2PPP/R1BQKB1R b - - 0 1","rn2kbnr/pp2pppp/2p5/q7/3P2b1/2N2N2/PPP2PPP/R1BQKB1R w - - 0 1","rn2kbnr/pp2pppp/2p5/q7/3P1Bb1/2N2N2/PPP2PPP/R2QKB1R b - - 0 1","rn2kbnr/pp3ppp/2p1p3/q7/3P1Bb1/2N2N2/PPP2PPP/R2QKB1R w - - 0 1","rn2kbnr/pp3ppp/2p1p3/q7/3P1Bb1/2N2N1P/PPP2PP1/R2QKB1R b - - 0 1","rn2kbnr/pp3ppp/2p1p3/q7/3P1B2/2N2b1P/PPP2PP1/R2QKB1R w - - 0 1","rn2kbnr/pp3ppp/2p1p3/q7/3P1B2/2N2Q1P/PPP2PP1/R3KB1R b - - 0 1","rn2k1nr/pp3ppp/2p1p3/q7/1b1P1B2/2N2Q1P/PPP2PP1/R3KB1R w - - 0 1","rn2k1nr/pp3ppp/2p1p3/q7/1b1P1B2/2N2Q1P/PPP1BPP1/R3K2R b - - 0 1","r3k1nr/pp1n1ppp/2p1p3/q7/1b1P1B2/2N2Q1P/PPP1BPP1/R3K2R w - - 0 1","r3k1nr/pp1n1ppp/2p1p3/q7/1b1P1B2/P1N2Q1P/1PP1BPP1/R3K2R b - - 0 1","2kr2nr/pp1n1ppp/2p1p3/q7/1b1P1B2/P1N2Q1P/1PP1BPP1/R3K2R w - - 0 1","2kr2nr/pp1n1ppp/2p1p3/q7/1P1P1B2/2N2Q1P/1PP1BPP1/R3K2R b - - 0 1","2kr2nr/pp1n1ppp/2p1p3/8/1P1P1B2/2N2Q1P/1PP1BPP1/q3K2R w - - 0 1","2kr2nr/pp1n1ppp/2p1p3/8/1P1P1B2/2N2Q1P/1PPKBPP1/q6R b - - 0 1","2kr2nr/pp1n1ppp/2p1p3/8/1P1P1B2/2N2Q1P/1PPKBPP1/7q w - - 0 1","2kr2nr/pp1n1ppp/2Q1p3/8/1P1P1B2/2N4P/1PPKBPP1/7q b - - 0 1","2kr2nr/p2n1ppp/2p1p3/8/1P1P1B2/2N4P/1PPKBPP1/7q w - - 0 1","2kr2nr/p2n1ppp/B1p1p3/8/1P1P1B2/2N4P/1PPK1PP1/7q b - - 0 1"],"resultado":"xeque-mate","tabuleiro_final":"8 . . k r . . n r\n7 p . . n . p p p\n6 B . p . p . . .\n5 . . . . . . . .\n4 . P . P . B . .\n3 . . N . . . . P\n2 . P P K . P P .\n1 . . . . . . . q\n  a b c d e f g h"},"origem":"wiki","palavras_chave":[],"sinonimos":[],"tipo":"conceito","titulo":"A Imortal Peruana (Canal, 1934)"}
\section{A Imortal Peruana (Canal, 1934)}
Canal sacrifica torre, dama e a outra torre numa simultânea e conclui com Ba6\# — o bispo dá mate sozinho no rei preto engaiolado. A 'Imortal Peruana'. Jogadores: Esteban Canal $\times$ N.N. (1934). A órbita percorre 28 posições e termina em xeque-mate após 27 lances.
\prov{relações: eh\_um orbita\_de\_partida; usa abertura\_xadrez; relaciona xeque\_mate; usa sacrificio\_xadrez; relaciona canal; usa lyapunov\_da\_orbita}
\prov{proveniência: wiki \textperiodcentered{} partidas reais \textperiodcentered{} fato \textperiodcentered{} confiança 1.0 \textperiodcentered{} domínio xadrez \textperiodcentered{} história 21 versões no jornal}

%CRISTAL {"arestas":[{"alvo":"orbita_de_partida","confianca":1.0,"epistemico":"fato","origem":"wiki","rel":"eh_um"},{"alvo":"abertura_xadrez","confianca":1.0,"epistemico":"fato","origem":"wiki","rel":"usa"},{"alvo":"xeque_mate","confianca":1.0,"epistemico":"fato","origem":"wiki","rel":"relaciona"},{"alvo":"sacrificio_xadrez","confianca":1.0,"epistemico":"fato","origem":"wiki","rel":"usa"},{"alvo":"reti","confianca":1.0,"epistemico":"fato","origem":"wiki","rel":"relaciona"},{"alvo":"lyapunov_da_orbita","confianca":1.0,"epistemico":"fato","origem":"wiki","rel":"usa"}],"confianca":1.0,"contraexemplos":[],"descricao":"Uma das miniaturas mais belas: o sacrifício da dama (Dd8+!!) seguido do XEQUE DUPLO (Bg5+ descobre a torre) e mate em Bd8#. Onze lances de pura geometria. Jogadores: Richard Réti × Savielly Tartakower (1910). A órbita percorre 22 posições e termina em xeque-mate após 21 lances.","epistemico":"fato","exemplos":[],"id":"partida_reti","memoria":"semantica","meta":{"abertura":"Defesa Caro-Kann","ano":1910,"classe_orbita":"borda","dominio":"xadrez","fen_final":"rnbB1b1r/ppk2ppp/2p5/4q3/4n3/8/PPP2PPP/2KR1BNR b - - 0 1","fonte":"partidas reais","jogadores":"Richard Réti × Savielly Tartakower","lances_san":["e4","c6","d4","d5","Nc3","dxe4","Nxe4","Nf6","Qd3","e5","dxe5","Qa5+","Bd2","Qxe5","O-O-O","Nxe4","Qd8+","Kxd8","Bg5+","Kc7","Bd8#"],"lyapunov":0.4492,"mineral_H":0.223,"mineral_classe":"borda","mineral_theta_g":25.9,"n_lances":21,"orbita_fen":["rnbqkbnr/pppppppp/8/8/8/8/PPPPPPPP/RNBQKBNR w - - 0 1","rnbqkbnr/pppppppp/8/8/4P3/8/PPPP1PPP/RNBQKBNR b - e3 0 1","rnbqkbnr/pp1ppppp/2p5/8/4P3/8/PPPP1PPP/RNBQKBNR w - - 0 1","rnbqkbnr/pp1ppppp/2p5/8/3PP3/8/PPP2PPP/RNBQKBNR b - d3 0 1","rnbqkbnr/pp2pppp/2p5/3p4/3PP3/8/PPP2PPP/RNBQKBNR w - d6 0 1","rnbqkbnr/pp2pppp/2p5/3p4/3PP3/2N5/PPP2PPP/R1BQKBNR b - - 0 1","rnbqkbnr/pp2pppp/2p5/8/3Pp3/2N5/PPP2PPP/R1BQKBNR w - - 0 1","rnbqkbnr/pp2pppp/2p5/8/3PN3/8/PPP2PPP/R1BQKBNR b - - 0 1","rnbqkb1r/pp2pppp/2p2n2/8/3PN3/8/PPP2PPP/R1BQKBNR w - - 0 1","rnbqkb1r/pp2pppp/2p2n2/8/3PN3/3Q4/PPP2PPP/R1B1KBNR b - - 0 1","rnbqkb1r/pp3ppp/2p2n2/4p3/3PN3/3Q4/PPP2PPP/R1B1KBNR w - e6 0 1","rnbqkb1r/pp3ppp/2p2n2/4P3/4N3/3Q4/PPP2PPP/R1B1KBNR b - - 0 1","rnb1kb1r/pp3ppp/2p2n2/q3P3/4N3/3Q4/PPP2PPP/R1B1KBNR w - - 0 1","rnb1kb1r/pp3ppp/2p2n2/q3P3/4N3/3Q4/PPPB1PPP/R3KBNR b - - 0 1","rnb1kb1r/pp3ppp/2p2n2/4q3/4N3/3Q4/PPPB1PPP/R3KBNR w - - 0 1","rnb1kb1r/pp3ppp/2p2n2/4q3/4N3/3Q4/PPPB1PPP/2KR1BNR b - - 0 1","rnb1kb1r/pp3ppp/2p5/4q3/4n3/3Q4/PPPB1PPP/2KR1BNR w - - 0 1","rnbQkb1r/pp3ppp/2p5/4q3/4n3/8/PPPB1PPP/2KR1BNR b - - 0 1","rnbk1b1r/pp3ppp/2p5/4q3/4n3/8/PPPB1PPP/2KR1BNR w - - 0 1","rnbk1b1r/pp3ppp/2p5/4q1B1/4n3/8/PPP2PPP/2KR1BNR b - - 0 1","rnb2b1r/ppk2ppp/2p5/4q1B1/4n3/8/PPP2PPP/2KR1BNR w - - 0 1","rnbB1b1r/ppk2ppp/2p5/4q3/4n3/8/PPP2PPP/2KR1BNR b - - 0 1"],"resultado":"xeque-mate","tabuleiro_final":"8 r n b B . b . r\n7 p p k . . p p p\n6 . . p . . . . .\n5 . . . . q . . .\n4 . . . . n . . .\n3 . . . . . . . .\n2 P P P . . P P P\n1 . . K R . B N R\n  a b c d e f g h"},"origem":"wiki","palavras_chave":[],"sinonimos":[],"tipo":"conceito","titulo":"Réti × Tartakower (Viena, 1910)"}
\section{Réti $\times$ Tartakower (Viena, 1910)}
Uma das miniaturas mais belas: o sacrifício da dama (Dd8+!!) seguido do XEQUE DUPLO (Bg5+ descobre a torre) e mate em Bd8\#. Onze lances de pura geometria. Jogadores: Richard Réti $\times$ Savielly Tartakower (1910). A órbita percorre 22 posições e termina em xeque-mate após 21 lances.
\prov{relações: eh\_um orbita\_de\_partida; usa abertura\_xadrez; relaciona xeque\_mate; usa sacrificio\_xadrez; relaciona reti; usa lyapunov\_da\_orbita}
\prov{proveniência: wiki \textperiodcentered{} partidas reais \textperiodcentered{} fato \textperiodcentered{} confiança 1.0 \textperiodcentered{} domínio xadrez \textperiodcentered{} história 41 versões no jornal}

%CRISTAL {"arestas":[{"alvo":"partida_imortal","confianca":1.0,"epistemico":"fato","origem":"wiki","rel":"semelhante"}],"confianca":1.0,"contraexemplos":[],"descricao":"Outra joia romântica de Anderssen: uma combinação de mate deslumbrante e correta. Batizada por Steinitz como a folha 'sempre-verde' da coroa de Anderssen.","epistemico":"fato","exemplos":[],"id":"partida_sempre_viva","memoria":"semantica","meta":{"autor":"Anderssen","dominio":"xadrez","fonte":"história do xadrez"},"origem":"wiki","palavras_chave":[],"sinonimos":[],"tipo":"conceito","titulo":"A Sempre-Viva (Evergreen, Anderssen–Dufresne, 1852)"}
\section{A Sempre-Viva (Evergreen, Anderssen–Dufresne, 1852)}
Outra joia romântica de Anderssen: uma combinação de mate deslumbrante e correta. Batizada por Steinitz como a folha 'sempre-verde' da coroa de Anderssen.
\prov{relações: semelhante partida\_imortal}
\prov{proveniência: wiki \textperiodcentered{} história do xadrez \textperiodcentered{} fato \textperiodcentered{} confiança 1.0 \textperiodcentered{} domínio xadrez \textperiodcentered{} história 2 versões no jornal}

%CRISTAL {"arestas":[{"alvo":"orbita_de_partida","confianca":1.0,"epistemico":"fato","origem":"wiki","rel":"eh_um"},{"alvo":"auto_jogo","confianca":1.0,"epistemico":"fato","origem":"wiki","rel":"usa"},{"alvo":"minimax_e_colapso","confianca":1.0,"epistemico":"fato","origem":"wiki","rel":"usa"},{"alvo":"dsl_bai_xadrez","confianca":1.0,"epistemico":"fato","origem":"wiki","rel":"relaciona"},{"alvo":"xadrez_bai","confianca":1.0,"epistemico":"fato","origem":"wiki","rel":"relaciona"},{"alvo":"lyapunov_da_orbita","confianca":1.0,"epistemico":"fato","origem":"wiki","rel":"usa"}],"confianca":1.0,"contraexemplos":[],"descricao":"A Tiffany contra si mesma: o colapso Ψ (minimax) ao orçamento δ=4 jogando os dois lados. Ao lado das históricas, esta é a órbita de uma MÁQUINA que decide como o BAI — cada lance um colapso podando ramos por deny-overrides. Ganhou uma peça e chegou a posição vencedora. A órbita percorre 41 posições e, após 40 lances a δ=4, para em em jogo. Lyapunov λ=0.357 (λ/logφ=0.74, classe borda).","epistemico":"fato","exemplos":[],"id":"partida_tiffany_d4","memoria":"semantica","meta":{"abertura":"Abertura do Cavalo-dama","ano":2026,"classe_orbita":"borda","delta":4,"dominio":"xadrez","fen_final":"1r1k1b1r/p1pq2pp/Q7/2ppP3/8/P7/1PP2PPP/R1B1K2R w - - 0 1","fonte":"auto-jogo tiffany","jogadores":"Tiffany × Tiffany (o colapso jogando os dois lados)","lances_coord":["b1c3","g8f6","e2e3","d7d6","g1f3","e7e6","d2d4","c8d7","f1b5","d7c6","a2a3","b8d7","b5c6","b7c6","c3e2","f6d5","e2f4","d5f4","e3f4","c6c5","d4d5","d7f6","d5e6","f7e6","d1e2","e8d7","e2d3","f6g4","f3e5","g4e5","f4e5","d6d5","d3b5","d7c8","b5c6","a8b8","c6e6","d8d7","e6a6","c8d8"],"lyapunov":0.357,"n_lances":40,"orbita_fen":["rnbqkbnr/pppppppp/8/8/8/8/PPPPPPPP/RNBQKBNR w - - 0 1","rnbqkbnr/pppppppp/8/8/8/2N5/PPPPPPPP/R1BQKBNR b - - 0 1","rnbqkb1r/pppppppp/5n2/8/8/2N5/PPPPPPPP/R1BQKBNR w - - 0 1","rnbqkb1r/pppppppp/5n2/8/8/2N1P3/PPPP1PPP/R1BQKBNR b - - 0 1","rnbqkb1r/ppp1pppp/3p1n2/8/8/2N1P3/PPPP1PPP/R1BQKBNR w - - 0 1","rnbqkb1r/ppp1pppp/3p1n2/8/8/2N1PN2/PPPP1PPP/R1BQKB1R b - - 0 1","rnbqkb1r/ppp2ppp/3ppn2/8/8/2N1PN2/PPPP1PPP/R1BQKB1R w - - 0 1","rnbqkb1r/ppp2ppp/3ppn2/8/3P4/2N1PN2/PPP2PPP/R1BQKB1R b - - 0 1","rn1qkb1r/pppb1ppp/3ppn2/8/3P4/2N1PN2/PPP2PPP/R1BQKB1R w - - 0 1","rn1qkb1r/pppb1ppp/3ppn2/1B6/3P4/2N1PN2/PPP2PPP/R1BQK2R b - - 0 1","rn1qkb1r/ppp2ppp/2bppn2/1B6/3P4/2N1PN2/PPP2PPP/R1BQK2R w - - 0 1","rn1qkb1r/ppp2ppp/2bppn2/1B6/3P4/P1N1PN2/1PP2PPP/R1BQK2R b - - 0 1","r2qkb1r/pppn1ppp/2bppn2/1B6/3P4/P1N1PN2/1PP2PPP/R1BQK2R w - - 0 1","r2qkb1r/pppn1ppp/2Bppn2/8/3P4/P1N1PN2/1PP2PPP/R1BQK2R b - - 0 1","r2qkb1r/p1pn1ppp/2pppn2/8/3P4/P1N1PN2/1PP2PPP/R1BQK2R w - - 0 1","r2qkb1r/p1pn1ppp/2pppn2/8/3P4/P3PN2/1PP1NPPP/R1BQK2R b - - 0 1","r2qkb1r/p1pn1ppp/2ppp3/3n4/3P4/P3PN2/1PP1NPPP/R1BQK2R w - - 0 1","r2qkb1r/p1pn1ppp/2ppp3/3n4/3P1N2/P3PN2/1PP2PPP/R1BQK2R b - - 0 1","r2qkb1r/p1pn1ppp/2ppp3/8/3P1n2/P3PN2/1PP2PPP/R1BQK2R w - - 0 1","r2qkb1r/p1pn1ppp/2ppp3/8/3P1P2/P4N2/1PP2PPP/R1BQK2R b - - 0 1","r2qkb1r/p1pn1ppp/3pp3/2p5/3P1P2/P4N2/1PP2PPP/R1BQK2R w - - 0 1","r2qkb1r/p1pn1ppp/3pp3/2pP4/5P2/P4N2/1PP2PPP/R1BQK2R b - - 0 1","r2qkb1r/p1p2ppp/3ppn2/2pP4/5P2/P4N2/1PP2PPP/R1BQK2R w - - 0 1","r2qkb1r/p1p2ppp/3pPn2/2p5/5P2/P4N2/1PP2PPP/R1BQK2R b - - 0 1","r2qkb1r/p1p3pp/3ppn2/2p5/5P2/P4N2/1PP2PPP/R1BQK2R w - - 0 1","r2qkb1r/p1p3pp/3ppn2/2p5/5P2/P4N2/1PP1QPPP/R1B1K2R b - - 0 1","r2q1b1r/p1pk2pp/3ppn2/2p5/5P2/P4N2/1PP1QPPP/R1B1K2R w - - 0 1","r2q1b1r/p1pk2pp/3ppn2/2p5/5P2/P2Q1N2/1PP2PPP/R1B1K2R b - - 0 1","r2q1b1r/p1pk2pp/3pp3/2p5/5Pn1/P2Q1N2/1PP2PPP/R1B1K2R w - - 0 1","r2q1b1r/p1pk2pp/3pp3/2p1N3/5Pn1/P2Q4/1PP2PPP/R1B1K2R b - - 0 1","r2q1b1r/p1pk2pp/3pp3/2p1n3/5P2/P2Q4/1PP2PPP/R1B1K2R w - - 0 1","r2q1b1r/p1pk2pp/3pp3/2p1P3/8/P2Q4/1PP2PPP/R1B1K2R b - - 0 1","r2q1b1r/p1pk2pp/4p3/2ppP3/8/P2Q4/1PP2PPP/R1B1K2R w - - 0 1","r2q1b1r/p1pk2pp/4p3/1QppP3/8/P7/1PP2PPP/R1B1K2R b - - 0 1","r1kq1b1r/p1p3pp/4p3/1QppP3/8/P7/1PP2PPP/R1B1K2R w - - 0 1","r1kq1b1r/p1p3pp/2Q1p3/2ppP3/8/P7/1PP2PPP/R1B1K2R b - - 0 1","1rkq1b1r/p1p3pp/2Q1p3/2ppP3/8/P7/1PP2PPP/R1B1K2R w - - 0 1","1rkq1b1r/p1p3pp/4Q3/2ppP3/8/P7/1PP2PPP/R1B1K2R b - - 0 1","1rk2b1r/p1pq2pp/4Q3/2ppP3/8/P7/1PP2PPP/R1B1K2R w - - 0 1","1rk2b1r/p1pq2pp/Q7/2ppP3/8/P7/1PP2PPP/R1B1K2R b - - 0 1","1r1k1b1r/p1pq2pp/Q7/2ppP3/8/P7/1PP2PPP/R1B1K2R w - - 0 1"],"resultado":"em jogo","tabuleiro_final":"8 . r . k . b . r\n7 p . p q . . p p\n6 Q . . . . . . .\n5 . . p p P . . .\n4 . . . . . . . .\n3 P . . . . . . .\n2 . P P . . P P P\n1 R . B . K . . R\n  a b c d e f g h"},"origem":"wiki","palavras_chave":[],"sinonimos":[],"tipo":"conceito","titulo":"A Órbita da Tiffany (auto-jogo, δ=4)"}
\section{A Órbita da Tiffany (auto-jogo, \ensuremath{\delta}=4)}
A Tiffany contra si mesma: o colapso \ensuremath{\Psi} (minimax) ao orçamento \ensuremath{\delta}=4 jogando os dois lados. Ao lado das históricas, esta é a órbita de uma MÁQUINA que decide como o BAI — cada lance um colapso podando ramos por deny-overrides. Ganhou uma peça e chegou a posição vencedora. A órbita percorre 41 posições e, após 40 lances a \ensuremath{\delta}=4, para em em jogo. Lyapunov \ensuremath{\lambda}=0.357 (\ensuremath{\lambda}/log\ensuremath{\varphi}=0.74, classe borda).
\prov{relações: eh\_um orbita\_de\_partida; usa auto\_jogo; usa minimax\_e\_colapso; relaciona dsl\_bai\_xadrez; relaciona xadrez\_bai; usa lyapunov\_da\_orbita}
\prov{proveniência: wiki \textperiodcentered{} auto-jogo tiffany \textperiodcentered{} fato \textperiodcentered{} confiança 1.0 \textperiodcentered{} domínio xadrez \textperiodcentered{} história 41 versões no jornal}

%CRISTAL {"arestas":[{"alvo":"orbita_de_partida","confianca":1.0,"epistemico":"fato","origem":"wiki","rel":"eh_um"},{"alvo":"auto_jogo","confianca":1.0,"epistemico":"fato","origem":"wiki","rel":"usa"},{"alvo":"minimax_e_colapso","confianca":1.0,"epistemico":"fato","origem":"wiki","rel":"usa"},{"alvo":"dsl_bai_xadrez","confianca":1.0,"epistemico":"fato","origem":"wiki","rel":"relaciona"},{"alvo":"xadrez_bai","confianca":1.0,"epistemico":"fato","origem":"wiki","rel":"relaciona"},{"alvo":"lyapunov_da_orbita","confianca":1.0,"epistemico":"fato","origem":"wiki","rel":"usa"}],"confianca":1.0,"contraexemplos":[],"descricao":"A Tiffany a δ=5 — busca mais funda. Aqui ela ABRE com e4 (o lance principista, ao contrário do Cc3 raso dos δ menores): a profundidade escolhe melhor. Jogou 35 min uma partida afiadíssima (órbita borda, não a técnica-cristal da δ=4) — busca mais funda deu abertura melhor, não órbita mais previsível: escolheu a linha mais combativa. A órbita percorre 31 posições e, após 30 lances a δ=5, para em em jogo. Lyapunov λ=0.320 (λ/logφ=0.67, classe borda).","epistemico":"fato","exemplos":[],"id":"partida_tiffany_d5","memoria":"semantica","meta":{"abertura":"Abertura do Peão-rei (e4)","ano":2026,"classe_orbita":"borda","delta":5,"dominio":"xadrez","fen_final":"r1bB1knr/pp1p1ppp/4p3/1B2P3/1q6/2N2Q2/2PK1PPP/R6R w - - 0 1","fonte":"auto-jogo tiffany","jogadores":"Tiffany × Tiffany (o colapso jogando os dois lados)","lances_coord":["e2e4","b8c6","b1c3","e7e6","g1f3","d8f6","f1b5","c6e5","f3e5","f6e5","d2d4","e5d6","e4e5","d6b4","a2a3","b4a5","c1g5","c7c5","d4c5","f8c5","b2b4","c5b4","a3b4","a5b4","d1f3","e8f8","g5d8","b4b2","e1d2","b2b4"],"lyapunov":0.3203,"n_lances":30,"orbita_fen":["rnbqkbnr/pppppppp/8/8/8/8/PPPPPPPP/RNBQKBNR w - - 0 1","rnbqkbnr/pppppppp/8/8/4P3/8/PPPP1PPP/RNBQKBNR b - - 0 1","r1bqkbnr/pppppppp/2n5/8/4P3/8/PPPP1PPP/RNBQKBNR w - - 0 1","r1bqkbnr/pppppppp/2n5/8/4P3/2N5/PPPP1PPP/R1BQKBNR b - - 0 1","r1bqkbnr/pppp1ppp/2n1p3/8/4P3/2N5/PPPP1PPP/R1BQKBNR w - - 0 1","r1bqkbnr/pppp1ppp/2n1p3/8/4P3/2N2N2/PPPP1PPP/R1BQKB1R b - - 0 1","r1b1kbnr/pppp1ppp/2n1pq2/8/4P3/2N2N2/PPPP1PPP/R1BQKB1R w - - 0 1","r1b1kbnr/pppp1ppp/2n1pq2/1B6/4P3/2N2N2/PPPP1PPP/R1BQK2R b - - 0 1","r1b1kbnr/pppp1ppp/4pq2/1B2n3/4P3/2N2N2/PPPP1PPP/R1BQK2R w - - 0 1","r1b1kbnr/pppp1ppp/4pq2/1B2N3/4P3/2N5/PPPP1PPP/R1BQK2R b - - 0 1","r1b1kbnr/pppp1ppp/4p3/1B2q3/4P3/2N5/PPPP1PPP/R1BQK2R w - - 0 1","r1b1kbnr/pppp1ppp/4p3/1B2q3/3PP3/2N5/PPP2PPP/R1BQK2R b - - 0 1","r1b1kbnr/pppp1ppp/3qp3/1B6/3PP3/2N5/PPP2PPP/R1BQK2R w - - 0 1","r1b1kbnr/pppp1ppp/3qp3/1B2P3/3P4/2N5/PPP2PPP/R1BQK2R b - - 0 1","r1b1kbnr/pppp1ppp/4p3/1B2P3/1q1P4/2N5/PPP2PPP/R1BQK2R w - - 0 1","r1b1kbnr/pppp1ppp/4p3/1B2P3/1q1P4/P1N5/1PP2PPP/R1BQK2R b - - 0 1","r1b1kbnr/pppp1ppp/4p3/qB2P3/3P4/P1N5/1PP2PPP/R1BQK2R w - - 0 1","r1b1kbnr/pppp1ppp/4p3/qB2P1B1/3P4/P1N5/1PP2PPP/R2QK2R b - - 0 1","r1b1kbnr/pp1p1ppp/4p3/qBp1P1B1/3P4/P1N5/1PP2PPP/R2QK2R w - - 0 1","r1b1kbnr/pp1p1ppp/4p3/qBP1P1B1/8/P1N5/1PP2PPP/R2QK2R b - - 0 1","r1b1k1nr/pp1p1ppp/4p3/qBb1P1B1/8/P1N5/1PP2PPP/R2QK2R w - - 0 1","r1b1k1nr/pp1p1ppp/4p3/qBb1P1B1/1P6/P1N5/2P2PPP/R2QK2R b - - 0 1","r1b1k1nr/pp1p1ppp/4p3/qB2P1B1/1b6/P1N5/2P2PPP/R2QK2R w - - 0 1","r1b1k1nr/pp1p1ppp/4p3/qB2P1B1/1P6/2N5/2P2PPP/R2QK2R b - - 0 1","r1b1k1nr/pp1p1ppp/4p3/1B2P1B1/1q6/2N5/2P2PPP/R2QK2R w - - 0 1","r1b1k1nr/pp1p1ppp/4p3/1B2P1B1/1q6/2N2Q2/2P2PPP/R3K2R b - - 0 1","r1b2knr/pp1p1ppp/4p3/1B2P1B1/1q6/2N2Q2/2P2PPP/R3K2R w - - 0 1","r1bB1knr/pp1p1ppp/4p3/1B2P3/1q6/2N2Q2/2P2PPP/R3K2R b - - 0 1","r1bB1knr/pp1p1ppp/4p3/1B2P3/8/2N2Q2/1qP2PPP/R3K2R w - - 0 1","r1bB1knr/pp1p1ppp/4p3/1B2P3/8/2N2Q2/1qPK1PPP/R6R b - - 0 1","r1bB1knr/pp1p1ppp/4p3/1B2P3/1q6/2N2Q2/2PK1PPP/R6R w - - 0 1"],"resultado":"em jogo","tabuleiro_final":"8 r . b B . k n r\n7 p p . p . p p p\n6 . . . . p . . .\n5 . B . . P . . .\n4 . q . . . . . .\n3 . . N . . Q . .\n2 . . P K . P P P\n1 R . . . . . . R\n  a b c d e f g h"},"origem":"wiki","palavras_chave":[],"sinonimos":[],"tipo":"conceito","titulo":"A Órbita da Tiffany (auto-jogo, δ=5)"}
\section{A Órbita da Tiffany (auto-jogo, \ensuremath{\delta}=5)}
A Tiffany a \ensuremath{\delta}=5 — busca mais funda. Aqui ela ABRE com e4 (o lance principista, ao contrário do Cc3 raso dos \ensuremath{\delta} menores): a profundidade escolhe melhor. Jogou 35 min uma partida afiadíssima (órbita borda, não a técnica-cristal da \ensuremath{\delta}=4) — busca mais funda deu abertura melhor, não órbita mais previsível: escolheu a linha mais combativa. A órbita percorre 31 posições e, após 30 lances a \ensuremath{\delta}=5, para em em jogo. Lyapunov \ensuremath{\lambda}=0.320 (\ensuremath{\lambda}/log\ensuremath{\varphi}=0.67, classe borda).
\prov{relações: eh\_um orbita\_de\_partida; usa auto\_jogo; usa minimax\_e\_colapso; relaciona dsl\_bai\_xadrez; relaciona xadrez\_bai; usa lyapunov\_da\_orbita}
\prov{proveniência: wiki \textperiodcentered{} auto-jogo tiffany \textperiodcentered{} fato \textperiodcentered{} confiança 1.0 \textperiodcentered{} domínio xadrez \textperiodcentered{} história 28 versões no jornal}

%CRISTAL {"arestas":[{"alvo":"orbita_de_partida","confianca":1.0,"epistemico":"fato","origem":"wiki","rel":"eh_um"},{"alvo":"auto_jogo","confianca":1.0,"epistemico":"fato","origem":"wiki","rel":"usa"},{"alvo":"minimax_e_colapso","confianca":1.0,"epistemico":"fato","origem":"wiki","rel":"usa"},{"alvo":"dsl_bai_xadrez","confianca":1.0,"epistemico":"fato","origem":"wiki","rel":"relaciona"},{"alvo":"xadrez_bai","confianca":1.0,"epistemico":"fato","origem":"wiki","rel":"relaciona"},{"alvo":"lyapunov_da_orbita","confianca":1.0,"epistemico":"fato","origem":"wiki","rel":"usa"}],"confianca":1.0,"contraexemplos":[],"descricao":"A Tiffany a δ=6 — a busca mais funda (2,7 h de cálculo). Em 20 lances ela CONVERGE a uma REPETIÇÃO (Cd5-f6 / Ce2-c3 em loop): o empate. A órbita mais CRISTALINA de todas (quase-periódica, entropia baixa) — busca mais funda tende ao valor teórico do xadrez, o empate. O oposto do romantismo caótico dos sacrifícios. A órbita percorre 21 posições e, após 20 lances a δ=6, para em em jogo. Lyapunov λ=0.290 (λ/logφ=0.60, classe borda).","epistemico":"fato","exemplos":[],"id":"partida_tiffany_d6","memoria":"semantica","meta":{"abertura":"Abertura do Peão-rei (e3)","ano":2026,"classe_orbita":"borda","delta":6,"dominio":"xadrez","fen_final":"r1b1kb1r/1pp1qppp/p1p1pn2/8/3P4/P3PN2/1PP1NPPP/R1BQK2R w - - 0 1","fonte":"auto-jogo tiffany","jogadores":"Tiffany × Tiffany (o colapso jogando os dois lados)","lances_coord":["e2e3","b8c6","g1f3","g8f6","b1c3","e7e6","f1b5","a7a6","b5c6","d7c6","a2a3","d8e7","d2d4","f6d5","c3e2","d5f6","e2c3","f6d5","c3e2","d5f6"],"lyapunov":0.2895,"n_lances":20,"orbita_fen":["rnbqkbnr/pppppppp/8/8/8/8/PPPPPPPP/RNBQKBNR w - - 0 1","rnbqkbnr/pppppppp/8/8/8/4P3/PPPP1PPP/RNBQKBNR b - - 0 1","r1bqkbnr/pppppppp/2n5/8/8/4P3/PPPP1PPP/RNBQKBNR w - - 0 1","r1bqkbnr/pppppppp/2n5/8/8/4PN2/PPPP1PPP/RNBQKB1R b - - 0 1","r1bqkb1r/pppppppp/2n2n2/8/8/4PN2/PPPP1PPP/RNBQKB1R w - - 0 1","r1bqkb1r/pppppppp/2n2n2/8/8/2N1PN2/PPPP1PPP/R1BQKB1R b - - 0 1","r1bqkb1r/pppp1ppp/2n1pn2/8/8/2N1PN2/PPPP1PPP/R1BQKB1R w - - 0 1","r1bqkb1r/pppp1ppp/2n1pn2/1B6/8/2N1PN2/PPPP1PPP/R1BQK2R b - - 0 1","r1bqkb1r/1ppp1ppp/p1n1pn2/1B6/8/2N1PN2/PPPP1PPP/R1BQK2R w - - 0 1","r1bqkb1r/1ppp1ppp/p1B1pn2/8/8/2N1PN2/PPPP1PPP/R1BQK2R b - - 0 1","r1bqkb1r/1pp2ppp/p1p1pn2/8/8/2N1PN2/PPPP1PPP/R1BQK2R w - - 0 1","r1bqkb1r/1pp2ppp/p1p1pn2/8/8/P1N1PN2/1PPP1PPP/R1BQK2R b - - 0 1","r1b1kb1r/1pp1qppp/p1p1pn2/8/8/P1N1PN2/1PPP1PPP/R1BQK2R w - - 0 1","r1b1kb1r/1pp1qppp/p1p1pn2/8/3P4/P1N1PN2/1PP2PPP/R1BQK2R b - - 0 1","r1b1kb1r/1pp1qppp/p1p1p3/3n4/3P4/P1N1PN2/1PP2PPP/R1BQK2R w - - 0 1","r1b1kb1r/1pp1qppp/p1p1p3/3n4/3P4/P3PN2/1PP1NPPP/R1BQK2R b - - 0 1","r1b1kb1r/1pp1qppp/p1p1pn2/8/3P4/P3PN2/1PP1NPPP/R1BQK2R w - - 0 1","r1b1kb1r/1pp1qppp/p1p1pn2/8/3P4/P1N1PN2/1PP2PPP/R1BQK2R b - - 0 1","r1b1kb1r/1pp1qppp/p1p1p3/3n4/3P4/P1N1PN2/1PP2PPP/R1BQK2R w - - 0 1","r1b1kb1r/1pp1qppp/p1p1p3/3n4/3P4/P3PN2/1PP1NPPP/R1BQK2R b - - 0 1","r1b1kb1r/1pp1qppp/p1p1pn2/8/3P4/P3PN2/1PP1NPPP/R1BQK2R w - - 0 1"],"resultado":"em jogo","tabuleiro_final":"8 r . b . k b . r\n7 . p p . q p p p\n6 p . p . p n . .\n5 . . . . . . . .\n4 . . . P . . . .\n3 P . . . P N . .\n2 . P P . N P P P\n1 R . B Q K . . R\n  a b c d e f g h"},"origem":"wiki","palavras_chave":[],"sinonimos":[],"tipo":"conceito","titulo":"A Órbita da Tiffany (auto-jogo, δ=6)"}
\section{A Órbita da Tiffany (auto-jogo, \ensuremath{\delta}=6)}
A Tiffany a \ensuremath{\delta}=6 — a busca mais funda (2,7 h de cálculo). Em 20 lances ela CONVERGE a uma REPETIÇÃO (Cd5-f6 / Ce2-c3 em loop): o empate. A órbita mais CRISTALINA de todas (quase-periódica, entropia baixa) — busca mais funda tende ao valor teórico do xadrez, o empate. O oposto do romantismo caótico dos sacrifícios. A órbita percorre 21 posições e, após 20 lances a \ensuremath{\delta}=6, para em em jogo. Lyapunov \ensuremath{\lambda}=0.290 (\ensuremath{\lambda}/log\ensuremath{\varphi}=0.60, classe borda).
\prov{relações: eh\_um orbita\_de\_partida; usa auto\_jogo; usa minimax\_e\_colapso; relaciona dsl\_bai\_xadrez; relaciona xadrez\_bai; usa lyapunov\_da\_orbita}
\prov{proveniência: wiki \textperiodcentered{} auto-jogo tiffany \textperiodcentered{} fato \textperiodcentered{} confiança 1.0 \textperiodcentered{} domínio xadrez \textperiodcentered{} história 7 versões no jornal}

%CRISTAL {"fusao":[{"arestas":[{"alvo":"xadrez","confianca":1.0,"epistemico":"fato","origem":"wiki","rel":"parte_de"}],"confianca":1.0,"contraexemplos":[],"descricao":"O gênio americano (1837–1884): entendeu desenvolvimento e iniciativa décadas antes de todos. Dominou a era romântica e abandonou o jogo aos 21. A ponte entre o ataque intuitivo e o jogo de princípios.","epistemico":"fato","exemplos":[],"id":"paul_morphy","memoria":"semantica","meta":{"autor":"Paul Morphy","dominio":"xadrez","fonte":"história do xadrez"},"origem":"wiki","palavras_chave":[],"sinonimos":[],"tipo":"conceito","titulo":"Paul Morphy"},{"arestas":[{"alvo":"xadrez","confianca":1.0,"epistemico":"fato","origem":"wiki","rel":"relaciona"}],"confianca":1.0,"contraexemplos":[],"descricao":"O maior enxadrista do século XIX (1837–1884), gênio do desenvolvimento rápido e da iniciativa. Antecipou a estratégia moderna décadas antes de Steinitz formalizá-la.","epistemico":"fato","exemplos":[],"id":"morphy","memoria":"semantica","meta":{"dominio":"xadrez","fonte":"jogadores"},"origem":"wiki","palavras_chave":[],"sinonimos":[],"tipo":"conceito","titulo":"Paul Morphy"}],"id":"paul_morphy","tipo":"conceito"}
\section{Paul Morphy}
O gênio americano (1837–1884): entendeu desenvolvimento e iniciativa décadas antes de todos. Dominou a era romântica e abandonou o jogo aos 21. A ponte entre o ataque intuitivo e o jogo de princípios.
\prov{relações: parte\_de xadrez}
\prov{proveniência: wiki \textperiodcentered{} história do xadrez \textperiodcentered{} fato \textperiodcentered{} confiança 1.0 \textperiodcentered{} domínio xadrez \textperiodcentered{} história 2 versões no jornal}
\prov{a segunda face --- morphy:}
O maior enxadrista do século XIX (1837–1884), gênio do desenvolvimento rápido e da iniciativa. Antecipou a estratégia moderna décadas antes de Steinitz formalizá-la.
\prov{fusão de curadoria (julgada): absorve morphy --- as duas faces acima, intactas no registo}

%CRISTAL {"arestas":[{"alvo":"estrutura_de_peoes","confianca":1.0,"epistemico":"fato","origem":"wiki","rel":"parte_de"}],"confianca":1.0,"contraexemplos":[],"descricao":"Um peão sem peões inimigos à frente na sua coluna nem nas vizinhas — candidato a promover a dama. 'Deve ser empurrado' (Nimzowitsch): uma ameaça crescente que amarra o defensor.","epistemico":"inferencia","exemplos":[],"id":"peao_passado","memoria":"semantica","meta":{"dominio":"xadrez","fonte":"xadrez"},"origem":"wiki","palavras_chave":[],"sinonimos":[],"tipo":"conceito","titulo":"Peão Passado"}
\section{Peão Passado}
Um peão sem peões inimigos à frente na sua coluna nem nas vizinhas — candidato a promover a dama. 'Deve ser empurrado' (Nimzowitsch): uma ameaça crescente que amarra o defensor.
\prov{relações: parte\_de estrutura\_de\_peoes}
\prov{proveniência: wiki \textperiodcentered{} xadrez \textperiodcentered{} inferencia \textperiodcentered{} confiança 1.0 \textperiodcentered{} domínio xadrez \textperiodcentered{} história 2 versões no jornal}

%CRISTAL {"arestas":[{"alvo":"partida_alekhine_vasic","confianca":1.0,"epistemico":"fato","origem":"wiki","rel":"parte_de"},{"alvo":"xeque_mate","confianca":1.0,"epistemico":"fato","origem":"wiki","rel":"relaciona"}],"confianca":1.0,"contraexemplos":[],"descricao":"Xeque-mate real, tabuleiro e peças:\n8 r . b q k . . r\n7 p . p n . . p .\n6 . p . . p n B p\n5 . . . . . . . .\n4 . . . P . . . .\n3 B . P . . . . .\n2 P . P . . P P P\n1 R . . . K . N R\n  a b c d e f g h\nFEN: r1bqk2r/p1pn2p1/1p2pnBp/8/3P4/B1P5/P1P2PPP/R3K1NR b - - 0 1","epistemico":"fato","exemplos":[],"id":"pos_alekhine_vasic_final","memoria":"semantica","meta":{"dominio":"xadrez","fen":"r1bqk2r/p1pn2p1/1p2pnBp/8/3P4/B1P5/P1P2PPP/R3K1NR b - - 0 1","fonte":"posições reais","partida":"alekhine_vasic","tabuleiro":"8 r . b q k . . r\n7 p . p n . . p .\n6 . p . . p n B p\n5 . . . . . . . .\n4 . . . P . . . .\n3 B . P . . . . .\n2 P . P . . P P P\n1 R . . . K . N R\n  a b c d e f g h"},"origem":"wiki","palavras_chave":[],"sinonimos":[],"tipo":"conceito","titulo":"Posição final — Alekhine × Vasic (1931)"}
\section{Posição final — Alekhine $\times$ Vasic (1931)}
Xeque-mate real, tabuleiro e peças:
8 r . b q k . . r
7 p . p n . . p .
6 . p . . p n B p
5 . . . . . . . .
4 . . . P . . . .
3 B . P . . . . .
2 P . P . . P P P
1 R . . . K . N R
  a b c d e f g h
FEN: r1bqk2r/p1pn2p1/1p2pnBp/8/3P4/B1P5/P1P2PPP/R3K1NR b - - 0 1
\prov{relações: parte\_de partida\_alekhine\_vasic; relaciona xeque\_mate}
\prov{proveniência: wiki \textperiodcentered{} posições reais \textperiodcentered{} fato \textperiodcentered{} confiança 1.0 \textperiodcentered{} domínio xadrez \textperiodcentered{} história 6 versões no jornal}

%CRISTAL {"arestas":[{"alvo":"partida_bobo","confianca":1.0,"epistemico":"fato","origem":"wiki","rel":"parte_de"},{"alvo":"xeque_mate","confianca":1.0,"epistemico":"fato","origem":"wiki","rel":"relaciona"}],"confianca":1.0,"contraexemplos":[],"descricao":"Xeque-mate real, tabuleiro e peças:\n8 r n b . k b n r\n7 p p p p . p p p\n6 . . . . . . . .\n5 . . . . p . . .\n4 . . . . . . P q\n3 . . . . . P . .\n2 P P P P P . . P\n1 R N B Q K B N R\n  a b c d e f g h\nFEN: rnb1kbnr/pppp1ppp/8/4p3/6Pq/5P2/PPPPP2P/RNBQKBNR w - - 0 1","epistemico":"fato","exemplos":[],"id":"pos_bobo_final","memoria":"semantica","meta":{"dominio":"xadrez","fen":"rnb1kbnr/pppp1ppp/8/4p3/6Pq/5P2/PPPPP2P/RNBQKBNR w - - 0 1","fonte":"posições reais","partida":"bobo","tabuleiro":"8 r n b . k b n r\n7 p p p p . p p p\n6 . . . . . . . .\n5 . . . . p . . .\n4 . . . . . . P q\n3 . . . . . P . .\n2 P P P P P . . P\n1 R N B Q K B N R\n  a b c d e f g h"},"origem":"wiki","palavras_chave":[],"sinonimos":[],"tipo":"conceito","titulo":"Posição final — Mate do Bobo (o mais rápido)"}
\section{Posição final — Mate do Bobo (o mais rápido)}
Xeque-mate real, tabuleiro e peças:
8 r n b . k b n r
7 p p p p . p p p
6 . . . . . . . .
5 . . . . p . . .
4 . . . . . . P q
3 . . . . . P . .
2 P P P P P . . P
1 R N B Q K B N R
  a b c d e f g h
FEN: rnb1kbnr/pppp1ppp/8/4p3/6Pq/5P2/PPPPP2P/RNBQKBNR w - - 0 1
\prov{relações: parte\_de partida\_bobo; relaciona xeque\_mate}
\prov{proveniência: wiki \textperiodcentered{} posições reais \textperiodcentered{} fato \textperiodcentered{} confiança 1.0 \textperiodcentered{} domínio xadrez \textperiodcentered{} história 25 versões no jornal}

%CRISTAL {"arestas":[{"alvo":"partida_deepblue","confianca":1.0,"epistemico":"fato","origem":"wiki","rel":"parte_de"},{"alvo":"xeque_mate","confianca":1.0,"epistemico":"fato","origem":"wiki","rel":"relaciona"}],"confianca":1.0,"contraexemplos":[],"descricao":"Xeque-mate real, tabuleiro e peças:\n8 r . k . . . . r\n7 p . . n b . p .\n6 . . b . . . . p\n5 . p . n . p . .\n4 . . P P . . . .\n3 . . . Q . N B .\n2 . P . . . P P P\n1 R . . . . . K .\n  a b c d e f g h\nFEN: r1k4r/p2nb1p1/2b4p/1p1n1p2/2PP4/3Q1NB1/1P3PPP/R5K1 b - c3 0 1","epistemico":"fato","exemplos":[],"id":"pos_deepblue_final","memoria":"semantica","meta":{"dominio":"xadrez","fen":"r1k4r/p2nb1p1/2b4p/1p1n1p2/2PP4/3Q1NB1/1P3PPP/R5K1 b - c3 0 1","fonte":"posições reais","partida":"deepblue","tabuleiro":"8 r . k . . . . r\n7 p . . n b . p .\n6 . . b . . . . p\n5 . p . n . p . .\n4 . . P P . . . .\n3 . . . Q . N B .\n2 . P . . . P P P\n1 R . . . . . K .\n  a b c d e f g h"},"origem":"wiki","palavras_chave":[],"sinonimos":[],"tipo":"conceito","titulo":"Posição final — Deep Blue × Kasparov (1997, Jogo 6)"}
\section{Posição final — Deep Blue $\times$ Kasparov (1997, Jogo 6)}
Xeque-mate real, tabuleiro e peças:
8 r . k . . . . r
7 p . . n b . p .
6 . . b . . . . p
5 . p . n . p . .
4 . . P P . . . .
3 . . . Q . N B .
2 . P . . . P P P
1 R . . . . . K .
  a b c d e f g h
FEN: r1k4r/p2nb1p1/2b4p/1p1n1p2/2PP4/3Q1NB1/1P3PPP/R5K1 b - c3 0 1
\prov{relações: parte\_de partida\_deepblue; relaciona xeque\_mate}
\prov{proveniência: wiki \textperiodcentered{} posições reais \textperiodcentered{} fato \textperiodcentered{} confiança 1.0 \textperiodcentered{} domínio xadrez \textperiodcentered{} história 18 versões no jornal}

%CRISTAL {"arestas":[{"alvo":"partida_evergreen","confianca":1.0,"epistemico":"fato","origem":"wiki","rel":"parte_de"},{"alvo":"xeque_mate","confianca":1.0,"epistemico":"fato","origem":"wiki","rel":"relaciona"}],"confianca":1.0,"contraexemplos":[],"descricao":"Xeque-mate real, tabuleiro e peças:\n8 . r . . . k r .\n7 p b p B B p . p\n6 . b . . . P . .\n5 . . . . . . . .\n4 . . . . . . . .\n3 . . P . . q . .\n2 P . . . . P P P\n1 . . . R . . K .\n  a b c d e f g h\nFEN: 1r3kr1/pbpBBp1p/1b3P2/8/8/2P2q2/P4PPP/3R2K1 b - - 0 1","epistemico":"fato","exemplos":[],"id":"pos_evergreen_final","memoria":"semantica","meta":{"dominio":"xadrez","fen":"1r3kr1/pbpBBp1p/1b3P2/8/8/2P2q2/P4PPP/3R2K1 b - - 0 1","fonte":"posições reais","partida":"evergreen","tabuleiro":"8 . r . . . k r .\n7 p b p B B p . p\n6 . b . . . P . .\n5 . . . . . . . .\n4 . . . . . . . .\n3 . . P . . q . .\n2 P . . . . P P P\n1 . . . R . . K .\n  a b c d e f g h"},"origem":"wiki","palavras_chave":[],"sinonimos":[],"tipo":"conceito","titulo":"Posição final — A Sempreviva (Anderssen, 1852)"}
\section{Posição final — A Sempreviva (Anderssen, 1852)}
Xeque-mate real, tabuleiro e peças:
8 . r . . . k r .
7 p b p B B p . p
6 . b . . . P . .
5 . . . . . . . .
4 . . . . . . . .
3 . . P . . q . .
2 P . . . . P P P
1 . . . R . . K .
  a b c d e f g h
FEN: 1r3kr1/pbpBBp1p/1b3P2/8/8/2P2q2/P4PPP/3R2K1 b - - 0 1
\prov{relações: parte\_de partida\_evergreen; relaciona xeque\_mate}
\prov{proveniência: wiki \textperiodcentered{} posições reais \textperiodcentered{} fato \textperiodcentered{} confiança 1.0 \textperiodcentered{} domínio xadrez \textperiodcentered{} história 18 versões no jornal}

%CRISTAL {"arestas":[{"alvo":"partida_fischer_byrne","confianca":1.0,"epistemico":"fato","origem":"wiki","rel":"parte_de"},{"alvo":"xeque_mate","confianca":1.0,"epistemico":"fato","origem":"wiki","rel":"relaciona"}],"confianca":1.0,"contraexemplos":[],"descricao":"Xeque-mate real, tabuleiro e peças:\n8 . Q . . . . . .\n7 . . . . . p k .\n6 . . p . . . p .\n5 . p . . N . . p\n4 . b . . . . . P\n3 . b n . . . . .\n2 . . r . . . P .\n1 . . K . . . . .\n  a b c d e f g h\nFEN: 1Q6/5pk1/2p3p1/1p2N2p/1b5P/1bn5/2r3P1/2K5 w - - 0 1","epistemico":"fato","exemplos":[],"id":"pos_fischer_byrne_final","memoria":"semantica","meta":{"dominio":"xadrez","fen":"1Q6/5pk1/2p3p1/1p2N2p/1b5P/1bn5/2r3P1/2K5 w - - 0 1","fonte":"posições reais","partida":"fischer_byrne","tabuleiro":"8 . Q . . . . . .\n7 . . . . . p k .\n6 . . p . . . p .\n5 . p . . N . . p\n4 . b . . . . . P\n3 . b n . . . . .\n2 . . r . . . P .\n1 . . K . . . . .\n  a b c d e f g h"},"origem":"wiki","palavras_chave":[],"sinonimos":[],"tipo":"conceito","titulo":"Posição final — A Partida do Século (Byrne × Fischer, 1956)"}
\section{Posição final — A Partida do Século (Byrne $\times$ Fischer, 1956)}
Xeque-mate real, tabuleiro e peças:
8 . Q . . . . . .
7 . . . . . p k .
6 . . p . . . p .
5 . p . . N . . p
4 . b . . . . . P
3 . b n . . . . .
2 . . r . . . P .
1 . . K . . . . .
  a b c d e f g h
FEN: 1Q6/5pk1/2p3p1/1p2N2p/1b5P/1bn5/2r3P1/2K5 w - - 0 1
\prov{relações: parte\_de partida\_fischer\_byrne; relaciona xeque\_mate}
\prov{proveniência: wiki \textperiodcentered{} posições reais \textperiodcentered{} fato \textperiodcentered{} confiança 1.0 \textperiodcentered{} domínio xadrez \textperiodcentered{} história 9 versões no jornal}

%CRISTAL {"arestas":[{"alvo":"partida_imortal","confianca":1.0,"epistemico":"fato","origem":"wiki","rel":"parte_de"},{"alvo":"xeque_mate","confianca":1.0,"epistemico":"fato","origem":"wiki","rel":"relaciona"}],"confianca":1.0,"contraexemplos":[],"descricao":"Xeque-mate real, tabuleiro e peças:\n8 r . b k . . . r\n7 p . . p B p N p\n6 n . . . . n . .\n5 . p . N P . . P\n4 . . . . . . P .\n3 . . . P . . . .\n2 P . P . K . . .\n1 q . . . . . b .\n  a b c d e f g h\nFEN: r1bk3r/p2pBpNp/n4n2/1p1NP2P/6P1/3P4/P1P1K3/q5b1 b - - 0 1","epistemico":"fato","exemplos":[],"id":"pos_imortal_final","memoria":"semantica","meta":{"dominio":"xadrez","fen":"r1bk3r/p2pBpNp/n4n2/1p1NP2P/6P1/3P4/P1P1K3/q5b1 b - - 0 1","fonte":"posições reais","partida":"imortal","tabuleiro":"8 r . b k . . . r\n7 p . . p B p N p\n6 n . . . . n . .\n5 . p . N P . . P\n4 . . . . . . P .\n3 . . . P . . . .\n2 P . P . K . . .\n1 q . . . . . b .\n  a b c d e f g h"},"origem":"wiki","palavras_chave":[],"sinonimos":[],"tipo":"conceito","titulo":"Posição final — A Imortal (Anderssen, 1851)"}
\section{Posição final — A Imortal (Anderssen, 1851)}
Xeque-mate real, tabuleiro e peças:
8 r . b k . . . r
7 p . . p B p N p
6 n . . . . n . .
5 . p . N P . . P
4 . . . . . . P .
3 . . . P . . . .
2 P . P . K . . .
1 q . . . . . b .
  a b c d e f g h
FEN: r1bk3r/p2pBpNp/n4n2/1p1NP2P/6P1/3P4/P1P1K3/q5b1 b - - 0 1
\prov{relações: parte\_de partida\_imortal; relaciona xeque\_mate}
\prov{proveniência: wiki \textperiodcentered{} posições reais \textperiodcentered{} fato \textperiodcentered{} confiança 1.0 \textperiodcentered{} domínio xadrez \textperiodcentered{} história 25 versões no jornal}

%CRISTAL {"arestas":[{"alvo":"partida_lasker_thomas","confianca":1.0,"epistemico":"fato","origem":"wiki","rel":"parte_de"},{"alvo":"xeque_mate","confianca":1.0,"epistemico":"fato","origem":"wiki","rel":"relaciona"}],"confianca":1.0,"contraexemplos":[],"descricao":"Xeque-mate real, tabuleiro e peças:\n8 r n . . . r . .\n7 p b p p q . p .\n6 . p . . p N . .\n5 . . . . . . . .\n4 . . . P . . N P\n3 . . . . . . P .\n2 P P P K B P . R\n1 R . . . . . k .\n  a b c d e f g h\nFEN: rn3r2/pbppq1p1/1p2pN2/8/3P2NP/6P1/PPPKBP1R/R5k1 b - - 0 1","epistemico":"fato","exemplos":[],"id":"pos_lasker_thomas_final","memoria":"semantica","meta":{"dominio":"xadrez","fen":"rn3r2/pbppq1p1/1p2pN2/8/3P2NP/6P1/PPPKBP1R/R5k1 b - - 0 1","fonte":"posições reais","partida":"lasker_thomas","tabuleiro":"8 r n . . . r . .\n7 p b p p q . p .\n6 . p . . p N . .\n5 . . . . . . . .\n4 . . . P . . N P\n3 . . . . . . P .\n2 P P P K B P . R\n1 R . . . . . k .\n  a b c d e f g h"},"origem":"wiki","palavras_chave":[],"sinonimos":[],"tipo":"conceito","titulo":"Posição final — A Caminhada do Rei (Ed. Lasker × Thomas, 1912)"}
\section{Posição final — A Caminhada do Rei (Ed. Lasker $\times$ Thomas, 1912)}
Xeque-mate real, tabuleiro e peças:
8 r n . . . r . .
7 p b p p q . p .
6 . p . . p N . .
5 . . . . . . . .
4 . . . P . . N P
3 . . . . . . P .
2 P P P K B P . R
1 R . . . . . k .
  a b c d e f g h
FEN: rn3r2/pbppq1p1/1p2pN2/8/3P2NP/6P1/PPPKBP1R/R5k1 b - - 0 1
\prov{relações: parte\_de partida\_lasker\_thomas; relaciona xeque\_mate}
\prov{proveniência: wiki \textperiodcentered{} posições reais \textperiodcentered{} fato \textperiodcentered{} confiança 1.0 \textperiodcentered{} domínio xadrez \textperiodcentered{} história 6 versões no jornal}

%CRISTAL {"arestas":[{"alvo":"partida_legal","confianca":1.0,"epistemico":"fato","origem":"wiki","rel":"parte_de"},{"alvo":"xeque_mate","confianca":1.0,"epistemico":"fato","origem":"wiki","rel":"relaciona"}],"confianca":1.0,"contraexemplos":[],"descricao":"Xeque-mate real, tabuleiro e peças:\n8 r n . q . b n r\n7 p p p . k B . p\n6 . . . p . . p .\n5 . . . N N . . .\n4 . . . . P . . .\n3 . . . . . . . .\n2 P P P P . P P P\n1 R . B b K . . R\n  a b c d e f g h\nFEN: rn1q1bnr/ppp1kB1p/3p2p1/3NN3/4P3/8/PPPP1PPP/R1BbK2R b - - 0 1","epistemico":"fato","exemplos":[],"id":"pos_legal_final","memoria":"semantica","meta":{"dominio":"xadrez","fen":"rn1q1bnr/ppp1kB1p/3p2p1/3NN3/4P3/8/PPPP1PPP/R1BbK2R b - - 0 1","fonte":"posições reais","partida":"legal","tabuleiro":"8 r n . q . b n r\n7 p p p . k B . p\n6 . . . p . . p .\n5 . . . N N . . .\n4 . . . . P . . .\n3 . . . . . . . .\n2 P P P P . P P P\n1 R . B b K . . R\n  a b c d e f g h"},"origem":"wiki","palavras_chave":[],"sinonimos":[],"tipo":"conceito","titulo":"Posição final — Mate de Légal (1750)"}
\section{Posição final — Mate de Légal (1750)}
Xeque-mate real, tabuleiro e peças:
8 r n . q . b n r
7 p p p . k B . p
6 . . . p . . p .
5 . . . N N . . .
4 . . . . P . . .
3 . . . . . . . .
2 P P P P . P P P
1 R . B b K . . R
  a b c d e f g h
FEN: rn1q1bnr/ppp1kB1p/3p2p1/3NN3/4P3/8/PPPP1PPP/R1BbK2R b - - 0 1
\prov{relações: parte\_de partida\_legal; relaciona xeque\_mate}
\prov{proveniência: wiki \textperiodcentered{} posições reais \textperiodcentered{} fato \textperiodcentered{} confiança 1.0 \textperiodcentered{} domínio xadrez \textperiodcentered{} história 25 versões no jornal}

%CRISTAL {"arestas":[{"alvo":"partida_opera","confianca":1.0,"epistemico":"fato","origem":"wiki","rel":"parte_de"},{"alvo":"xeque_mate","confianca":1.0,"epistemico":"fato","origem":"wiki","rel":"relaciona"}],"confianca":1.0,"contraexemplos":[],"descricao":"Xeque-mate real, tabuleiro e peças:\n8 . n . R k b . r\n7 p . . . . p p p\n6 . . . . q . . .\n5 . . . . p . B .\n4 . . . . P . . .\n3 . . . . . . . .\n2 P P P . . P P P\n1 . . K . . . . .\n  a b c d e f g h\nFEN: 1n1Rkb1r/p4ppp/4q3/4p1B1/4P3/8/PPP2PPP/2K5 b - - 0 1","epistemico":"fato","exemplos":[],"id":"pos_opera_final","memoria":"semantica","meta":{"dominio":"xadrez","fen":"1n1Rkb1r/p4ppp/4q3/4p1B1/4P3/8/PPP2PPP/2K5 b - - 0 1","fonte":"posições reais","partida":"opera","tabuleiro":"8 . n . R k b . r\n7 p . . . . p p p\n6 . . . . q . . .\n5 . . . . p . B .\n4 . . . . P . . .\n3 . . . . . . . .\n2 P P P . . P P P\n1 . . K . . . . .\n  a b c d e f g h"},"origem":"wiki","palavras_chave":[],"sinonimos":[],"tipo":"conceito","titulo":"Posição final — Partida da Ópera (Morphy, 1858)"}
\section{Posição final — Partida da Ópera (Morphy, 1858)}
Xeque-mate real, tabuleiro e peças:
8 . n . R k b . r
7 p . . . . p p p
6 . . . . q . . .
5 . . . . p . B .
4 . . . . P . . .
3 . . . . . . . .
2 P P P . . P P P
1 . . K . . . . .
  a b c d e f g h
FEN: 1n1Rkb1r/p4ppp/4q3/4p1B1/4P3/8/PPP2PPP/2K5 b - - 0 1
\prov{relações: parte\_de partida\_opera; relaciona xeque\_mate}
\prov{proveniência: wiki \textperiodcentered{} posições reais \textperiodcentered{} fato \textperiodcentered{} confiança 1.0 \textperiodcentered{} domínio xadrez \textperiodcentered{} história 25 versões no jornal}

%CRISTAL {"arestas":[{"alvo":"partida_pastor","confianca":1.0,"epistemico":"fato","origem":"wiki","rel":"parte_de"},{"alvo":"xeque_mate","confianca":1.0,"epistemico":"fato","origem":"wiki","rel":"relaciona"}],"confianca":1.0,"contraexemplos":[],"descricao":"Xeque-mate real, tabuleiro e peças:\n8 r . b q k b . r\n7 p p p p . Q p p\n6 . . n . . n . .\n5 . . . . p . . .\n4 . . B . P . . .\n3 . . . . . . . .\n2 P P P P . P P P\n1 R N B . K . N R\n  a b c d e f g h\nFEN: r1bqkb1r/pppp1Qpp/2n2n2/4p3/2B1P3/8/PPPP1PPP/RNB1K1NR b - - 0 1","epistemico":"fato","exemplos":[],"id":"pos_pastor_final","memoria":"semantica","meta":{"dominio":"xadrez","fen":"r1bqkb1r/pppp1Qpp/2n2n2/4p3/2B1P3/8/PPPP1PPP/RNB1K1NR b - - 0 1","fonte":"posições reais","partida":"pastor","tabuleiro":"8 r . b q k b . r\n7 p p p p . Q p p\n6 . . n . . n . .\n5 . . . . p . . .\n4 . . B . P . . .\n3 . . . . . . . .\n2 P P P P . P P P\n1 R N B . K . N R\n  a b c d e f g h"},"origem":"wiki","palavras_chave":[],"sinonimos":[],"tipo":"conceito","titulo":"Posição final — Mate do Pastor"}
\section{Posição final — Mate do Pastor}
Xeque-mate real, tabuleiro e peças:
8 r . b q k b . r
7 p p p p . Q p p
6 . . n . . n . .
5 . . . . p . . .
4 . . B . P . . .
3 . . . . . . . .
2 P P P P . P P P
1 R N B . K . N R
  a b c d e f g h
FEN: r1bqkb1r/pppp1Qpp/2n2n2/4p3/2B1P3/8/PPPP1PPP/RNB1K1NR b - - 0 1
\prov{relações: parte\_de partida\_pastor; relaciona xeque\_mate}
\prov{proveniência: wiki \textperiodcentered{} posições reais \textperiodcentered{} fato \textperiodcentered{} confiança 1.0 \textperiodcentered{} domínio xadrez \textperiodcentered{} história 25 versões no jornal}

%CRISTAL {"arestas":[{"alvo":"partida_peruana","confianca":1.0,"epistemico":"fato","origem":"wiki","rel":"parte_de"},{"alvo":"xeque_mate","confianca":1.0,"epistemico":"fato","origem":"wiki","rel":"relaciona"}],"confianca":1.0,"contraexemplos":[],"descricao":"Xeque-mate real, tabuleiro e peças:\n8 . . k r . . n r\n7 p . . n . p p p\n6 B . p . p . . .\n5 . . . . . . . .\n4 . P . P . B . .\n3 . . N . . . . P\n2 . P P K . P P .\n1 . . . . . . . q\n  a b c d e f g h\nFEN: 2kr2nr/p2n1ppp/B1p1p3/8/1P1P1B2/2N4P/1PPK1PP1/7q b - - 0 1","epistemico":"fato","exemplos":[],"id":"pos_peruana_final","memoria":"semantica","meta":{"dominio":"xadrez","fen":"2kr2nr/p2n1ppp/B1p1p3/8/1P1P1B2/2N4P/1PPK1PP1/7q b - - 0 1","fonte":"posições reais","partida":"peruana","tabuleiro":"8 . . k r . . n r\n7 p . . n . p p p\n6 B . p . p . . .\n5 . . . . . . . .\n4 . P . P . B . .\n3 . . N . . . . P\n2 . P P K . P P .\n1 . . . . . . . q\n  a b c d e f g h"},"origem":"wiki","palavras_chave":[],"sinonimos":[],"tipo":"conceito","titulo":"Posição final — A Imortal Peruana (Canal, 1934)"}
\section{Posição final — A Imortal Peruana (Canal, 1934)}
Xeque-mate real, tabuleiro e peças:
8 . . k r . . n r
7 p . . n . p p p
6 B . p . p . . .
5 . . . . . . . .
4 . P . P . B . .
3 . . N . . . . P
2 . P P K . P P .
1 . . . . . . . q
  a b c d e f g h
FEN: 2kr2nr/p2n1ppp/B1p1p3/8/1P1P1B2/2N4P/1PPK1PP1/7q b - - 0 1
\prov{relações: parte\_de partida\_peruana; relaciona xeque\_mate}
\prov{proveniência: wiki \textperiodcentered{} posições reais \textperiodcentered{} fato \textperiodcentered{} confiança 1.0 \textperiodcentered{} domínio xadrez \textperiodcentered{} história 9 versões no jornal}

%CRISTAL {"arestas":[{"alvo":"partida_reti","confianca":1.0,"epistemico":"fato","origem":"wiki","rel":"parte_de"},{"alvo":"xeque_mate","confianca":1.0,"epistemico":"fato","origem":"wiki","rel":"relaciona"}],"confianca":1.0,"contraexemplos":[],"descricao":"Xeque-mate real, tabuleiro e peças:\n8 r n b B . b . r\n7 p p k . . p p p\n6 . . p . . . . .\n5 . . . . q . . .\n4 . . . . n . . .\n3 . . . . . . . .\n2 P P P . . P P P\n1 . . K R . B N R\n  a b c d e f g h\nFEN: rnbB1b1r/ppk2ppp/2p5/4q3/4n3/8/PPP2PPP/2KR1BNR b - - 0 1","epistemico":"fato","exemplos":[],"id":"pos_reti_final","memoria":"semantica","meta":{"dominio":"xadrez","fen":"rnbB1b1r/ppk2ppp/2p5/4q3/4n3/8/PPP2PPP/2KR1BNR b - - 0 1","fonte":"posições reais","partida":"reti","tabuleiro":"8 r n b B . b . r\n7 p p k . . p p p\n6 . . p . . . . .\n5 . . . . q . . .\n4 . . . . n . . .\n3 . . . . . . . .\n2 P P P . . P P P\n1 . . K R . B N R\n  a b c d e f g h"},"origem":"wiki","palavras_chave":[],"sinonimos":[],"tipo":"conceito","titulo":"Posição final — Réti × Tartakower (Viena, 1910)"}
\section{Posição final — Réti $\times$ Tartakower (Viena, 1910)}
Xeque-mate real, tabuleiro e peças:
8 r n b B . b . r
7 p p k . . p p p
6 . . p . . . . .
5 . . . . q . . .
4 . . . . n . . .
3 . . . . . . . .
2 P P P . . P P P
1 . . K R . B N R
  a b c d e f g h
FEN: rnbB1b1r/ppk2ppp/2p5/4q3/4n3/8/PPP2PPP/2KR1BNR b - - 0 1
\prov{relações: parte\_de partida\_reti; relaciona xeque\_mate}
\prov{proveniência: wiki \textperiodcentered{} posições reais \textperiodcentered{} fato \textperiodcentered{} confiança 1.0 \textperiodcentered{} domínio xadrez \textperiodcentered{} história 18 versões no jornal}

%CRISTAL {"arestas":[{"alvo":"partida_tiffany_d4","confianca":1.0,"epistemico":"fato","origem":"wiki","rel":"parte_de"}],"confianca":1.0,"contraexemplos":[],"descricao":"Posição final (em jogo), tabuleiro e peças:\n8 . r . k . b . r\n7 p . p q . . p p\n6 Q . . . . . . .\n5 . . p p P . . .\n4 . . . . . . . .\n3 P . . . . . . .\n2 . P P . . P P P\n1 R . B . K . . R\n  a b c d e f g h\nFEN: 1r1k1b1r/p1pq2pp/Q7/2ppP3/8/P7/1PP2PPP/R1B1K2R w - - 0 1","epistemico":"fato","exemplos":[],"id":"pos_tiffany_d4_final","memoria":"semantica","meta":{"dominio":"xadrez","fen":"1r1k1b1r/p1pq2pp/Q7/2ppP3/8/P7/1PP2PPP/R1B1K2R w - - 0 1","fonte":"posições reais","partida":"tiffany_d4","tabuleiro":"8 . r . k . b . r\n7 p . p q . . p p\n6 Q . . . . . . .\n5 . . p p P . . .\n4 . . . . . . . .\n3 P . . . . . . .\n2 . P P . . P P P\n1 R . B . K . . R\n  a b c d e f g h"},"origem":"wiki","palavras_chave":[],"sinonimos":[],"tipo":"conceito","titulo":"Posição final — A Órbita da Tiffany (auto-jogo, δ=4)"}
\section{Posição final — A Órbita da Tiffany (auto-jogo, \ensuremath{\delta}=4)}
Posição final (em jogo), tabuleiro e peças:
8 . r . k . b . r
7 p . p q . . p p
6 Q . . . . . . .
5 . . p p P . . .
4 . . . . . . . .
3 P . . . . . . .
2 . P P . . P P P
1 R . B . K . . R
  a b c d e f g h
FEN: 1r1k1b1r/p1pq2pp/Q7/2ppP3/8/P7/1PP2PPP/R1B1K2R w - - 0 1
\prov{relações: parte\_de partida\_tiffany\_d4}
\prov{proveniência: wiki \textperiodcentered{} posições reais \textperiodcentered{} fato \textperiodcentered{} confiança 1.0 \textperiodcentered{} domínio xadrez \textperiodcentered{} história 12 versões no jornal}

%CRISTAL {"arestas":[{"alvo":"partida_tiffany_d5","confianca":1.0,"epistemico":"fato","origem":"wiki","rel":"parte_de"}],"confianca":1.0,"contraexemplos":[],"descricao":"Posição final (em jogo), tabuleiro e peças:\n8 r . b B . k n r\n7 p p . p . p p p\n6 . . . . p . . .\n5 . B . . P . . .\n4 . q . . . . . .\n3 . . N . . Q . .\n2 . . P K . P P P\n1 R . . . . . . R\n  a b c d e f g h\nFEN: r1bB1knr/pp1p1ppp/4p3/1B2P3/1q6/2N2Q2/2PK1PPP/R6R w - - 0 1","epistemico":"fato","exemplos":[],"id":"pos_tiffany_d5_final","memoria":"semantica","meta":{"dominio":"xadrez","fen":"r1bB1knr/pp1p1ppp/4p3/1B2P3/1q6/2N2Q2/2PK1PPP/R6R w - - 0 1","fonte":"posições reais","partida":"tiffany_d5","tabuleiro":"8 r . b B . k n r\n7 p p . p . p p p\n6 . . . . p . . .\n5 . B . . P . . .\n4 . q . . . . . .\n3 . . N . . Q . .\n2 . . P K . P P P\n1 R . . . . . . R\n  a b c d e f g h"},"origem":"wiki","palavras_chave":[],"sinonimos":[],"tipo":"conceito","titulo":"Posição final — A Órbita da Tiffany (auto-jogo, δ=5)"}
\section{Posição final — A Órbita da Tiffany (auto-jogo, \ensuremath{\delta}=5)}
Posição final (em jogo), tabuleiro e peças:
8 r . b B . k n r
7 p p . p . p p p
6 . . . . p . . .
5 . B . . P . . .
4 . q . . . . . .
3 . . N . . Q . .
2 . . P K . P P P
1 R . . . . . . R
  a b c d e f g h
FEN: r1bB1knr/pp1p1ppp/4p3/1B2P3/1q6/2N2Q2/2PK1PPP/R6R w - - 0 1
\prov{relações: parte\_de partida\_tiffany\_d5}
\prov{proveniência: wiki \textperiodcentered{} posições reais \textperiodcentered{} fato \textperiodcentered{} confiança 1.0 \textperiodcentered{} domínio xadrez \textperiodcentered{} história 8 versões no jornal}

%CRISTAL {"arestas":[{"alvo":"partida_tiffany_d6","confianca":1.0,"epistemico":"fato","origem":"wiki","rel":"parte_de"}],"confianca":1.0,"contraexemplos":[],"descricao":"Posição final (em jogo), tabuleiro e peças:\n8 r . b . k b . r\n7 . p p . q p p p\n6 p . p . p n . .\n5 . . . . . . . .\n4 . . . P . . . .\n3 P . . . P N . .\n2 . P P . N P P P\n1 R . B Q K . . R\n  a b c d e f g h\nFEN: r1b1kb1r/1pp1qppp/p1p1pn2/8/3P4/P3PN2/1PP1NPPP/R1BQK2R w - - 0 1","epistemico":"fato","exemplos":[],"id":"pos_tiffany_d6_final","memoria":"semantica","meta":{"dominio":"xadrez","fen":"r1b1kb1r/1pp1qppp/p1p1pn2/8/3P4/P3PN2/1PP1NPPP/R1BQK2R w - - 0 1","fonte":"posições reais","partida":"tiffany_d6","tabuleiro":"8 r . b . k b . r\n7 . p p . q p p p\n6 p . p . p n . .\n5 . . . . . . . .\n4 . . . P . . . .\n3 P . . . P N . .\n2 . P P . N P P P\n1 R . B Q K . . R\n  a b c d e f g h"},"origem":"wiki","palavras_chave":[],"sinonimos":[],"tipo":"conceito","titulo":"Posição final — A Órbita da Tiffany (auto-jogo, δ=6)"}
\section{Posição final — A Órbita da Tiffany (auto-jogo, \ensuremath{\delta}=6)}
Posição final (em jogo), tabuleiro e peças:
8 r . b . k b . r
7 . p p . q p p p
6 p . p . p n . .
5 . . . . . . . .
4 . . . P . . . .
3 P . . . P N . .
2 . P P . N P P P
1 R . B Q K . . R
  a b c d e f g h
FEN: r1b1kb1r/1pp1qppp/p1p1pn2/8/3P4/P3PN2/1PP1NPPP/R1BQK2R w - - 0 1
\prov{relações: parte\_de partida\_tiffany\_d6}
\prov{proveniência: wiki \textperiodcentered{} posições reais \textperiodcentered{} fato \textperiodcentered{} confiança 1.0 \textperiodcentered{} domínio xadrez \textperiodcentered{} história 2 versões no jornal}

%CRISTAL {"arestas":[{"alvo":"budismo","confianca":1.0,"epistemico":"fato","origem":"wiki","rel":"parte_de"},{"alvo":"autopoiese","confianca":1.0,"epistemico":"fato","origem":"wiki","rel":"relaciona"}],"confianca":1.0,"contraexemplos":[],"descricao":"Buda / Nagarjuna: todos os fenômenos surgem em interdependência condicionada; nada existe por si, isolado ou autossuficiente — tudo é co-construído e interconectado.","epistemico":"inferencia","exemplos":[],"id":"pratityasamutpada","memoria":"semantica","meta":{"autor":"Nagarjuna","dominio":"mestres","fonte":""},"origem":"wiki","palavras_chave":[],"sinonimos":[],"tipo":"conceito","titulo":"Pratityasamutpada — a Originação Dependente"}
\section{Pratityasamutpada — a Originação Dependente}
Buda / Nagarjuna: todos os fenômenos surgem em interdependência condicionada; nada existe por si, isolado ou autossuficiente — tudo é co-construído e interconectado.
\prov{relações: parte\_de budismo; relaciona autopoiese}
\prov{proveniência: wiki \textperiodcentered{} inferencia \textperiodcentered{} confiança 1.0 \textperiodcentered{} domínio mestres \textperiodcentered{} história 3 versões no jornal}

%CRISTAL {"arestas":[{"alvo":"abertura_xadrez","confianca":1.0,"epistemico":"fato","origem":"wiki","rel":"usa"}],"confianca":1.0,"contraexemplos":[],"descricao":"Centro, desenvolvimento, segurança do rei: ocupe/controle o centro, tire as peças (não mova a mesma duas vezes), role, não saia com a dama cedo. As leis que a teoria concreta particulariza.","epistemico":"inferencia","exemplos":[],"id":"principios_de_abertura","memoria":"semantica","meta":{"dominio":"xadrez","fonte":"xadrez"},"origem":"wiki","palavras_chave":[],"sinonimos":[],"tipo":"conceito","titulo":"Princípios de Abertura"}
\section{Princípios de Abertura}
Centro, desenvolvimento, segurança do rei: ocupe/controle o centro, tire as peças (não mova a mesma duas vezes), role, não saia com a dama cedo. As leis que a teoria concreta particulariza.
\prov{relações: usa abertura\_xadrez}
\prov{proveniência: wiki \textperiodcentered{} xadrez \textperiodcentered{} inferencia \textperiodcentered{} confiança 1.0 \textperiodcentered{} domínio xadrez \textperiodcentered{} história 2 versões no jornal}

%CRISTAL {"arestas":[{"alvo":"consciencia_pura","confianca":1.0,"epistemico":"fato","origem":"wiki","rel":"relaciona"},{"alvo":"taoismo","confianca":1.0,"epistemico":"fato","origem":"wiki","rel":"parte_de"}],"confianca":1.0,"contraexemplos":[],"descricao":"Lao Tzu: o estado original, íntegro e não-processado, anterior ao recorte pela cultura e pelo desejo. 'A madeira bruta partida transforma-se em instrumentos' — cindir é já perder a inteireza.","epistemico":"inferencia","exemplos":[],"id":"pu","memoria":"semantica","meta":{"autor":"Lao Tzu","dominio":"mestres","fonte":""},"origem":"wiki","palavras_chave":[],"sinonimos":[],"tipo":"conceito","titulo":"Pu — a Simplicidade (o bloco não-esculpido)"}
\section{Pu — a Simplicidade (o bloco não-esculpido)}
Lao Tzu: o estado original, íntegro e não-processado, anterior ao recorte pela cultura e pelo desejo. 'A madeira bruta partida transforma-se em instrumentos' — cindir é já perder a inteireza.
\prov{relações: relaciona consciencia\_pura; parte\_de taoismo}
\prov{proveniência: wiki \textperiodcentered{} inferencia \textperiodcentered{} confiança 1.0 \textperiodcentered{} domínio mestres \textperiodcentered{} história 3 versões no jornal}

%CRISTAL {"arestas":[{"alvo":"budismo","confianca":1.0,"epistemico":"fato","origem":"wiki","rel":"parte_de"},{"alvo":"dor_vs_sofrimento_gentil","confianca":1.0,"epistemico":"fato","origem":"wiki","rel":"relaciona"}],"confianca":1.0,"contraexemplos":[],"descricao":"Buda: o diagnóstico — (1) há sofrimento (dukkha); (2) ele tem origem (o desejo/apego, tanha); (3) há cessação; (4) há um caminho (o Nobre Caminho Óctuplo) que conduz à cessação.","epistemico":"inferencia","exemplos":[],"id":"quatro_nobres_verdades","memoria":"semantica","meta":{"autor":"Buda","dominio":"mestres","fonte":""},"origem":"wiki","palavras_chave":[],"sinonimos":[],"tipo":"conceito","titulo":"As Quatro Nobres Verdades"}
\section{As Quatro Nobres Verdades}
Buda: o diagnóstico — (1) há sofrimento (dukkha); (2) ele tem origem (o desejo/apego, tanha); (3) há cessação; (4) há um caminho (o Nobre Caminho Óctuplo) que conduz à cessação.
\prov{relações: parte\_de budismo; relaciona dor\_vs\_sofrimento\_gentil}
\prov{proveniência: wiki \textperiodcentered{} inferencia \textperiodcentered{} confiança 1.0 \textperiodcentered{} domínio mestres \textperiodcentered{} história 3 versões no jornal}

%CRISTAL {"arestas":[{"alvo":"regua_de_gentil","confianca":1.0,"epistemico":"fato","origem":"wiki","rel":"usa"},{"alvo":"retorno_taoista","confianca":1.0,"epistemico":"fato","origem":"wiki","rel":"explica"},{"alvo":"tao_da_matematica","confianca":1.0,"epistemico":"fato","origem":"wiki","rel":"parte_de"}],"confianca":1.0,"contraexemplos":[],"descricao":"Lopes constrói uma métrica σ sobre [0,1[ (a 'régua divina', escala 0→0,5→0 de volta) em que, passada a metade, quanto mais o ponto se afasta da origem, menor sua distância — realizando matematicamente 'afastar-se significa retornar' (Lao Tsé) e 'está perto e ao mesmo tempo distante' (Gita).","epistemico":"inferencia","exemplos":[],"id":"regua_divina_afastar_retornar","memoria":"semantica","meta":{"autor":"Gentil Lopes","dominio":"mestres","fonte":""},"origem":"wiki","palavras_chave":[],"sinonimos":[],"tipo":"conceito","titulo":"A Régua Divina — afastar-se é retornar (Lopes)"}
\section{A Régua Divina — afastar-se é retornar (Lopes)}
Lopes constrói uma métrica \ensuremath{\sigma} sobre [0,1[ (a 'régua divina', escala 0\ensuremath{\to}0,5\ensuremath{\to}0 de volta) em que, passada a metade, quanto mais o ponto se afasta da origem, menor sua distância — realizando matematicamente 'afastar-se significa retornar' (Lao Tsé) e 'está perto e ao mesmo tempo distante' (Gita).
\prov{relações: usa regua\_de\_gentil; explica retorno\_taoista; parte\_de tao\_da\_matematica}
\prov{proveniência: wiki \textperiodcentered{} inferencia \textperiodcentered{} confiança 1.0 \textperiodcentered{} domínio mestres \textperiodcentered{} história 5 versões no jornal}

%CRISTAL {"arestas":[{"alvo":"confucionismo","confianca":1.0,"epistemico":"fato","origem":"wiki","rel":"parte_de"},{"alvo":"ethics_of_care","confianca":1.0,"epistemico":"fato","origem":"wiki","rel":"relaciona"},{"alvo":"li_ritual","confianca":1.0,"epistemico":"fato","origem":"wiki","rel":"usa"}],"confianca":1.0,"contraexemplos":[],"descricao":"Confúcio: a virtude moral suprema, o amor ao próximo — 'superar o eu e voltar-se à observância dos ritos'. Ideal quase inatingível, mas sempre ao alcance de quem o deseja: 'tão logo a desejo, ela está aqui'.","epistemico":"inferencia","exemplos":[],"id":"ren","memoria":"semantica","meta":{"autor":"Confúcio","dominio":"mestres","fonte":""},"origem":"wiki","palavras_chave":[],"sinonimos":[],"tipo":"conceito","titulo":"Ren — Benevolência / Humanidade"}
\section{Ren — Benevolência / Humanidade}
Confúcio: a virtude moral suprema, o amor ao próximo — 'superar o eu e voltar-se à observância dos ritos'. Ideal quase inatingível, mas sempre ao alcance de quem o deseja: 'tão logo a desejo, ela está aqui'.
\prov{relações: parte\_de confucionismo; relaciona ethics\_of\_care; usa li\_ritual}
\prov{proveniência: wiki \textperiodcentered{} inferencia \textperiodcentered{} confiança 1.0 \textperiodcentered{} domínio mestres \textperiodcentered{} história 4 versões no jornal}

%CRISTAL {"arestas":[{"alvo":"xadrez","confianca":1.0,"epistemico":"fato","origem":"wiki","rel":"relaciona"}],"confianca":1.0,"contraexemplos":[],"descricao":"Grande mestre e teórico do hipermodernismo (1889–1929), autor de miniaturas brilhantes e da célebre abertura que leva seu nome.","epistemico":"fato","exemplos":[],"id":"reti","memoria":"semantica","meta":{"dominio":"xadrez","fonte":"jogadores"},"origem":"wiki","palavras_chave":[],"sinonimos":[],"tipo":"conceito","titulo":"Richard Réti"}
\section{Richard Réti}
Grande mestre e teórico do hipermodernismo (1889–1929), autor de miniaturas brilhantes e da célebre abertura que leva seu nome.
\prov{relações: relaciona xadrez}
\prov{proveniência: wiki \textperiodcentered{} jogadores \textperiodcentered{} fato \textperiodcentered{} confiança 1.0 \textperiodcentered{} domínio xadrez \textperiodcentered{} história 12 versões no jornal}

%CRISTAL {"arestas":[{"alvo":"tao","confianca":1.0,"epistemico":"fato","origem":"wiki","rel":"usa"},{"alvo":"contemplacao","confianca":1.0,"epistemico":"fato","origem":"wiki","rel":"relaciona"},{"alvo":"consciencia_pura","confianca":1.0,"epistemico":"fato","origem":"wiki","rel":"relaciona"}],"confianca":1.0,"contraexemplos":[],"descricao":"Lao Tzu: o movimento próprio do Tao é o retorno — tudo regressa à quietude/origem. 'O retorno é o movimento do Caminho'; 'o regresso à raiz se chama quietude; conhecer a constância se chama iluminação'. Alcançado pelo vazio extremo e a quietude contemplativa.","epistemico":"inferencia","exemplos":[],"id":"retorno_taoista","memoria":"semantica","meta":{"autor":"Lao Tzu","dominio":"mestres","fonte":""},"origem":"wiki","palavras_chave":[],"sinonimos":[],"tipo":"conceito","titulo":"O Retorno à Raiz (fu)"}
\section{O Retorno à Raiz (fu)}
Lao Tzu: o movimento próprio do Tao é o retorno — tudo regressa à quietude/origem. 'O retorno é o movimento do Caminho'; 'o regresso à raiz se chama quietude; conhecer a constância se chama iluminação'. Alcançado pelo vazio extremo e a quietude contemplativa.
\prov{relações: usa tao; relaciona contemplacao; relaciona consciencia\_pura}
\prov{proveniência: wiki \textperiodcentered{} inferencia \textperiodcentered{} confiança 1.0 \textperiodcentered{} domínio mestres \textperiodcentered{} história 4 versões no jornal}

%CRISTAL {"arestas":[{"alvo":"abertura_xadrez","confianca":1.0,"epistemico":"fato","origem":"wiki","rel":"especializa"}],"confianca":1.0,"contraexemplos":[],"descricao":"O bispo pressiona o cavalo que defende o peão central — pressão posicional lenta e profunda. Um dos pilares do repertório clássico, favorita de Fischer e Kasparov.","epistemico":"fato","exemplos":[],"id":"ruy_lopez","memoria":"semantica","meta":{"dominio":"xadrez","fonte":"xadrez"},"origem":"wiki","palavras_chave":[],"sinonimos":[],"tipo":"conceito","titulo":"Abertura Espanhola (Ruy López, 1.e4 e5 2.Cf3 Cc6 3.Bb5)"}
\section{Abertura Espanhola (Ruy López, 1.e4 e5 2.Cf3 Cc6 3.Bb5)}
O bispo pressiona o cavalo que defende o peão central — pressão posicional lenta e profunda. Um dos pilares do repertório clássico, favorita de Fischer e Kasparov.
\prov{relações: especializa abertura\_xadrez}
\prov{proveniência: wiki \textperiodcentered{} xadrez \textperiodcentered{} fato \textperiodcentered{} confiança 1.0 \textperiodcentered{} domínio xadrez \textperiodcentered{} história 2 versões no jornal}

%CRISTAL {"arestas":[{"alvo":"filosofia","confianca":1.0,"epistemico":"fato","origem":"wiki","rel":"parte_de"},{"alvo":"contemplacao","confianca":1.0,"epistemico":"fato","origem":"wiki","rel":"relaciona"}],"confianca":1.0,"contraexemplos":[],"descricao":"As grandes tradições contemplativas e sapienciais do Oriente — taoísmo, confucionismo, vedanta, budismo, zen — que pensam o caminho (não só a verdade), o vazio, a não-dualidade e o cultivo de si.","epistemico":"fato","exemplos":[],"id":"sabedoria_oriental","memoria":"semantica","meta":{"dominio":"mestres","fonte":""},"origem":"wiki","palavras_chave":[],"sinonimos":[],"tipo":"conceito","titulo":"Sabedoria Oriental"}
\section{Sabedoria Oriental}
As grandes tradições contemplativas e sapienciais do Oriente — taoísmo, confucionismo, vedanta, budismo, zen — que pensam o caminho (não só a verdade), o vazio, a não-dualidade e o cultivo de si.
\prov{relações: parte\_de filosofia; relaciona contemplacao}
\prov{proveniência: wiki \textperiodcentered{} fato \textperiodcentered{} confiança 1.0 \textperiodcentered{} domínio mestres \textperiodcentered{} história 3 versões no jornal}

%CRISTAL {"arestas":[{"alvo":"tatica_xadrez","confianca":1.0,"epistemico":"fato","origem":"wiki","rel":"especializa"}],"confianca":1.0,"contraexemplos":[],"descricao":"Entregar material por uma vantagem maior — mate, ganho posterior ou ataque irresistível. A COMBINAÇÃO é a sequência tática que o sacrifício abre; a beleza do jogo mora aqui.","epistemico":"inferencia","exemplos":[],"id":"sacrificio_xadrez","memoria":"semantica","meta":{"dominio":"xadrez","fonte":"xadrez"},"origem":"wiki","palavras_chave":[],"sinonimos":[],"tipo":"conceito","titulo":"Sacrifício e Combinação"}
\section{Sacrifício e Combinação}
Entregar material por uma vantagem maior — mate, ganho posterior ou ataque irresistível. A COMBINAÇÃO é a sequência tática que o sacrifício abre; a beleza do jogo mora aqui.
\prov{relações: especializa tatica\_xadrez}
\prov{proveniência: wiki \textperiodcentered{} xadrez \textperiodcentered{} inferencia \textperiodcentered{} confiança 1.0 \textperiodcentered{} domínio xadrez \textperiodcentered{} história 2 versões no jornal}

%CRISTAL {"arestas":[{"alvo":"ren","confianca":1.0,"epistemico":"fato","origem":"wiki","rel":"usa"},{"alvo":"imperativo_categorico","confianca":1.0,"epistemico":"fato","origem":"wiki","rel":"relaciona"}],"confianca":1.0,"contraexemplos":[],"descricao":"Confúcio: 'não imponha aos outros o que não deseja para si' — o método negativo da benevolência (pôr-se no lugar do outro). Emparelha com zhong (dar o melhor de si): juntos, a via prática do ren.","epistemico":"inferencia","exemplos":[],"id":"shu","memoria":"semantica","meta":{"autor":"Confúcio","dominio":"mestres","fonte":""},"origem":"wiki","palavras_chave":[],"sinonimos":[],"tipo":"conceito","titulo":"Shu — a Regra de Ouro (Reciprocidade)"}
\section{Shu — a Regra de Ouro (Reciprocidade)}
Confúcio: 'não imponha aos outros o que não deseja para si' — o método negativo da benevolência (pôr-se no lugar do outro). Emparelha com zhong (dar o melhor de si): juntos, a via prática do ren.
\prov{relações: usa ren; relaciona imperativo\_categorico}
\prov{proveniência: wiki \textperiodcentered{} inferencia \textperiodcentered{} confiança 1.0 \textperiodcentered{} domínio mestres \textperiodcentered{} história 3 versões no jornal}

%CRISTAL {"arestas":[{"alvo":"estrategia","confianca":1.0,"epistemico":"fato","origem":"wiki","rel":"explica"},{"alvo":"estrategia_indireta","confianca":1.0,"epistemico":"fato","origem":"wiki","rel":"relaciona"}],"confianca":1.0,"contraexemplos":[],"descricao":"Vencer sem lutar é o cume da habilidade: conhecer a si e ao inimigo, o terreno, o momento; a guerra por engano, economia e posição. O clássico da estratégia INDIRETA, ~500 a.C.","epistemico":"fato","exemplos":[],"id":"sun_tzu_arte_guerra","memoria":"semantica","meta":{"autor":"Sun Tzu","dominio":"estrategia","fonte":"estratégia"},"origem":"wiki","palavras_chave":[],"sinonimos":[],"tipo":"conceito","titulo":"A Arte da Guerra (Sun Tzu)"}
\section{A Arte da Guerra (Sun Tzu)}
Vencer sem lutar é o cume da habilidade: conhecer a si e ao inimigo, o terreno, o momento; a guerra por engano, economia e posição. O clássico da estratégia INDIRETA, \~{}500 a.C.
\prov{relações: explica estrategia; relaciona estrategia\_indireta}
\prov{proveniência: wiki \textperiodcentered{} estratégia \textperiodcentered{} fato \textperiodcentered{} confiança 1.0 \textperiodcentered{} domínio estrategia \textperiodcentered{} história 3 versões no jornal}

%CRISTAL {"arestas":[{"alvo":"budismo","confianca":1.0,"epistemico":"fato","origem":"wiki","rel":"parte_de"},{"alvo":"borda_critica_vazio","confianca":1.0,"epistemico":"fato","origem":"wiki","rel":"usa"},{"alvo":"pratityasamutpada","confianca":1.0,"epistemico":"fato","origem":"wiki","rel":"usa"},{"alvo":"vazio_util","confianca":1.0,"epistemico":"fato","origem":"wiki","rel":"relaciona"}],"confianca":1.0,"contraexemplos":[],"descricao":"Nagarjuna (Madhyamaka): as coisas são vazias de essência/natureza intrínseca (svabhava) porque tudo é dependentemente originado — e a própria vacuidade também é vazia. Vazio não é inexistência, mas ausência de fundamento intrínseco.","epistemico":"inferencia","exemplos":[],"id":"sunyata","memoria":"semantica","meta":{"autor":"Nagarjuna","dominio":"mestres","fonte":""},"origem":"wiki","palavras_chave":[],"sinonimos":[],"tipo":"conceito","titulo":"Sunyata — a Vacuidade"}
\section{Sunyata — a Vacuidade}
Nagarjuna (Madhyamaka): as coisas são vazias de essência/natureza intrínseca (svabhava) porque tudo é dependentemente originado — e a própria vacuidade também é vazia. Vazio não é inexistência, mas ausência de fundamento intrínseco.
\prov{relações: parte\_de budismo; usa borda\_critica\_vazio; usa pratityasamutpada; relaciona vazio\_util}
\prov{proveniência: wiki \textperiodcentered{} inferencia \textperiodcentered{} confiança 1.0 \textperiodcentered{} domínio mestres \textperiodcentered{} história 5 versões no jornal}

%CRISTAL {"arestas":[{"alvo":"taoismo","confianca":1.0,"epistemico":"fato","origem":"wiki","rel":"parte_de"},{"alvo":"simbolo_vs_significado","confianca":1.0,"epistemico":"fato","origem":"wiki","rel":"relaciona"},{"alvo":"nao_dualidade","confianca":1.0,"epistemico":"fato","origem":"wiki","rel":"relaciona"},{"alvo":"sentido_da_vida","confianca":1.0,"epistemico":"fato","origem":"wiki","rel":"relaciona"}],"confianca":1.0,"contraexemplos":[],"descricao":"Lao Tzu: o princípio último, indizível e inesgotável, fonte de todos os seres, anterior a toda distinção — 'o caminho que pode ser nomeado não é o Caminho constante'. O que se nomeia já não é o Tao eterno.","epistemico":"inferencia","exemplos":[],"id":"tao","memoria":"semantica","meta":{"autor":"Lao Tzu","dominio":"mestres","fonte":""},"origem":"wiki","palavras_chave":[],"sinonimos":[],"tipo":"conceito","titulo":"O Tao (o Caminho não-nomeável)"}
\section{O Tao (o Caminho não-nomeável)}
Lao Tzu: o princípio último, indizível e inesgotável, fonte de todos os seres, anterior a toda distinção — 'o caminho que pode ser nomeado não é o Caminho constante'. O que se nomeia já não é o Tao eterno.
\prov{relações: parte\_de taoismo; relaciona simbolo\_vs\_significado; relaciona nao\_dualidade; relaciona sentido\_da\_vida}
\prov{proveniência: wiki \textperiodcentered{} inferencia \textperiodcentered{} confiança 1.0 \textperiodcentered{} domínio mestres \textperiodcentered{} história 5 versões no jornal}

%CRISTAL {"arestas":[{"alvo":"sabedoria_oriental","confianca":1.0,"epistemico":"fato","origem":"wiki","rel":"relaciona"},{"alvo":"matematica_arte_engenharia","confianca":1.0,"epistemico":"fato","origem":"wiki","rel":"usa"},{"alvo":"tao","confianca":1.0,"epistemico":"fato","origem":"wiki","rel":"usa"},{"alvo":"capra_tao_fisica","confianca":1.0,"epistemico":"fato","origem":"wiki","rel":"usa"}],"confianca":1.0,"contraexemplos":[],"descricao":"Lopes (livro de 2011): ecoando 'O Tao da Física' de Capra, lê o Tao Te Ching como 'O Caminho que leva à Divindade' e seu livro como 'o Caminho que leva à Divindade através da Matemática' (Tao ≡ Deus). A matemática como via contemplativa e espiritual.","epistemico":"inferencia","exemplos":[],"id":"tao_da_matematica","memoria":"semantica","meta":{"autor":"Gentil Lopes","dominio":"mestres","fonte":""},"origem":"wiki","palavras_chave":[],"sinonimos":[],"tipo":"conceito","titulo":"O Tao da Matemática (Lopes)"}
\section{O Tao da Matemática (Lopes)}
Lopes (livro de 2011): ecoando 'O Tao da Física' de Capra, lê o Tao Te Ching como 'O Caminho que leva à Divindade' e seu livro como 'o Caminho que leva à Divindade através da Matemática' (Tao \ensuremath{\equiv} Deus). A matemática como via contemplativa e espiritual.
\prov{relações: relaciona sabedoria\_oriental; usa matematica\_arte\_engenharia; usa tao; usa capra\_tao\_fisica}
\prov{proveniência: wiki \textperiodcentered{} inferencia \textperiodcentered{} confiança 1.0 \textperiodcentered{} domínio mestres \textperiodcentered{} história 6 versões no jornal}

%CRISTAL {"arestas":[{"alvo":"sabedoria_oriental","confianca":1.0,"epistemico":"fato","origem":"wiki","rel":"parte_de"},{"alvo":"word","confianca":0.5,"epistemico":"fato","origem":"wiki","rel":"relaciona"},{"alvo":"virtualizacao","confianca":0.5,"epistemico":"fato","origem":"wiki","rel":"relaciona"}],"confianca":1.0,"contraexemplos":[],"descricao":"A tradição do Tao Te Ching — viver em harmonia com o Tao (o Caminho), pelo não-forçar (wu wei), a naturalidade e o vazio fecundo.","epistemico":"fato","exemplos":[],"id":"taoismo","memoria":"semantica","meta":{"autor":"Lao Tzu","dominio":"mestres","fonte":""},"origem":"wiki","palavras_chave":[],"sinonimos":[],"tipo":"conceito","titulo":"Taoísmo (Lao Tzu)"}
\section{Taoísmo (Lao Tzu)}
A tradição do Tao Te Ching — viver em harmonia com o Tao (o Caminho), pelo não-forçar (wu wei), a naturalidade e o vazio fecundo.
\prov{relações: parte\_de sabedoria\_oriental; relaciona word; relaciona virtualizacao}
\prov{proveniência: wiki \textperiodcentered{} fato \textperiodcentered{} confiança 1.0 \textperiodcentered{} domínio mestres \textperiodcentered{} história 4 versões no jornal}

%CRISTAL {"arestas":[{"alvo":"xadrez","confianca":1.0,"epistemico":"fato","origem":"wiki","rel":"parte_de"}],"confianca":1.0,"contraexemplos":[],"descricao":"Sequências forçadas de curto prazo que ganham material ou dão mate — o cálculo concreto de variantes. 'A tática é o que fazer quando há o que fazer' (Tartakower). O oposto da estratégia lenta.","epistemico":"inferencia","exemplos":[],"id":"tatica_xadrez","memoria":"semantica","meta":{"dominio":"xadrez","fonte":"xadrez"},"origem":"wiki","palavras_chave":[],"sinonimos":[],"tipo":"conceito","titulo":"Tática"}
\section{Tática}
Sequências forçadas de curto prazo que ganham material ou dão mate — o cálculo concreto de variantes. 'A tática é o que fazer quando há o que fazer' (Tartakower). O oposto da estratégia lenta.
\prov{relações: parte\_de xadrez}
\prov{proveniência: wiki \textperiodcentered{} xadrez \textperiodcentered{} inferencia \textperiodcentered{} confiança 1.0 \textperiodcentered{} domínio xadrez \textperiodcentered{} história 2 versões no jornal}

%CRISTAL {"arestas":[{"alvo":"tao","confianca":1.0,"epistemico":"fato","origem":"wiki","rel":"usa"}],"confianca":1.0,"contraexemplos":[],"descricao":"Lao Tzu: a manifestação do Tao no ser concreto — potência que gera, nutre e sustenta sem se apossar nem dominar. 'A Virtude Superior não é virtude, assim possui a Virtude': virtude oculta, não intencional.","epistemico":"inferencia","exemplos":[],"id":"te_virtude","memoria":"semantica","meta":{"autor":"Lao Tzu","dominio":"mestres","fonte":""},"origem":"wiki","palavras_chave":[],"sinonimos":[],"tipo":"conceito","titulo":"Te — a Virtude / Poder"}
\section{Te — a Virtude / Poder}
Lao Tzu: a manifestação do Tao no ser concreto — potência que gera, nutre e sustenta sem se apossar nem dominar. 'A Virtude Superior não é virtude, assim possui a Virtude': virtude oculta, não intencional.
\prov{relações: usa tao}
\prov{proveniência: wiki \textperiodcentered{} inferencia \textperiodcentered{} confiança 1.0 \textperiodcentered{} domínio mestres \textperiodcentered{} história 2 versões no jornal}

%CRISTAL {"arestas":[{"alvo":"tecnicas_de_mate","confianca":1.0,"epistemico":"fato","origem":"wiki","rel":"parte_de"},{"alvo":"xeque_mate","confianca":1.0,"epistemico":"fato","origem":"wiki","rel":"relaciona"}],"confianca":1.0,"contraexemplos":[],"descricao":"O padrão mais antigo com nome: cavalo em f6 cobre g8, a torre entra em h7 — Th7#. Torre e cavalo, a dupla clássica do canto.\n8 . . . . . . . k\n7 . . . . . . . R\n6 . . . . . N . .\n5 . . . . . . . .\n4 . . . . . . . .\n3 . . . . . . . .\n2 . . . . . . . .\n1 . . . . . . K .\n  a b c d e f g h\nFEN inicial: 7k/8/5N2/8/8/8/8/6KR w - - 0 1","epistemico":"fato","exemplos":[],"id":"tecnica_mate_arabe","memoria":"semantica","meta":{"dominio":"xadrez","fen_final":"7k/7R/5N2/8/8/8/8/6K1 b - - 0 1","fen_inicial":"7k/8/5N2/8/8/8/8/6KR w - - 0 1","fonte":"técnicas de mate","lances_coord":["h1h7"],"orbita_fen":["7k/8/5N2/8/8/8/8/6KR w - - 0 1","7k/7R/5N2/8/8/8/8/6K1 b - - 0 1"],"tabuleiro_final":"8 . . . . . . . k\n7 . . . . . . . R\n6 . . . . . N . .\n5 . . . . . . . .\n4 . . . . . . . .\n3 . . . . . . . .\n2 . . . . . . . .\n1 . . . . . . K .\n  a b c d e f g h"},"origem":"wiki","palavras_chave":[],"sinonimos":[],"tipo":"conceito","titulo":"Mate Árabe (torre + cavalo)"}
\section{Mate Árabe (torre + cavalo)}
O padrão mais antigo com nome: cavalo em f6 cobre g8, a torre entra em h7 — Th7\#. Torre e cavalo, a dupla clássica do canto.
8 . . . . . . . k
7 . . . . . . . R
6 . . . . . N . .
5 . . . . . . . .
4 . . . . . . . .
3 . . . . . . . .
2 . . . . . . . .
1 . . . . . . K .
  a b c d e f g h
FEN inicial: 7k/8/5N2/8/8/8/8/6KR w - - 0 1
\prov{relações: parte\_de tecnicas\_de\_mate; relaciona xeque\_mate}
\prov{proveniência: wiki \textperiodcentered{} técnicas de mate \textperiodcentered{} fato \textperiodcentered{} confiança 1.0 \textperiodcentered{} domínio xadrez \textperiodcentered{} história 3 versões no jornal}

%CRISTAL {"arestas":[{"alvo":"tecnicas_de_mate","confianca":1.0,"epistemico":"fato","origem":"wiki","rel":"parte_de"},{"alvo":"xeque_mate","confianca":1.0,"epistemico":"fato","origem":"wiki","rel":"relaciona"}],"confianca":1.0,"contraexemplos":[],"descricao":"O mais difícil dos mates básicos: encurralar o rei no canto DA COR do bispo (aqui a8, claro), com a manobra em W do cavalo. O net final: Be4# — bispo dá o xeque, cavalo e rei cobrem as fugas.\n8 k . . . . . . .\n7 . . . . . . . .\n6 N K . . . . . .\n5 . . . . . . . .\n4 . . . . B . . .\n3 . . . . . . . .\n2 . . . . . . . .\n1 . . . . . . . .\n  a b c d e f g h\nFEN inicial: k7/8/NK6/8/8/5B2/8/8 w - - 0 1","epistemico":"fato","exemplos":[],"id":"tecnica_mate_bispo_cavalo","memoria":"semantica","meta":{"dominio":"xadrez","fen_final":"k7/8/NK6/8/4B3/8/8/8 b - - 0 1","fen_inicial":"k7/8/NK6/8/8/5B2/8/8 w - - 0 1","fonte":"técnicas de mate","lances_coord":["f3e4"],"orbita_fen":["k7/8/NK6/8/8/5B2/8/8 w - - 0 1","k7/8/NK6/8/4B3/8/8/8 b - - 0 1"],"tabuleiro_final":"8 k . . . . . . .\n7 . . . . . . . .\n6 N K . . . . . .\n5 . . . . . . . .\n4 . . . . B . . .\n3 . . . . . . . .\n2 . . . . . . . .\n1 . . . . . . . .\n  a b c d e f g h"},"origem":"wiki","palavras_chave":[],"sinonimos":[],"tipo":"conceito","titulo":"Mate de Bispo e Cavalo (KBN vs K)"}
\section{Mate de Bispo e Cavalo (KBN vs K)}
O mais difícil dos mates básicos: encurralar o rei no canto DA COR do bispo (aqui a8, claro), com a manobra em W do cavalo. O net final: Be4\# — bispo dá o xeque, cavalo e rei cobrem as fugas.
8 k . . . . . . .
7 . . . . . . . .
6 N K . . . . . .
5 . . . . . . . .
4 . . . . B . . .
3 . . . . . . . .
2 . . . . . . . .
1 . . . . . . . .
  a b c d e f g h
FEN inicial: k7/8/NK6/8/8/5B2/8/8 w - - 0 1
\prov{relações: parte\_de tecnicas\_de\_mate; relaciona xeque\_mate}
\prov{proveniência: wiki \textperiodcentered{} técnicas de mate \textperiodcentered{} fato \textperiodcentered{} confiança 1.0 \textperiodcentered{} domínio xadrez \textperiodcentered{} história 3 versões no jornal}

%CRISTAL {"arestas":[{"alvo":"tecnicas_de_mate","confianca":1.0,"epistemico":"fato","origem":"wiki","rel":"parte_de"},{"alvo":"xeque_mate","confianca":1.0,"epistemico":"fato","origem":"wiki","rel":"relaciona"}],"confianca":1.0,"contraexemplos":[],"descricao":"A dama encurrala o rei na borda e o rei branco a apoia: Dg7# — a dama toca o rei, o rei branco cobre a fuga. O mate básico mais simples.\n8 . . . . . . . k\n7 . . . . . . Q .\n6 . . . . . K . .\n5 . . . . . . . .\n4 . . . . . . . .\n3 . . . . . . . .\n2 . . . . . . . .\n1 . . . . . . . .\n  a b c d e f g h\nFEN inicial: 7k/8/5K2/8/8/8/8/6Q1 w - - 0 1","epistemico":"fato","exemplos":[],"id":"tecnica_mate_dama","memoria":"semantica","meta":{"dominio":"xadrez","fen_final":"7k/6Q1/5K2/8/8/8/8/8 b - - 0 1","fen_inicial":"7k/8/5K2/8/8/8/8/6Q1 w - - 0 1","fonte":"técnicas de mate","lances_coord":["g1g7"],"orbita_fen":["7k/8/5K2/8/8/8/8/6Q1 w - - 0 1","7k/6Q1/5K2/8/8/8/8/8 b - - 0 1"],"tabuleiro_final":"8 . . . . . . . k\n7 . . . . . . Q .\n6 . . . . . K . .\n5 . . . . . . . .\n4 . . . . . . . .\n3 . . . . . . . .\n2 . . . . . . . .\n1 . . . . . . . .\n  a b c d e f g h"},"origem":"wiki","palavras_chave":[],"sinonimos":[],"tipo":"conceito","titulo":"Mate da Dama (KQ vs K)"}
\section{Mate da Dama (KQ vs K)}
A dama encurrala o rei na borda e o rei branco a apoia: Dg7\# — a dama toca o rei, o rei branco cobre a fuga. O mate básico mais simples.
8 . . . . . . . k
7 . . . . . . Q .
6 . . . . . K . .
5 . . . . . . . .
4 . . . . . . . .
3 . . . . . . . .
2 . . . . . . . .
1 . . . . . . . .
  a b c d e f g h
FEN inicial: 7k/8/5K2/8/8/8/8/6Q1 w - - 0 1
\prov{relações: parte\_de tecnicas\_de\_mate; relaciona xeque\_mate}
\prov{proveniência: wiki \textperiodcentered{} técnicas de mate \textperiodcentered{} fato \textperiodcentered{} confiança 1.0 \textperiodcentered{} domínio xadrez \textperiodcentered{} história 3 versões no jornal}

%CRISTAL {"arestas":[{"alvo":"tecnicas_de_mate","confianca":1.0,"epistemico":"fato","origem":"wiki","rel":"parte_de"},{"alvo":"xeque_mate","confianca":1.0,"epistemico":"fato","origem":"wiki","rel":"relaciona"}],"confianca":1.0,"contraexemplos":[],"descricao":"Os dois bispos varrem duas diagonais paralelas, empurrando o rei ao canto (qualquer canto): Bd5# — um bispo dá o xeque, o outro corta a fuga, o rei branco fecha.\n8 k . . . . . . .\n7 . . . . . . . .\n6 . K . . . . . .\n5 . . . B . . . .\n4 . . . . . B . .\n3 . . . . . . . .\n2 . . . . . . . .\n1 . . . . . . . .\n  a b c d e f g h\nFEN inicial: k7/8/1K6/8/5B2/8/8/7B w - - 0 1","epistemico":"fato","exemplos":[],"id":"tecnica_mate_dois_bispos","memoria":"semantica","meta":{"dominio":"xadrez","fen_final":"k7/8/1K6/3B4/5B2/8/8/8 b - - 0 1","fen_inicial":"k7/8/1K6/8/5B2/8/8/7B w - - 0 1","fonte":"técnicas de mate","lances_coord":["h1d5"],"orbita_fen":["k7/8/1K6/8/5B2/8/8/7B w - - 0 1","k7/8/1K6/3B4/5B2/8/8/8 b - - 0 1"],"tabuleiro_final":"8 k . . . . . . .\n7 . . . . . . . .\n6 . K . . . . . .\n5 . . . B . . . .\n4 . . . . . B . .\n3 . . . . . . . .\n2 . . . . . . . .\n1 . . . . . . . .\n  a b c d e f g h"},"origem":"wiki","palavras_chave":[],"sinonimos":[],"tipo":"conceito","titulo":"Mate dos Dois Bispos (KBB vs K)"}
\section{Mate dos Dois Bispos (KBB vs K)}
Os dois bispos varrem duas diagonais paralelas, empurrando o rei ao canto (qualquer canto): Bd5\# — um bispo dá o xeque, o outro corta a fuga, o rei branco fecha.
8 k . . . . . . .
7 . . . . . . . .
6 . K . . . . . .
5 . . . B . . . .
4 . . . . . B . .
3 . . . . . . . .
2 . . . . . . . .
1 . . . . . . . .
  a b c d e f g h
FEN inicial: k7/8/1K6/8/5B2/8/8/7B w - - 0 1
\prov{relações: parte\_de tecnicas\_de\_mate; relaciona xeque\_mate}
\prov{proveniência: wiki \textperiodcentered{} técnicas de mate \textperiodcentered{} fato \textperiodcentered{} confiança 1.0 \textperiodcentered{} domínio xadrez \textperiodcentered{} história 3 versões no jornal}

%CRISTAL {"arestas":[{"alvo":"tecnicas_de_mate","confianca":1.0,"epistemico":"fato","origem":"wiki","rel":"parte_de"},{"alvo":"xeque_mate","confianca":1.0,"epistemico":"fato","origem":"wiki","rel":"relaciona"}],"confianca":1.0,"contraexemplos":[],"descricao":"O cavalo dá mate com o rei totalmente preso pelas PRÓPRIAS peças: Nf7#. O rei em h8 cercado por torre g8 e peões g7/h7 não tem para onde ir — só o cavalo alcança.\n8 . . . . . . r k\n7 . . . . . N p p\n6 . . . . . . . .\n5 . . . . . . . .\n4 . . . . . . . .\n3 . . . . . . . .\n2 . . . . . . . .\n1 . . . . . . K .\n  a b c d e f g h\nFEN inicial: 6rk/6pp/8/6N1/8/8/8/6K1 w - - 0 1","epistemico":"fato","exemplos":[],"id":"tecnica_mate_sufocado","memoria":"semantica","meta":{"dominio":"xadrez","fen_final":"6rk/5Npp/8/8/8/8/8/6K1 b - - 0 1","fen_inicial":"6rk/6pp/8/6N1/8/8/8/6K1 w - - 0 1","fonte":"técnicas de mate","lances_coord":["g5f7"],"orbita_fen":["6rk/6pp/8/6N1/8/8/8/6K1 w - - 0 1","6rk/5Npp/8/8/8/8/8/6K1 b - - 0 1"],"tabuleiro_final":"8 . . . . . . r k\n7 . . . . . N p p\n6 . . . . . . . .\n5 . . . . . . . .\n4 . . . . . . . .\n3 . . . . . . . .\n2 . . . . . . . .\n1 . . . . . . K .\n  a b c d e f g h"},"origem":"wiki","palavras_chave":[],"sinonimos":[],"tipo":"conceito","titulo":"Mate Sufocado (Philidor)"}
\section{Mate Sufocado (Philidor)}
O cavalo dá mate com o rei totalmente preso pelas PRÓPRIAS peças: Nf7\#. O rei em h8 cercado por torre g8 e peões g7/h7 não tem para onde ir — só o cavalo alcança.
8 . . . . . . r k
7 . . . . . N p p
6 . . . . . . . .
5 . . . . . . . .
4 . . . . . . . .
3 . . . . . . . .
2 . . . . . . . .
1 . . . . . . K .
  a b c d e f g h
FEN inicial: 6rk/6pp/8/6N1/8/8/8/6K1 w - - 0 1
\prov{relações: parte\_de tecnicas\_de\_mate; relaciona xeque\_mate}
\prov{proveniência: wiki \textperiodcentered{} técnicas de mate \textperiodcentered{} fato \textperiodcentered{} confiança 1.0 \textperiodcentered{} domínio xadrez \textperiodcentered{} história 3 versões no jornal}

%CRISTAL {"arestas":[{"alvo":"tecnicas_de_mate","confianca":1.0,"epistemico":"fato","origem":"wiki","rel":"parte_de"},{"alvo":"xeque_mate","confianca":1.0,"epistemico":"fato","origem":"wiki","rel":"relaciona"}],"confianca":1.0,"contraexemplos":[],"descricao":"O método da caixa: o rei branco corta as casas, a torre dá o xeque na borda — Th8#. A torre não toca o rei; quem prende é o rei branco em oposição.\n8 k . . . . . . R\n7 . . . . . . . .\n6 . K . . . . . .\n5 . . . . . . . .\n4 . . . . . . . .\n3 . . . . . . . .\n2 . . . . . . . .\n1 . . . . . . . .\n  a b c d e f g h\nFEN inicial: k7/8/1K6/8/8/8/8/7R w - - 0 1","epistemico":"fato","exemplos":[],"id":"tecnica_mate_torre","memoria":"semantica","meta":{"dominio":"xadrez","fen_final":"k6R/8/1K6/8/8/8/8/8 b - - 0 1","fen_inicial":"k7/8/1K6/8/8/8/8/7R w - - 0 1","fonte":"técnicas de mate","lances_coord":["h1h8"],"orbita_fen":["k7/8/1K6/8/8/8/8/7R w - - 0 1","k6R/8/1K6/8/8/8/8/8 b - - 0 1"],"tabuleiro_final":"8 k . . . . . . R\n7 . . . . . . . .\n6 . K . . . . . .\n5 . . . . . . . .\n4 . . . . . . . .\n3 . . . . . . . .\n2 . . . . . . . .\n1 . . . . . . . .\n  a b c d e f g h"},"origem":"wiki","palavras_chave":[],"sinonimos":[],"tipo":"conceito","titulo":"Mate da Torre (KR vs K)"}
\section{Mate da Torre (KR vs K)}
O método da caixa: o rei branco corta as casas, a torre dá o xeque na borda — Th8\#. A torre não toca o rei; quem prende é o rei branco em oposição.
8 k . . . . . . R
7 . . . . . . . .
6 . K . . . . . .
5 . . . . . . . .
4 . . . . . . . .
3 . . . . . . . .
2 . . . . . . . .
1 . . . . . . . .
  a b c d e f g h
FEN inicial: k7/8/1K6/8/8/8/8/7R w - - 0 1
\prov{relações: parte\_de tecnicas\_de\_mate; relaciona xeque\_mate}
\prov{proveniência: wiki \textperiodcentered{} técnicas de mate \textperiodcentered{} fato \textperiodcentered{} confiança 1.0 \textperiodcentered{} domínio xadrez \textperiodcentered{} história 3 versões no jornal}

%CRISTAL {"arestas":[{"alvo":"tecnicas_de_mate","confianca":1.0,"epistemico":"fato","origem":"wiki","rel":"parte_de"},{"alvo":"xeque_mate","confianca":1.0,"epistemico":"fato","origem":"wiki","rel":"relaciona"}],"confianca":1.0,"contraexemplos":[],"descricao":"O rei sufocado pelos próprios peões (sem 'luft'): Ta8# na oitava fileira. O mate mais comum das partidas — a lição de fazer o buraco de escape.\n8 R . . . . . . k\n7 . . . . . p p p\n6 . . . . . . . .\n5 . . . . . . . .\n4 . . . . . . . .\n3 . . . . . . . .\n2 . . . . . . . .\n1 . . . . . . K .\n  a b c d e f g h\nFEN inicial: 7k/5ppp/8/8/8/8/8/R5K1 w - - 0 1","epistemico":"fato","exemplos":[],"id":"tecnica_mate_ultima_fileira","memoria":"semantica","meta":{"dominio":"xadrez","fen_final":"R6k/5ppp/8/8/8/8/8/6K1 b - - 0 1","fen_inicial":"7k/5ppp/8/8/8/8/8/R5K1 w - - 0 1","fonte":"técnicas de mate","lances_coord":["a1a8"],"orbita_fen":["7k/5ppp/8/8/8/8/8/R5K1 w - - 0 1","R6k/5ppp/8/8/8/8/8/6K1 b - - 0 1"],"tabuleiro_final":"8 R . . . . . . k\n7 . . . . . p p p\n6 . . . . . . . .\n5 . . . . . . . .\n4 . . . . . . . .\n3 . . . . . . . .\n2 . . . . . . . .\n1 . . . . . . K .\n  a b c d e f g h"},"origem":"wiki","palavras_chave":[],"sinonimos":[],"tipo":"conceito","titulo":"Mate da Última Fileira"}
\section{Mate da Última Fileira}
O rei sufocado pelos próprios peões (sem 'luft'): Ta8\# na oitava fileira. O mate mais comum das partidas — a lição de fazer o buraco de escape.
8 R . . . . . . k
7 . . . . . p p p
6 . . . . . . . .
5 . . . . . . . .
4 . . . . . . . .
3 . . . . . . . .
2 . . . . . . . .
1 . . . . . . K .
  a b c d e f g h
FEN inicial: 7k/5ppp/8/8/8/8/8/R5K1 w - - 0 1
\prov{relações: parte\_de tecnicas\_de\_mate; relaciona xeque\_mate}
\prov{proveniência: wiki \textperiodcentered{} técnicas de mate \textperiodcentered{} fato \textperiodcentered{} confiança 1.0 \textperiodcentered{} domínio xadrez \textperiodcentered{} história 3 versões no jornal}

%CRISTAL {"arestas":[{"alvo":"final_de_xadrez","confianca":1.0,"epistemico":"fato","origem":"wiki","rel":"parte_de"},{"alvo":"xadrez","confianca":1.0,"epistemico":"fato","origem":"wiki","rel":"parte_de"},{"alvo":"xeque_mate","confianca":1.0,"epistemico":"fato","origem":"wiki","rel":"usa"},{"alvo":"xadrez_mineral","confianca":1.0,"epistemico":"fato","origem":"wiki","rel":"relaciona"}],"confianca":1.0,"contraexemplos":[],"descricao":"O vocabulário dos padrões de mate — as técnicas: dama, torre, última fileira, sufocado, bispo e cavalo, dois bispos, árabe. Cada uma é uma posição montada onde o lance de mate é verificado pelo motor. São o alfabeto que as partidas concretizam.","epistemico":"inferencia","exemplos":[],"id":"tecnicas_de_mate","memoria":"semantica","meta":{"dominio":"xadrez","fonte":"técnicas de mate"},"origem":"wiki","palavras_chave":[],"sinonimos":[],"tipo":"conceito","titulo":"Técnicas de Mate"}
\section{Técnicas de Mate}
O vocabulário dos padrões de mate — as técnicas: dama, torre, última fileira, sufocado, bispo e cavalo, dois bispos, árabe. Cada uma é uma posição montada onde o lance de mate é verificado pelo motor. São o alfabeto que as partidas concretizam.
\prov{relações: parte\_de final\_de\_xadrez; parte\_de xadrez; usa xeque\_mate; relaciona xadrez\_mineral}
\prov{proveniência: wiki \textperiodcentered{} técnicas de mate \textperiodcentered{} inferencia \textperiodcentered{} confiança 1.0 \textperiodcentered{} domínio xadrez \textperiodcentered{} história 5 versões no jornal}

%CRISTAL {"arestas":[{"alvo":"tao_da_matematica","confianca":1.0,"epistemico":"fato","origem":"wiki","rel":"usa"},{"alvo":"convencao_vs_verdade_gentil","confianca":1.0,"epistemico":"fato","origem":"wiki","rel":"relaciona"},{"alvo":"deus_postulado","confianca":1.0,"epistemico":"fato","origem":"wiki","rel":"relaciona"}],"confianca":1.0,"contraexemplos":[],"descricao":"Disciplina que Lopes cunhou: 'aproximar-se dos problemas do espírito pela via de uma experimentação abstrata, resultante de processos lógicos da mais moderna físico-matemática'. Escolheu a matemática como fulcro/alicerce — 'é minha religião': se a matemática afirma o absurdo, crê; se a Bíblia afirma o óbvio, desconfia.","epistemico":"inferencia","exemplos":[],"id":"teomatematica","memoria":"semantica","meta":{"autor":"Gentil Lopes","dominio":"mestres","fonte":""},"origem":"wiki","palavras_chave":[],"sinonimos":[],"tipo":"conceito","titulo":"Teomatemática (Lopes)"}
\section{Teomatemática (Lopes)}
Disciplina que Lopes cunhou: 'aproximar-se dos problemas do espírito pela via de uma experimentação abstrata, resultante de processos lógicos da mais moderna físico-matemática'. Escolheu a matemática como fulcro/alicerce — 'é minha religião': se a matemática afirma o absurdo, crê; se a Bíblia afirma o óbvio, desconfia.
\prov{relações: usa tao\_da\_matematica; relaciona convencao\_vs\_verdade\_gentil; relaciona deus\_postulado}
\prov{proveniência: wiki \textperiodcentered{} inferencia \textperiodcentered{} confiança 1.0 \textperiodcentered{} domínio mestres \textperiodcentered{} história 5 versões no jornal}

%CRISTAL {"arestas":[{"alvo":"confucionismo","confianca":1.0,"epistemico":"fato","origem":"wiki","rel":"parte_de"},{"alvo":"contrato_social","confianca":1.0,"epistemico":"fato","origem":"wiki","rel":"relaciona"}],"confianca":1.0,"contraexemplos":[],"descricao":"Confúcio / tradição Chou: a legitimidade do governante depende de sua virtude — o Céu retira o mandato de quem governa em benefício próprio e o confere a alguém mais merecedor. Um imperativo moral (distinto de ming, o Destino).","epistemico":"inferencia","exemplos":[],"id":"tianming","memoria":"semantica","meta":{"autor":"Confúcio","dominio":"mestres","fonte":""},"origem":"wiki","palavras_chave":[],"sinonimos":[],"tipo":"conceito","titulo":"Tianming — o Mandato do Céu"}
\section{Tianming — o Mandato do Céu}
Confúcio / tradição Chou: a legitimidade do governante depende de sua virtude — o Céu retira o mandato de quem governa em benefício próprio e o confere a alguém mais merecedor. Um imperativo moral (distinto de ming, o Destino).
\prov{relações: parte\_de confucionismo; relaciona contrato\_social}
\prov{proveniência: wiki \textperiodcentered{} inferencia \textperiodcentered{} confiança 1.0 \textperiodcentered{} domínio mestres \textperiodcentered{} história 3 versões no jornal}

%CRISTAL {"arestas":[{"alvo":"taoismo","confianca":1.0,"epistemico":"fato","origem":"wiki","rel":"parte_de"},{"alvo":"ethics_of_care","confianca":1.0,"epistemico":"fato","origem":"wiki","rel":"relaciona"}],"confianca":1.0,"contraexemplos":[],"descricao":"Lao Tzu: as três posses que ele valoriza — compaixão, moderação (frugalidade) e humildade ('não ousar ser o primeiro do mundo'). Da compaixão vem a coragem; da moderação, a amplitude.","epistemico":"inferencia","exemplos":[],"id":"tres_tesouros","memoria":"semantica","meta":{"autor":"Lao Tzu","dominio":"mestres","fonte":""},"origem":"wiki","palavras_chave":[],"sinonimos":[],"tipo":"conceito","titulo":"Os Três Tesouros"}
\section{Os Três Tesouros}
Lao Tzu: as três posses que ele valoriza — compaixão, moderação (frugalidade) e humildade ('não ousar ser o primeiro do mundo'). Da compaixão vem a coragem; da moderação, a amplitude.
\prov{relações: parte\_de taoismo; relaciona ethics\_of\_care}
\prov{proveniência: wiki \textperiodcentered{} inferencia \textperiodcentered{} confiança 1.0 \textperiodcentered{} domínio mestres \textperiodcentered{} história 3 versões no jornal}

%CRISTAL {"arestas":[{"alvo":"nao_dualidade","confianca":1.0,"epistemico":"fato","origem":"wiki","rel":"usa"},{"alvo":"tao","confianca":1.0,"epistemico":"fato","origem":"wiki","rel":"usa"}],"confianca":1.0,"contraexemplos":[],"descricao":"Lao Tzu: a raiz da integridade de tudo — 'o céu adquiriu o Um e tornou-se transparente'; quem 'abraça a unidade pode tornar-se indivisível'. A não-dualidade como fonte.","epistemico":"inferencia","exemplos":[],"id":"unidade_taoista","memoria":"semantica","meta":{"autor":"Lao Tzu","dominio":"mestres","fonte":""},"origem":"wiki","palavras_chave":[],"sinonimos":[],"tipo":"conceito","titulo":"Abraçar o Um"}
\section{Abraçar o Um}
Lao Tzu: a raiz da integridade de tudo — 'o céu adquiriu o Um e tornou-se transparente'; quem 'abraça a unidade pode tornar-se indivisível'. A não-dualidade como fonte.
\prov{relações: usa nao\_dualidade; usa tao}
\prov{proveniência: wiki \textperiodcentered{} inferencia \textperiodcentered{} confiança 1.0 \textperiodcentered{} domínio mestres \textperiodcentered{} história 3 versões no jornal}

%CRISTAL {"arestas":[{"alvo":"xadrez","confianca":1.0,"epistemico":"fato","origem":"wiki","rel":"parte_de"}],"confianca":1.0,"contraexemplos":[],"descricao":"A régua material: peão=1, cavalo≈bispo=3, torre=5, dama=9, rei=∞ (a perda encerra o jogo). Aproximação que orienta trocas — mas posição e iniciativa frequentemente valem mais que material.","epistemico":"inferencia","exemplos":[],"id":"valor_das_pecas","memoria":"semantica","meta":{"dominio":"xadrez","fonte":"xadrez"},"origem":"wiki","palavras_chave":[],"sinonimos":[],"tipo":"conceito","titulo":"Valor das Peças"}
\section{Valor das Peças}
A régua material: peão=1, cavalo\ensuremath{\approx}bispo=3, torre=5, dama=9, rei=\ensuremath{\infty} (a perda encerra o jogo). Aproximação que orienta trocas — mas posição e iniciativa frequentemente valem mais que material.
\prov{relações: parte\_de xadrez}
\prov{proveniência: wiki \textperiodcentered{} xadrez \textperiodcentered{} inferencia \textperiodcentered{} confiança 1.0 \textperiodcentered{} domínio xadrez \textperiodcentered{} história 2 versões no jornal}

%CRISTAL {"arestas":[{"alvo":"estrategia_empresarial","confianca":1.0,"epistemico":"fato","origem":"wiki","rel":"parte_de"}],"confianca":1.0,"contraexemplos":[],"descricao":"O que faz uma organização render acima da média de forma sustentável — custo baixo ou diferenciação, ancorados em recursos difíceis de imitar. O 'centro de gravidade' da empresa.","epistemico":"inferencia","exemplos":[],"id":"vantagem_competitiva","memoria":"semantica","meta":{"dominio":"estrategia","fonte":"estratégia"},"origem":"wiki","palavras_chave":[],"sinonimos":[],"tipo":"conceito","titulo":"Vantagem Competitiva"}
\section{Vantagem Competitiva}
O que faz uma organização render acima da média de forma sustentável — custo baixo ou diferenciação, ancorados em recursos difíceis de imitar. O 'centro de gravidade' da empresa.
\prov{relações: parte\_de estrategia\_empresarial}
\prov{proveniência: wiki \textperiodcentered{} estratégia \textperiodcentered{} inferencia \textperiodcentered{} confiança 1.0 \textperiodcentered{} domínio estrategia \textperiodcentered{} história 2 versões no jornal}

%CRISTAL {"arestas":[{"alvo":"borda_critica_vazio","confianca":1.0,"epistemico":"fato","origem":"wiki","rel":"relaciona"},{"alvo":"taoismo","confianca":1.0,"epistemico":"fato","origem":"wiki","rel":"parte_de"},{"alvo":"consciencia_pura","confianca":1.0,"epistemico":"fato","origem":"wiki","rel":"relaciona"}],"confianca":1.0,"contraexemplos":[],"descricao":"Lao Tzu: é o não-ser que confere utilidade ao ser — 'trinta raios convergem ao vazio do centro da roda; da não-existência vem a utilidade'. O vazio do vaso, da roda e da sala é o que os torna úteis. O Espírito do Vale: o vazio não é ausência estéril, mas matriz geradora.","epistemico":"inferencia","exemplos":[],"id":"vazio_util","memoria":"semantica","meta":{"autor":"Lao Tzu","dominio":"mestres","fonte":""},"origem":"wiki","palavras_chave":[],"sinonimos":[],"tipo":"conceito","titulo":"O Vazio Útil (wu — a utilidade do não-ser)"}
\section{O Vazio Útil (wu — a utilidade do não-ser)}
Lao Tzu: é o não-ser que confere utilidade ao ser — 'trinta raios convergem ao vazio do centro da roda; da não-existência vem a utilidade'. O vazio do vaso, da roda e da sala é o que os torna úteis. O Espírito do Vale: o vazio não é ausência estéril, mas matriz geradora.
\prov{relações: relaciona borda\_critica\_vazio; parte\_de taoismo; relaciona consciencia\_pura}
\prov{proveniência: wiki \textperiodcentered{} inferencia \textperiodcentered{} confiança 1.0 \textperiodcentered{} domínio mestres \textperiodcentered{} história 4 versões no jornal}

%CRISTAL {"arestas":[{"alvo":"sabedoria_oriental","confianca":1.0,"epistemico":"fato","origem":"wiki","rel":"parte_de"},{"alvo":"virtualizacao","confianca":0.5,"epistemico":"fato","origem":"wiki","rel":"relaciona"},{"alvo":"word","confianca":0.5,"epistemico":"fato","origem":"wiki","rel":"relaciona"}],"confianca":1.0,"contraexemplos":[],"descricao":"A tradição das Upanishads — sobretudo o Advaita de Shankara: a identidade radical entre o self (atman) e o Absoluto (brahman), velada pela ilusão (maya).","epistemico":"fato","exemplos":[],"id":"vedanta","memoria":"semantica","meta":{"autor":"Shankara","dominio":"mestres","fonte":""},"origem":"wiki","palavras_chave":[],"sinonimos":[],"tipo":"conceito","titulo":"Vedanta (Hinduísmo)"}
\section{Vedanta (Hinduísmo)}
A tradição das Upanishads — sobretudo o Advaita de Shankara: a identidade radical entre o self (atman) e o Absoluto (brahman), velada pela ilusão (maya).
\prov{relações: parte\_de sabedoria\_oriental; relaciona virtualizacao; relaciona word}
\prov{proveniência: wiki \textperiodcentered{} fato \textperiodcentered{} confiança 1.0 \textperiodcentered{} domínio mestres \textperiodcentered{} história 4 versões no jornal}

%CRISTAL {"arestas":[{"alvo":"estrategia_posicional","confianca":1.0,"epistemico":"fato","origem":"wiki","rel":"explica"}],"confianca":1.0,"contraexemplos":[],"descricao":"1º Campeão Mundial (1886) e pai da ESCOLA POSICIONAL: mostrou que o ataque precisa de base posicional — acumular pequenas vantagens antes de golpear. Fundou a teoria moderna.","epistemico":"fato","exemplos":[],"id":"wilhelm_steinitz","memoria":"semantica","meta":{"autor":"Steinitz","dominio":"xadrez","fonte":"história do xadrez"},"origem":"wiki","palavras_chave":[],"sinonimos":[],"tipo":"conceito","titulo":"Wilhelm Steinitz"}
\section{Wilhelm Steinitz}
1º Campeão Mundial (1886) e pai da ESCOLA POSICIONAL: mostrou que o ataque precisa de base posicional — acumular pequenas vantagens antes de golpear. Fundou a teoria moderna.
\prov{relações: explica estrategia\_posicional}
\prov{proveniência: wiki \textperiodcentered{} história do xadrez \textperiodcentered{} fato \textperiodcentered{} confiança 1.0 \textperiodcentered{} domínio xadrez \textperiodcentered{} história 2 versões no jornal}

%CRISTAL {"arestas":[{"alvo":"tao","confianca":1.0,"epistemico":"fato","origem":"wiki","rel":"usa"},{"alvo":"simbolo_vs_significado","confianca":1.0,"epistemico":"fato","origem":"wiki","rel":"relaciona"}],"confianca":1.0,"contraexemplos":[],"descricao":"Lao Tzu: 'Sem-Nome é o princípio do céu e da terra; Com-Nome é a mãe de dez mil coisas'. Nomear é recortar e dividir o real; as distinções (belo/feio, bem/mal) geram-se mutuamente e velam a unidade de origem.","epistemico":"inferencia","exemplos":[],"id":"wu_ming","memoria":"semantica","meta":{"autor":"Lao Tzu","dominio":"mestres","fonte":""},"origem":"wiki","palavras_chave":[],"sinonimos":[],"tipo":"conceito","titulo":"Wu Ming — o Sem-Nome (crítica à linguagem)"}
\section{Wu Ming — o Sem-Nome (crítica à linguagem)}
Lao Tzu: 'Sem-Nome é o princípio do céu e da terra; Com-Nome é a mãe de dez mil coisas'. Nomear é recortar e dividir o real; as distinções (belo/feio, bem/mal) geram-se mutuamente e velam a unidade de origem.
\prov{relações: usa tao; relaciona simbolo\_vs\_significado}
\prov{proveniência: wiki \textperiodcentered{} inferencia \textperiodcentered{} confiança 1.0 \textperiodcentered{} domínio mestres \textperiodcentered{} história 3 versões no jornal}

%CRISTAL {"arestas":[{"alvo":"taoismo","confianca":1.0,"epistemico":"fato","origem":"wiki","rel":"parte_de"},{"alvo":"contemplacao","confianca":1.0,"epistemico":"fato","origem":"wiki","rel":"relaciona"},{"alvo":"estetica_do_fluir_gentil","confianca":1.0,"epistemico":"fato","origem":"wiki","rel":"relaciona"},{"alvo":"vazio_util","confianca":1.0,"epistemico":"fato","origem":"wiki","rel":"usa"},{"alvo":"nirvana","confianca":1.0,"epistemico":"fato","origem":"wiki","rel":"relaciona"}],"confianca":1.0,"contraexemplos":[],"descricao":"Lao Tzu: agir conforme o fluxo natural, sem forçar nem impor intenção — 'o Sábio realiza a obra pela não-ação'. Não é passividade, mas eficácia sem atrito: 'o Caminho é uma constante não-ação que nada deixa por realizar'. O fazer pelo não-fazer.","epistemico":"inferencia","exemplos":[],"id":"wu_wei","memoria":"semantica","meta":{"autor":"Lao Tzu","dominio":"mestres","fonte":""},"origem":"wiki","palavras_chave":[],"sinonimos":[],"tipo":"conceito","titulo":"Wu Wei — a Não-Ação / ação sem esforço"}
\section{Wu Wei — a Não-Ação / ação sem esforço}
Lao Tzu: agir conforme o fluxo natural, sem forçar nem impor intenção — 'o Sábio realiza a obra pela não-ação'. Não é passividade, mas eficácia sem atrito: 'o Caminho é uma constante não-ação que nada deixa por realizar'. O fazer pelo não-fazer.
\prov{relações: parte\_de taoismo; relaciona contemplacao; relaciona estetica\_do\_fluir\_gentil; usa vazio\_util; relaciona nirvana}
\prov{proveniência: wiki \textperiodcentered{} inferencia \textperiodcentered{} confiança 1.0 \textperiodcentered{} domínio mestres \textperiodcentered{} história 6 versões no jornal}

%CRISTAL {"arestas":[{"alvo":"jogo_de_soma_zero","confianca":1.0,"epistemico":"fato","origem":"wiki","rel":"eh_um"},{"alvo":"jogo_estrategico","confianca":1.0,"epistemico":"fato","origem":"wiki","rel":"eh_um"}],"confianca":1.0,"contraexemplos":[],"descricao":"Jogo de tabuleiro 8×8, dois jogadores, soma zero e INFORMAÇÃO PERFEITA: dar xeque-mate ao rei inimigo. Resolvível em teoria por minimax, mas o espaço de jogo é colossal — a fronteira eterna entre cálculo e intuição, e o laboratório da IA por 70 anos.","epistemico":"fato","exemplos":[],"id":"xadrez","memoria":"semantica","meta":{"dominio":"xadrez","fonte":"xadrez"},"origem":"wiki","palavras_chave":[],"sinonimos":[],"tipo":"conceito","titulo":"Xadrez"}
\section{Xadrez}
Jogo de tabuleiro 8$\times$8, dois jogadores, soma zero e INFORMAÇÃO PERFEITA: dar xeque-mate ao rei inimigo. Resolvível em teoria por minimax, mas o espaço de jogo é colossal — a fronteira eterna entre cálculo e intuição, e o laboratório da IA por 70 anos.
\prov{relações: eh\_um jogo\_de\_soma\_zero; eh\_um jogo\_estrategico}
\prov{proveniência: wiki \textperiodcentered{} xadrez \textperiodcentered{} fato \textperiodcentered{} confiança 1.0 \textperiodcentered{} domínio xadrez \textperiodcentered{} história 3 versões no jornal}

%CRISTAL {"arestas":[{"alvo":"xadrez","confianca":1.0,"epistemico":"fato","origem":"wiki","rel":"eh_um"},{"alvo":"final_de_xadrez","confianca":1.0,"epistemico":"fato","origem":"wiki","rel":"relaciona"},{"alvo":"orbita_de_partida","confianca":1.0,"epistemico":"fato","origem":"wiki","rel":"relaciona"}],"confianca":1.0,"contraexemplos":[],"descricao":"O rei em xeque sem nenhum lance legal que o salve: o fim do jogo. É o TERMINAL da árvore de decisão e o ponto fixo da órbita de uma partida — perder a raiz (o rei = autoridade intrínseca) encerra tudo. No BAI, o colapso resolvido.","epistemico":"fato","exemplos":[],"id":"xeque_mate","memoria":"semantica","meta":{"dominio":"xadrez","fonte":"partidas↔órbitas"},"origem":"wiki","palavras_chave":[],"sinonimos":[],"tipo":"conceito","titulo":"Xeque-mate"}
\section{Xeque-mate}
O rei em xeque sem nenhum lance legal que o salve: o fim do jogo. É o TERMINAL da árvore de decisão e o ponto fixo da órbita de uma partida — perder a raiz (o rei = autoridade intrínseca) encerra tudo. No BAI, o colapso resolvido.
\prov{relações: eh\_um xadrez; relaciona final\_de\_xadrez; relaciona orbita\_de\_partida}
\prov{proveniência: wiki \textperiodcentered{} partidas\ensuremath{\leftrightarrow}órbitas \textperiodcentered{} fato \textperiodcentered{} confiança 1.0 \textperiodcentered{} domínio xadrez \textperiodcentered{} história 28 versões no jornal}

%CRISTAL {"arestas":[{"alvo":"confucionismo","confianca":1.0,"epistemico":"fato","origem":"wiki","rel":"parte_de"},{"alvo":"li_ritual","confianca":1.0,"epistemico":"fato","origem":"wiki","rel":"usa"}],"confianca":1.0,"contraexemplos":[],"descricao":"Confúcio: reverência e cuidado para com pais e ancestrais, raiz de toda a ordem social; não basta o sustento material — sem reverência, 'qual a diferença?'. A hierarquia familiar como base do Estado.","epistemico":"inferencia","exemplos":[],"id":"xiao","memoria":"semantica","meta":{"autor":"Confúcio","dominio":"mestres","fonte":""},"origem":"wiki","palavras_chave":[],"sinonimos":[],"tipo":"conceito","titulo":"Xiao — Piedade Filial"}
\section{Xiao — Piedade Filial}
Confúcio: reverência e cuidado para com pais e ancestrais, raiz de toda a ordem social; não basta o sustento material — sem reverência, 'qual a diferença?'. A hierarquia familiar como base do Estado.
\prov{relações: parte\_de confucionismo; usa li\_ritual}
\prov{proveniência: wiki \textperiodcentered{} inferencia \textperiodcentered{} confiança 1.0 \textperiodcentered{} domínio mestres \textperiodcentered{} história 3 versões no jornal}

%CRISTAL {"arestas":[{"alvo":"confucionismo","confianca":1.0,"epistemico":"fato","origem":"wiki","rel":"parte_de"},{"alvo":"pedagogia","confianca":1.0,"epistemico":"fato","origem":"wiki","rel":"relaciona"}],"confianca":1.0,"contraexemplos":[],"descricao":"Confúcio: o aprender como formação moral contínua e inesgotável, não mero acúmulo de fatos — 'amar a benevolência sem amar o aprendizado leva à tolice'. Ele se dizia não um sábio nato, mas 'de coração a aprender'.","epistemico":"inferencia","exemplos":[],"id":"xue_aprendizado","memoria":"semantica","meta":{"autor":"Confúcio","dominio":"mestres","fonte":""},"origem":"wiki","palavras_chave":[],"sinonimos":[],"tipo":"conceito","titulo":"Xue — o Amor ao Aprendizado"}
\section{Xue — o Amor ao Aprendizado}
Confúcio: o aprender como formação moral contínua e inesgotável, não mero acúmulo de fatos — 'amar a benevolência sem amar o aprendizado leva à tolice'. Ele se dizia não um sábio nato, mas 'de coração a aprender'.
\prov{relações: parte\_de confucionismo; relaciona pedagogia}
\prov{proveniência: wiki \textperiodcentered{} inferencia \textperiodcentered{} confiança 1.0 \textperiodcentered{} domínio mestres \textperiodcentered{} história 3 versões no jornal}

%CRISTAL {"arestas":[{"alvo":"junzi","confianca":1.0,"epistemico":"fato","origem":"wiki","rel":"usa"},{"alvo":"confucionismo","confianca":1.0,"epistemico":"fato","origem":"wiki","rel":"parte_de"}],"confianca":1.0,"contraexemplos":[],"descricao":"Confúcio: a conveniência moral do ato nas circunstâncias — fazer o correto por ser correto, em oposição a li (利, lucro/vantagem). 'Diante de uma vantagem, o cavalheiro pensa se é conforme a retidão'. A base da retificação de si.","epistemico":"inferencia","exemplos":[],"id":"yi_retidao","memoria":"semantica","meta":{"autor":"Confúcio","dominio":"mestres","fonte":""},"origem":"wiki","palavras_chave":[],"sinonimos":[],"tipo":"conceito","titulo":"Yi — Retidão / Justiça"}
\section{Yi — Retidão / Justiça}
Confúcio: a conveniência moral do ato nas circunstâncias — fazer o correto por ser correto, em oposição a li (\texttt{U+5229}, lucro/vantagem). 'Diante de uma vantagem, o cavalheiro pensa se é conforme a retidão'. A base da retificação de si.
\prov{relações: usa junzi; parte\_de confucionismo}
\prov{proveniência: wiki \textperiodcentered{} inferencia \textperiodcentered{} confiança 1.0 \textperiodcentered{} domínio mestres \textperiodcentered{} história 3 versões no jornal}

%CRISTAL {"arestas":[{"alvo":"nao_dualidade","confianca":1.0,"epistemico":"fato","origem":"wiki","rel":"relaciona"},{"alvo":"taoismo","confianca":1.0,"epistemico":"fato","origem":"wiki","rel":"parte_de"}],"confianca":1.0,"contraexemplos":[],"descricao":"Lao Tzu / tradição: os opostos geram-se e completam-se mutuamente, e sua alternância é o motor do movimento — 'a existência e a inexistência geram-se uma pela outra'. Por trás dos opostos, uma única energia (o Tao).","epistemico":"inferencia","exemplos":[],"id":"yin_yang","memoria":"semantica","meta":{"autor":"—","dominio":"mestres","fonte":""},"origem":"wiki","palavras_chave":[],"sinonimos":[],"tipo":"conceito","titulo":"Yin-Yang — a Polaridade Complementar"}
\section{Yin-Yang — a Polaridade Complementar}
Lao Tzu / tradição: os opostos geram-se e completam-se mutuamente, e sua alternância é o motor do movimento — 'a existência e a inexistência geram-se uma pela outra'. Por trás dos opostos, uma única energia (o Tao).
\prov{relações: relaciona nao\_dualidade; parte\_de taoismo}
\prov{proveniência: wiki \textperiodcentered{} inferencia \textperiodcentered{} confiança 1.0 \textperiodcentered{} domínio mestres \textperiodcentered{} história 3 versões no jornal}

%CRISTAL {"arestas":[{"alvo":"confucionismo","confianca":1.0,"epistemico":"fato","origem":"wiki","rel":"parte_de"},{"alvo":"simbolo_vs_significado","confianca":1.0,"epistemico":"fato","origem":"wiki","rel":"usa"},{"alvo":"convencao_vs_verdade_gentil","confianca":1.0,"epistemico":"fato","origem":"wiki","rel":"relaciona"},{"alvo":"wu_ming","confianca":1.0,"epistemico":"fato","origem":"wiki","rel":"contradiz"}],"confianca":1.0,"contraexemplos":[],"descricao":"Confúcio: 'quando os nomes não são corretos, o que é dito não soa razoável, e os negócios não culminam em sucesso'. Cada coisa e papel deve corresponder ao seu nome ('que o governante seja governante, o pai seja pai') — a correspondência nome↔coisa como condição de verdade e de ação correta.","epistemico":"inferencia","exemplos":[],"id":"zhengming","memoria":"semantica","meta":{"autor":"Confúcio","dominio":"mestres","fonte":""},"origem":"wiki","palavras_chave":[],"sinonimos":[],"tipo":"conceito","titulo":"Zhengming — a Retificação dos Nomes"}
\section{Zhengming — a Retificação dos Nomes}
Confúcio: 'quando os nomes não são corretos, o que é dito não soa razoável, e os negócios não culminam em sucesso'. Cada coisa e papel deve corresponder ao seu nome ('que o governante seja governante, o pai seja pai') — a correspondência nome\ensuremath{\leftrightarrow}coisa como condição de verdade e de ação correta.
\prov{relações: parte\_de confucionismo; usa simbolo\_vs\_significado; relaciona convencao\_vs\_verdade\_gentil; contradiz wu\_ming}
\prov{proveniência: wiki \textperiodcentered{} inferencia \textperiodcentered{} confiança 1.0 \textperiodcentered{} domínio mestres \textperiodcentered{} história 5 versões no jornal}

%CRISTAL {"arestas":[{"alvo":"confucionismo","confianca":1.0,"epistemico":"fato","origem":"wiki","rel":"parte_de"},{"alvo":"eudaimonia","confianca":1.0,"epistemico":"fato","origem":"wiki","rel":"relaciona"}],"confianca":1.0,"contraexemplos":[],"descricao":"Confúcio: o 'meio inabalável' — nem excesso nem deficiência (zhong = não pender; yong = constante). Não é mediania medíocre, mas o ponto ótimo de equilíbrio onde a virtude reside; ecoa o meio-termo (mesótes) de Aristóteles.","epistemico":"inferencia","exemplos":[],"id":"zhong_yong","memoria":"semantica","meta":{"autor":"Confúcio","dominio":"mestres","fonte":""},"origem":"wiki","palavras_chave":[],"sinonimos":[],"tipo":"conceito","titulo":"Zhong Yong — a Doutrina do Meio"}
\section{Zhong Yong — a Doutrina do Meio}
Confúcio: o 'meio inabalável' — nem excesso nem deficiência (zhong = não pender; yong = constante). Não é mediania medíocre, mas o ponto ótimo de equilíbrio onde a virtude reside; ecoa o meio-termo (mesótes) de Aristóteles.
\prov{relações: parte\_de confucionismo; relaciona eudaimonia}
\prov{proveniência: wiki \textperiodcentered{} inferencia \textperiodcentered{} confiança 1.0 \textperiodcentered{} domínio mestres \textperiodcentered{} história 3 versões no jornal}

%CRISTAL {"arestas":[{"alvo":"wu_wei","confianca":1.0,"epistemico":"fato","origem":"wiki","rel":"relaciona"},{"alvo":"nao_dualidade","confianca":1.0,"epistemico":"fato","origem":"wiki","rel":"relaciona"}],"confianca":1.0,"contraexemplos":[],"descricao":"Lao Tzu: o 'assim-de-si-mesmo' — cada ser seguindo a própria natureza sem imposição externa. O wu wei é a prática que deixa o ziran acontecer; 'os dez mil seres se transformam por si'.","epistemico":"inferencia","exemplos":[],"id":"ziran","memoria":"semantica","meta":{"autor":"Lao Tzu","dominio":"mestres","fonte":""},"origem":"wiki","palavras_chave":[],"sinonimos":[],"tipo":"conceito","titulo":"Ziran — a Naturalidade / Espontaneidade"}
\section{Ziran — a Naturalidade / Espontaneidade}
Lao Tzu: o 'assim-de-si-mesmo' — cada ser seguindo a própria natureza sem imposição externa. O wu wei é a prática que deixa o ziran acontecer; 'os dez mil seres se transformam por si'.
\prov{relações: relaciona wu\_wei; relaciona nao\_dualidade}
\prov{proveniência: wiki \textperiodcentered{} inferencia \textperiodcentered{} confiança 1.0 \textperiodcentered{} domínio mestres \textperiodcentered{} história 3 versões no jornal}

\end{document}
